% PV-Finder paper deck, 13 November 2026. Build: bash build.sh
% Every number is a macro from numbers.tex; red TBD = not final yet.
% Figures come from images.txt through images_manifest.tex (written by build.sh).
\documentclass[xcolor=table,aspectratio=169]{beamer}

\usepackage[defaultsans,scale=0.95]{opensans}
\usepackage{FiraMono}
\usepackage{fourier}
\usefonttheme{professionalfonts}
\DeclareMathAlphabet{\mathcal}{OMS}{cmsy}{m}{n}
\newcommand\bmmax{2}
\usepackage{bm}
\usepackage[T1]{fontenc}

\usepackage{graphicx}
\usepackage{caption}
\captionsetup[figure]{labelformat=empty}
\usepackage{subcaption}
\usepackage{multicol}

\usepackage{tikz}
\usetikzlibrary{decorations.pathreplacing}
\usepackage{wrapfig}

% COLORS
\usepackage{xcolor}
\colorlet{myred}{red!80!black}
\colorlet{myblue}{blue!80!black}
\colorlet{mygreen}{green!60!black}
\colorlet{myorange}{orange!70!red!60!black}
\colorlet{mydarkred}{red!30!black}
\colorlet{mydarkblue}{blue!40!black}
\colorlet{mydarkgreen}{green!30!black}
\colorlet{mylightblue}{blue!60!white}
\colorlet{mydarkyellow}{yellow!70!black}
\colorlet{myyellowred}{yellow!70!red}

\usepackage{float}
\usepackage{hyperref}

\usepackage{amsmath, amsthm, amssymb, amsfonts}
\usepackage{esint}
\usepackage{mathtools}
\usepackage{listings}
\lstset{basicstyle=\ttfamily}

\usepackage{physics}
\usepackage[makeroom]{cancel}
% siunitx + mhchem commented out due to version clash with chemgreek on this TL2017 install
\usepackage[italic]{heppennames}
\usepackage{xpatch}
\makeatletter
\xpatchcmd\@HepConStyle
{\edef\@upcode{\updefault}}
{\ifdefined\shapedefault\edef\@upcode{\shapedefault}\else\edef\@upcode{\updefault}\fi}
{}{}
\makeatother
\usepackage{tensor}
\usepackage{feynmp-auto}
\usepackage{listofitems}
\usetikzlibrary{arrows.meta}
\usepackage[outline]{contour}
\usepackage{booktabs}

%%% ---- shared TikZ vocabulary for the Part I method figures ----
\usetikzlibrary{positioning,calc,backgrounds}
\tikzset{
    pvflow/.style   ={-{Latex[length=2.4mm,width=1.9mm]}, line width=0.9pt, draw=black!55},
    pvpoint/.style  ={-{Latex[length=2.0mm,width=1.6mm]}, line width=0.9pt},
    pvaxis/.style   ={-{Latex[length=1.9mm,width=1.5mm]}, line width=0.7pt, draw=black!60},
    pvdim/.style    ={{Latex[length=1.7mm,width=1.4mm]}-{Latex[length=1.7mm,width=1.4mm]},
                      line width=0.7pt},
    pvcurve/.style  ={line width=1.0pt, line join=round, line cap=round},
    pvthresh/.style ={dash pattern=on 2.6pt off 2.2pt, line width=0.8pt,
                      draw=myorange!85!black},
    pvsplit/.style  ={dash pattern=on 3pt off 2.4pt, line width=1.0pt, draw=pennred},
    pvguide/.style  ={line width=0.5pt, draw=black!15},
    pvcard/.style   ={rounded corners=6pt, draw=black!15, fill=black!2, line width=0.7pt},
    pvbox/.style    ={rounded corners=6pt, align=center, line width=0.9pt},
}
\newcommand{\pvcap}[1]{{\scriptsize\textcolor{black!55}{#1}}}


\addtobeamertemplate{frametitle}{}{\vspace{-15.5pt}}

%%%%%%%%%%%%%%%%%%%%%%%%%%%%%%%%%%%%%%

\newcommand{\backupbegin}{
    \newcounter{finalframe}
    \setcounter{finalframe}{\value{framenumber}}
    \setbeamertemplate{footline}{}
    \begin{frame}[plain]
        \centering
        \vfill
        {\color{pennred!28}\rule{0.30\linewidth}{1.1pt}}\\[0.8em]
        {\LARGE\textbf{\textcolor{pennred}{Backup}}}\\[0.8em]
        {\color{pennred!28}\rule{0.30\linewidth}{1.1pt}}
        \vfill
    \end{frame}
}
\newcommand{\backupend}{
    \setcounter{framenumber}{\value{finalframe}}
    \setbeamertemplate{footline}[frame number]
}

\newcommand{\ntrk}{\ensuremath{n_\mathrm{trk}}}

\def\thisframelogos{}
\newcommand{\framelogo}[1]{\def\thisframelogos{#1}}

\addtobeamertemplate{frametitle}{}{%
\begin{tikzpicture}[remember picture,overlay]
\node[anchor=north east,yshift=3.5pt] at (current page.north east) {%
    \foreach \img in \thisframelogos {%
        \hspace{.5ex}%
        \includegraphics[height=0.8cm, keepaspectratio]{\img}%
    }%
};
\end{tikzpicture}}

% Asset folders are symlinked locally (e.g. logo-ATLAS -> ../../legacy/...)

% ---------------------------------------------------------------------------
% (Copied from the EB deck, 11_13_2026_EB; the paper deck uses the reading build.)
% Two builds from one source.  See build.sh.
%
%   slides.pdf       reading version   (default, \talkbuild = 0)
%   slides_talk.pdf  presentation      (pdflatex "\def\talkbuild{1}% PV-Finder paper deck, 13 November 2026. Build: bash build.sh
% Every number is a macro from numbers.tex; red TBD = not final yet.
% Figures come from images.txt through images_manifest.tex (written by build.sh).
\documentclass[xcolor=table,aspectratio=169]{beamer}

\usepackage[defaultsans,scale=0.95]{opensans}
\usepackage{FiraMono}
\usepackage{fourier}
\usefonttheme{professionalfonts}
\DeclareMathAlphabet{\mathcal}{OMS}{cmsy}{m}{n}
\newcommand\bmmax{2}
\usepackage{bm}
\usepackage[T1]{fontenc}

\usepackage{graphicx}
\usepackage{caption}
\captionsetup[figure]{labelformat=empty}
\usepackage{subcaption}
\usepackage{multicol}

\usepackage{tikz}
\usetikzlibrary{decorations.pathreplacing}
\usepackage{wrapfig}

% COLORS
\usepackage{xcolor}
\colorlet{myred}{red!80!black}
\colorlet{myblue}{blue!80!black}
\colorlet{mygreen}{green!60!black}
\colorlet{myorange}{orange!70!red!60!black}
\colorlet{mydarkred}{red!30!black}
\colorlet{mydarkblue}{blue!40!black}
\colorlet{mydarkgreen}{green!30!black}
\colorlet{mylightblue}{blue!60!white}
\colorlet{mydarkyellow}{yellow!70!black}
\colorlet{myyellowred}{yellow!70!red}

\usepackage{float}
\usepackage{hyperref}

\usepackage{amsmath, amsthm, amssymb, amsfonts}
\usepackage{esint}
\usepackage{mathtools}
\usepackage{listings}
\lstset{basicstyle=\ttfamily}

\usepackage{physics}
\usepackage[makeroom]{cancel}
% siunitx + mhchem commented out due to version clash with chemgreek on this TL2017 install
\usepackage[italic]{heppennames}
\usepackage{xpatch}
\makeatletter
\xpatchcmd\@HepConStyle
{\edef\@upcode{\updefault}}
{\ifdefined\shapedefault\edef\@upcode{\shapedefault}\else\edef\@upcode{\updefault}\fi}
{}{}
\makeatother
\usepackage{tensor}
\usepackage{feynmp-auto}
\usepackage{listofitems}
\usetikzlibrary{arrows.meta}
\usepackage[outline]{contour}
\usepackage{booktabs}

%%% ---- shared TikZ vocabulary for the Part I method figures ----
\usetikzlibrary{positioning,calc,backgrounds}
\tikzset{
    pvflow/.style   ={-{Latex[length=2.4mm,width=1.9mm]}, line width=0.9pt, draw=black!55},
    pvpoint/.style  ={-{Latex[length=2.0mm,width=1.6mm]}, line width=0.9pt},
    pvaxis/.style   ={-{Latex[length=1.9mm,width=1.5mm]}, line width=0.7pt, draw=black!60},
    pvdim/.style    ={{Latex[length=1.7mm,width=1.4mm]}-{Latex[length=1.7mm,width=1.4mm]},
                      line width=0.7pt},
    pvcurve/.style  ={line width=1.0pt, line join=round, line cap=round},
    pvthresh/.style ={dash pattern=on 2.6pt off 2.2pt, line width=0.8pt,
                      draw=myorange!85!black},
    pvsplit/.style  ={dash pattern=on 3pt off 2.4pt, line width=1.0pt, draw=pennred},
    pvguide/.style  ={line width=0.5pt, draw=black!15},
    pvcard/.style   ={rounded corners=6pt, draw=black!15, fill=black!2, line width=0.7pt},
    pvbox/.style    ={rounded corners=6pt, align=center, line width=0.9pt},
}
\newcommand{\pvcap}[1]{{\scriptsize\textcolor{black!55}{#1}}}


\addtobeamertemplate{frametitle}{}{\vspace{-15.5pt}}

%%%%%%%%%%%%%%%%%%%%%%%%%%%%%%%%%%%%%%

\newcommand{\backupbegin}{
    \newcounter{finalframe}
    \setcounter{finalframe}{\value{framenumber}}
    \setbeamertemplate{footline}{}
    \begin{frame}[plain]
        \centering
        \vfill
        {\color{pennred!28}\rule{0.30\linewidth}{1.1pt}}\\[0.8em]
        {\LARGE\textbf{\textcolor{pennred}{Backup}}}\\[0.8em]
        {\color{pennred!28}\rule{0.30\linewidth}{1.1pt}}
        \vfill
    \end{frame}
}
\newcommand{\backupend}{
    \setcounter{framenumber}{\value{finalframe}}
    \setbeamertemplate{footline}[frame number]
}

\newcommand{\ntrk}{\ensuremath{n_\mathrm{trk}}}

\def\thisframelogos{}
\newcommand{\framelogo}[1]{\def\thisframelogos{#1}}

\addtobeamertemplate{frametitle}{}{%
\begin{tikzpicture}[remember picture,overlay]
\node[anchor=north east,yshift=3.5pt] at (current page.north east) {%
    \foreach \img in \thisframelogos {%
        \hspace{.5ex}%
        \includegraphics[height=0.8cm, keepaspectratio]{\img}%
    }%
};
\end{tikzpicture}}

% Asset folders are symlinked locally (e.g. logo-ATLAS -> ../../legacy/...)

% ---------------------------------------------------------------------------
% (Copied from the EB deck, 11_13_2026_EB; the paper deck uses the reading build.)
% Two builds from one source.  See build.sh.
%
%   slides.pdf       reading version   (default, \talkbuild = 0)
%   slides_talk.pdf  presentation      (pdflatex "\def\talkbuild{1}\input{slides}")
%
% \readingonly{...} keeps its argument only in the reading build; \talkonly{...}
% only in the talk build.  Every number lives in exactly one place, so the two
% versions cannot drift apart.
\providecommand{\talkbuild}{0}
\newcommand{\readingonly}[1]{\ifnum\talkbuild=0\relax#1\fi}
\newcommand{\talkonly}[1]{\ifnum\talkbuild=1\relax#1\fi}
% \bnum{reading}{talk} expands to a bare number, so it can be used inside a
% length: \includegraphics[height=\bnum{0.5}{0.6}\textheight]{...}.  The talk
% build drops prose, which frees vertical room for the figure that is left.
\newcommand{\bnum}[2]{\ifnum\talkbuild=0 #1\else #2\fi}
% ---------------------------------------------------------------------------

% Part-divider helper: \partdivider{numeral}{title}{one-line subtitle}
\newcommand{\partdivider}[3]{%
    \begin{frame}[plain]
        \centering
        \vfill
        \begin{tikzpicture}
            \fill[pennred] (0,0) circle (0.72cm);
            \node[text=white, font=\Large\bfseries] at (0,0.02) {#1};
        \end{tikzpicture}\\[0.85em]
        {\LARGE\textbf{\textcolor{pennred}{#2}}}\\[0.75em]
        {\color{pennred!28}\rule{0.30\linewidth}{1.1pt}}\\[0.75em]
        {\small\textcolor{black!50}{#3}}
        \vfill
    \end{frame}%
}

%%%%%%%%%%%%%%%%%%%%%%%%%%%%%%%%%%%%%%

\usetheme{Philadelphia}
\hypersetup{colorlinks,allcolors=pennred}

\title[PV-Finder]{PV-Finder: primary vertex reconstruction and track association with neural networks}
\subtitle{Paper deck: configuration, definitions, answers and figures}

\author[Mattia Mauro]{Qi-Bin Lei, Rocky B. Garg, Lauren Tompkins,\\\textcolor{pennred}{Mattia Mauro}, Trey Warner}

\institute[]{Stanford University}

\date{November 13, 2026}

\titlegraphic{
    \hspace{5pt}
    \href{https://stanford.edu/group/stanford_atlas/}{\includegraphics[height=1cm]{logo-Stanford/Block_S_2_color.png}}
    \href{https://atlas.cern/}{\includegraphics[trim={0 11.5cm 0 11.5cm},clip,height=1cm]{logo-ATLAS/ATLAS-Logo-Ref-Pantone.pdf}}
    \href{https://iris-hep.org/}{\includegraphics[height=1cm]{logo-iris-hep/Iris-hep-4-no-long-name.png}}
}


\newcommand{\readingnote}[1]{\readingonly{\par\vspace{.35em}%
  {\footnotesize\color{black!65}#1\par}}}
\newcommand{\takeaway}[1]{\par\vspace{.65em}{\centering
  \setlength{\fboxsep}{7pt}\colorbox{pennred!5}{\parbox{.94\linewidth}{%
  \centering\small\textcolor{pennred}{#1}}}\par}}
\newcommand{\context}[1]{{\centering\footnotesize\color{black!55}#1\par}\vspace{.65em}}
\newcommand{\bigvalue}[2]{{\LARGE\bfseries\textcolor{pennred}{#1}}\\[.4em]{\small #2}\par}
\newcommand{\pubnote}{\href{https://cds.cern.ch/record/2858348}{ATL-PHYS-PUB-2023-011}}
\newcommand{\notelink}{\href{https://gitlab.cern.ch/atlas-physics-office/IDTR/ANA-IDTR-2026-01/ANA-IDTR-2026-01-INT1}{ANA-IDTR-2026-01-INT1}}
\newcommand{\treysource}{\href{https://gitlab.cern.ch/atlas-physics-office/IDTR/ANA-IDTR-2026-01/ANA-IDTR-2026-01-INT1/-/blob/4c50ed51500c897dd6a1b5bb52b0f1e3c12e7575/sections/04b_timing.tex}{Note Sec. 5.1 (18 Sep 2026)}}

\newcommand{\svtx}{\sigma_{\mathrm{vtx-vtx}}}
\setbeamertemplate{navigation symbols}{}
\hypersetup{pdftitle={PV-Finder, paper deck, November 13, 2026},
  pdfauthor={Mattia Mauro, Qi-Bin Lei, Rocky B. Garg, Lauren Tompkins, Trey Warner}}

% ---------------------------------------------------------------------------
% Paper deck helpers
\usepackage{array}
% \imgsrc{name}{expected source path}{ext}: written by build.sh into
% images_manifest.tex. \plot{name}{height} includes images/name.ext when the
% figure exists and draws a framed placeholder with the expected path if not.
\makeatletter
\newcommand{\imgsrc}[3]{%
  \expandafter\def\csname pd@src@#1\endcsname{#2}%
  \expandafter\def\csname pd@ext@#1\endcsname{#3}}
\newcommand{\plot}[2]{%
  \expandafter\ifx\csname pd@ext@#1\endcsname\relax
    \pd@placeholder{#1}{#2}{(not in images.txt)}%
  \else
    \edef\pd@e{\csname pd@ext@#1\endcsname}%
    \ifx\pd@e\@empty
      \pd@placeholder{#1}{#2}{\csname pd@src@#1\endcsname.\{pdf,png\}}%
    \else
      \includegraphics[width=\linewidth,height=#2\textheight,keepaspectratio]{images/#1.\pd@e}%
    \fi
  \fi}
\newcommand{\pd@placeholder}[3]{%
  {\setlength{\fboxsep}{6pt}\setlength{\fboxrule}{0.8pt}%
   \fcolorbox{red!60!black}{red!3}{\parbox[c][#2\textheight][c]{.9\linewidth}{\centering
   {\small\bfseries\textcolor{red!60!black}{Figure not produced yet}}\\[.5em]
   {\scriptsize\ttfamily #3}\\[.4em]
   {\tiny\textcolor{black!55}{image \texttt{\detokenize{#1}}; rebuild with build.sh when it exists}}}}}}
\makeatother
% Role badges for the results section
\newcommand{\mainbadge}{{\setlength{\fboxsep}{2pt}\colorbox{mygreen!15}{\scriptsize\bfseries\textcolor{mygreen!70!black}{PAPER MAIN}}}}
\newcommand{\backupbadge}{{\setlength{\fboxsep}{2pt}\colorbox{black!8}{\scriptsize\bfseries\textcolor{black!60}{BACKUP / NOTE}}}}
\newcommand{\fcap}[1]{\par\vspace{.25em}{\footnotesize\color{black!70}#1\par}}
\newcommand{\src}[1]{{\raggedright\tiny\color{black!50}\texttt{#1}\par}}
\newcommand{\draftref}[1]{{\small\textcolor{pennred}{\textbf{Draft #1}}}}
\newcommand{\qtext}[1]{{\footnotesize\itshape\color{black!65}#1\par}\vspace{.4em}}
\newcommand{\final}[1]{\textcolor{mydarkblue}{\textbf{Final:}} #1}

% numbers.tex: every number of the paper deck, one macro each.
%
% Rule: \TBD (red) marks a number that is not final yet. Each TBD line names
% the doc and the JSON key it is filled from. "seen" gives the value present in
% that file on 12-13 Nov 2026, before the recompute agents finished: it is for
% checking the fill only, never for quoting.
%
% Sources:
%   HLDOC = docs/evaluation/hllhc_final_numbers_2026-11-12.md (HL-LHC agent)
%   LPDOC = docs/evaluation/lowpu_final_numbers_2026-11-12.md (low pile-up agent)
%   H = outputs/11_12_2026_output/hllhc_final    L = outputs/11_12_2026_output/lowpu_final
%   CH  = H/chain/chain_headline.json :: headline/PU200 (185-215)/...
%   CHI = H/chain/chain_headline.json :: headline/inclusive/...
%   SUP = H/supplementary/supplementary_summary.json
%   CX  = H/chain_extras/chain_extras_summary.json
%   H1  = H/h1_vertex_finding/h1_summary.json (r16443, PU200, 1,920 events)
%   RVM = H/resolution_vs_mu/results.json
%   CPU = H/cpu_timing/cpu_timing.json (and H/cpu_timing/README.md)
%   VF  = L/summary/vf_results.json   GM = L/gnn_mc/summary.json :: bins/inclusive/...
%   R3A = L/run3_assoc/summary.json :: by_t/0.97/...
% Final numbers without a TBD come from committed code or committed docs.
\providecommand{\TBD}{\textcolor{red}{\textnormal{TBD}}}

%======================================================================
% 1. Frozen configuration (src/pv_finder/utils/working_points.py, final)
%======================================================================
\newcommand{\winMM}{0.17}
\newcommand{\hsWinMM}{0.1}
\newcommand{\assignMM}{1}
\newcommand{\pairBins}{300}
\newcommand{\purityCut}{70\%}
% HL-LHC density model s8
\newcommand{\hlModel}{s8}
\newcommand{\hlCkpt}{hllhc\_alleta\_v6\_s8\_\allowbreak 1p1ep\_phase2\_\allowbreak step\_450000}
\newcommand{\hlUnetCh}{80}
\newcommand{\hlLatent}{1}
\newcommand{\hlHidden}{$5\times128$}
% training log outputs/logs/hllhc_alleta_v6_s8_1p1ep_20261109_211919.log
\newcommand{\hlParams}{439{,}593}
% derived: MLP 7-128x5-1000 = 196,072; U-Net = 439,593 - 196,072
\newcommand{\hlParamsMLP}{196{,}072}
\newcommand{\hlParamsUnet}{243{,}521}
\newcommand{\hlAmp}{0.01}
\newcommand{\hlWidth}{3}
\newcommand{\hlIntegral}{0.184}
\newcommand{\hlHeight}{0.0184}
\newcommand{\hlCentroid}{3}
% HL-LHC GNN
\newcommand{\hlGnnEpoch}{527}
\newcommand{\hlT}{0.97}
\newcommand{\gnnParams}{39{,}521}
\newcommand{\gnnK}{20}
% low pile-up
\newcommand{\lpModel}{PVF-e400}
\newcommand{\lpCkpt}{e2e\_mlpHist50\_e2e400\_\allowbreak 1latent\_mse\_\allowbreak phase2\_epoch\_400}
\newcommand{\lpParams}{358{,}713}
\newcommand{\lpUnetCh}{64}
\newcommand{\lpHidden}{$5\times100$}
\newcommand{\lpIntegral}{0.50}
\newcommand{\lpGnnCkpt}{ttva\_gat\_mu60\_v2\_aug epoch 115}
% working_points.LOWPU.gnn.t (frozen 12 Nov; LPDOC Sec. 3)
\newcommand{\lpT}{0.97}

%======================================================================
% 2. Datasets (final)
%======================================================================
% docs/additional_studies/paper_questions_2026-11-11.md Part 3;
% outputs/11_11_2026_output/datasets/hllhc_dataset_inventory.json
\newcommand{\nPoolEvt}{5{,}173{,}600}
\newcommand{\nTrainSub}{59{,}599{,}872}
\newcommand{\nValSub}{1{,}862{,}496}
\newcommand{\nTestSub}{620{,}832}
\newcommand{\nTrainEq}{4{,}966{,}656}
\newcommand{\nValEq}{155{,}208}
\newcommand{\nTestEq}{51{,}736}
\newcommand{\nEvtPU}{1{,}920}
\newcommand{\nEvtPUb}{1{,}888}
\newcommand{\nTruthPU}{213{,}738}
\newcommand{\truthEvt}{111.3}
\newcommand{\meanMuPU}{192.5}
% GNN v2: outputs/11_11_2026_output/gnn_fullmu_v2/split_manifest.json
\newcommand{\gnnTrainA}{60{,}000}
\newcommand{\gnnTrainB}{89{,}820}
\newcommand{\gnnTrain}{149{,}820}
\newcommand{\gnnVal}{7{,}500}
\newcommand{\gnnMon}{1{,}500}
\newcommand{\gnnResTest}{7{,}500}
\newcommand{\gnnResTestB}{9{,}980}
% resolution vs mu sample (docs/evaluation/resolution_vs_mu_2026-11-11.md)
\newcommand{\nEvtResMu}{44{,}999}
% low pile-up: outputs/11_12_2026_output/paper_deck/lowpu_mu/lowpu_mu_distributions.json
\newcommand{\lpMCTotal}{51{,}000}
\newcommand{\lpMCTrain}{35{,}700}
\newcommand{\lpMCVal}{12{,}750}
\newcommand{\lpMCTest}{2{,}550}
\newcommand{\lpMCNonTest}{48{,}450}
\newcommand{\lpMuMC}{38.9}
\newcommand{\lpMuRunTwo}{42.5}
\newcommand{\lpMuRunThree}{50.0}
\newcommand{\lpTruthMC}{72{,}189}

%======================================================================
% 3. HL-LHC before the GNN, PU200 (r16443, 1,920 events)
%======================================================================
% AMVF numbers do not depend on the density model: final, from
% docs/evaluation/model_freeze_candidates_2026-11-12.md (0.17 mm).
\newcommand{\effAMVF}{79.85}
\newcommand{\effAMVFb}{80.13}
\newcommand{\fakeFracAMVF}{20.08}
\newcommand{\fakeEvtAMVF}{19.65}
\newcommand{\splitEvtAMVF}{0.184}
\newcommand{\recoEvtAMVF}{97.9}
\newcommand{\sigAMVF}{0.305}
\newcommand{\sigAMVFErr}{0.006}
\newcommand{\pseAMVF}{56.45}
\newcommand{\hsRecoAMVF}{99.79}
% window derivation (v6 and AMVF, docs/evaluation/matching_window_audit_2026-11-09.md)
\newcommand{\pseVsix}{55.05}
% PV-Finder s8 at the frozen cuts: TBD
\newcommand{\effPVF}{\TBD}      % HLDOC | H1 Eq.(1) eff r16443 | seen: freeze doc 83.05 (v6 83.20 - 0.15)
\newcommand{\effPVFErr}{\TBD}   % HLDOC | H1 bootstrap error
\newcommand{\effDiff}{\TBD}     % HLDOC | H1 PVF - AMVF, paired [pt] | seen: about +3.2
\newcommand{\effDiffErr}{\TBD}  % HLDOC | H1 paired bootstrap error
\newcommand{\effPVFb}{\TBD}     % HLDOC | H1 Eq.(1) eff r16638 (1,888 events) | seen: 83.21
\newcommand{\effOneOnePVF}{\TBD}  % HLDOC | H1 one-to-one efficiency, PVF
\newcommand{\effOneOneAMVF}{\TBD} % HLDOC | H1 one-to-one efficiency, AMVF (0.1694 mm: 70.07)
\newcommand{\fakeFracPVF}{\TBD} % HLDOC | H1 fake / reco [%] | seen: about 18.5
\newcommand{\fakeEvtPVF}{\TBD}  % HLDOC | H1 fakes / event | seen: 18.64
\newcommand{\splitEvtPVF}{\TBD} % HLDOC | H1 splits / event
\newcommand{\recoEvtPVF}{\TBD}  % HLDOC | H1 reco / event
\newcommand{\sigPVF}{\TBD}      % HLDOC | H1 sigma_vtx-vtx [mm], counted peaks | seen: 0.217 at 0.182 cuts
\newcommand{\sigPVFErr}{\TBD}   % HLDOC | H1 fit error
\newcommand{\psePVF}{\TBD}      % HLDOC | H1 68% residual [um] | seen: 55.60 at 0.182 cuts
\newcommand{\hsRecoPVF}{\TBD}   % HLDOC | H1 / CH HS reco at 0.1 mm [%] | seen: 99.90 (1,918)
\newcommand{\chancePVF}{\TBD}   % HLDOC | H1 3 mm-shift chance match [%] (v6: 18.6)
\newcommand{\chanceAMVF}{\TBD}  % HLDOC | H1 3 mm-shift chance match [%] (0.1694 mm: 18.4)
\newcommand{\scanLo}{\TBD}      % HLDOC | H1 window scan: smallest PVF-AMVF over 0.10-0.20 mm (v6: +3.1)
\newcommand{\scanHi}{\TBD}      % HLDOC | H1 window scan: largest PVF-AMVF over 0.10-0.20 mm (v6: +3.4)
% resolution vs mu (s8 rerun of the 11 Nov study; AMVF final from that study)
\newcommand{\pseMuLoAMVF}{32.1}
\newcommand{\pseMuHiAMVF}{56.5}
\newcommand{\sigMuHundredAMVF}{275}
\newcommand{\sigMuTopAMVF}{304}
\newcommand{\pseMuLoPVF}{\TBD}      % HLDOC | RVM bin 0-10 p68 [um] (v6: 33.7)
\newcommand{\pseMuHiPVF}{\TBD}      % HLDOC | RVM bin 180-200 p68 [um] (v6: 54.9)
\newcommand{\sigMuHundredPVF}{\TBD} % HLDOC | RVM bin 100-120 sigma [um] (v6: 194.8)
\newcommand{\sigMuTopPVF}{\TBD}     % HLDOC | RVM bin 180-200 sigma [um] (v6: 216.2)
\newcommand{\sigRatioMu}{\TBD}      % HLDOC | RVM AMVF/PVF sigma range for mu >= 100 (v6: 1.41-1.43)

%======================================================================
% 4. HL-LHC after the GNN (s8 + GNN v2 e527, t = 0.97)
%======================================================================
% AMVF, model independent, final (gnn_v2_freeze_2026-11-12.md test r16443;
% metric_definitions.md Sec. 1; supplementary counts 66,127 / 80,104 / 41,594 / 135 over 1,920)
\newcommand{\cleanTruthAMVF}{30.94}
\newcommand{\cleanTruthAMVFincl}{40.49}
\newcommand{\cleanRecoAMVF}{35.2}
\newcommand{\mergedRecoAMVF}{42.6}
\newcommand{\splitRecoAMVF}{22.1}
\newcommand{\fakeRecoAMVF}{0.07}
\newcommand{\clnEvtAMVF}{34.4}
\newcommand{\mrgEvtAMVF}{41.7}
\newcommand{\splEvtAMVF}{21.7}
\newcommand{\fakEvtAMVF}{0.07}
\newcommand{\hsSelAMVF}{92.92}
\newcommand{\hsSelAMVFrule}{89.11}
\newcommand{\hsSelAMVFplain}{78.80}
\newcommand{\hsSelAMVFincl}{95.14}
\newcommand{\hsResAMVF}{10.24}
\newcommand{\hsResAMVFErr}{0.27}
\newcommand{\hsCleanAMVF}{92.5}
\newcommand{\hsTrkFracAMVF}{74.1}
% PV-Finder + GNN: TBD
\newcommand{\cleanTruthGNN}{\TBD}     % HLDOC | CH clean_per_truth/gnn | seen 66.72
\newcommand{\cleanTruthGNNErr}{\TBD}  % HLDOC | CH clean_per_truth/gnn_err | seen 0.13
\newcommand{\cleanTruthGNNincl}{\TBD} % HLDOC | CHI clean_per_truth/gnn | seen 73.09
\newcommand{\cleanRecoGNN}{\TBD}      % HLDOC | CH clean_per_reco/gnn | seen 79.03
\newcommand{\mergedRecoGNN}{\TBD}     % HLDOC | CH merged_per_reco/gnn | seen 13.55
\newcommand{\splitRecoGNN}{\TBD}      % HLDOC | CH split_per_reco/gnn | seen 7.24
\newcommand{\fakeRecoGNN}{\TBD}       % HLDOC | CH fake_per_reco/gnn | seen 0.18
\newcommand{\clnEvtGNN}{\TBD}         % HLDOC | SUP classes, PU200, gnn unweighted, Clean/event | seen 74.28
\newcommand{\mrgEvtGNN}{\TBD}         % HLDOC | SUP Merged/event | seen 12.73
\newcommand{\splEvtGNN}{\TBD}         % HLDOC | SUP Split/event | seen 6.81
\newcommand{\fakEvtGNN}{\TBD}         % HLDOC | SUP Fake/event | seen 0.17
\newcommand{\clnEvtGNNw}{\TBD}        % HLDOC | SUP gnn, GNN-score weights, Clean/event | seen 74.34
\newcommand{\clnEvtAMVFfit}{\TBD}     % HLDOC | SUP amvf fit tracks + fit weights, Clean/event | seen 46.26
\newcommand{\mrgEvtAMVFfit}{\TBD}     % HLDOC | SUP | seen 31.98
\newcommand{\splEvtAMVFfit}{\TBD}     % HLDOC | SUP | seen 19.56
\newcommand{\cleanLeadFit}{\TBD}      % HLDOC | SUP Clean ratio GNN/AMVF with AMVF fit weights | seen 1.61
\newcommand{\cleanLeadTTVA}{\TBD}     % HLDOC | SUP Clean ratio GNN/AMVF with TTVA links | seen 2.16
\newcommand{\emptyDropped}{\TBD}      % HLDOC | CH table4/empty_candidates_dropped | seen 13,204
\newcommand{\hsSelGNN}{\TBD}          % HLDOC | CH hs_selection/gnn_rule | seen 94.95
\newcommand{\hsSelGNNErr}{\TBD}       % HLDOC | CH hs_selection/gnn_rule_err | seen 0.53
\newcommand{\hsSelDiff}{\TBD}         % HLDOC | CH gnn_rule - amvf_stored [pt] | seen +2.03
\newcommand{\hsSelDiffErr}{\TBD}      % HLDOC | CH hs_selection/gnn_minus_stored_err | seen 0.56
\newcommand{\hsSelGNNincl}{\TBD}      % HLDOC | CHI hs_selection/gnn_rule | seen 96.94
\newcommand{\hsResGNN}{\TBD}          % HLDOC | CH hs_resolution_um/gnn | seen 9.48
\newcommand{\hsResGNNErr}{\TBD}       % HLDOC | CH hs_resolution_um/gnn_err | seen 0.28
\newcommand{\hsResN}{\TBD}            % HLDOC | CH hs_resolution_um/n_common | seen 1,915
\newcommand{\hsResGNNincl}{\TBD}      % HLDOC | CHI hs_resolution_um/gnn | seen 9.57
\newcommand{\hsCleanGNN}{\TBD}        % HLDOC | CH hs_clean/gnn | seen 94.7
\newcommand{\hsTrkFracGNN}{\TBD}      % HLDOC | CH hs_track_frac/gnn | seen 80.9
\newcommand{\trkEffGNN}{\TBD}         % HLDOC | CH trk_eff/gnn | seen 57.45
\newcommand{\trkPurGNN}{\TBD}         % HLDOC | CH trk_pur/gnn | seen 85.87
\newcommand{\trkFOneGNN}{\TBD}        % HLDOC | CH trk_f1/gnn | seen 68.84
\newcommand{\trkEffAMVF}{\TBD}        % HLDOC | CH trk_eff/amvf | seen 68.32
\newcommand{\trkPurAMVF}{\TBD}        % HLDOC | CH trk_pur/amvf | seen 68.65
\newcommand{\trkFOneAMVF}{\TBD}       % HLDOC | CH trk_f1/amvf | seen 68.48
\newcommand{\wrongGNN}{\TBD}          % HLDOC | CH trk_wrong_rate/gnn | seen 14.1
\newcommand{\wrongAMVF}{\TBD}         % HLDOC | CH trk_wrong_rate/amvf | seen 31.4
\newcommand{\assignedGNN}{\TBD}       % HLDOC | CH assigned_frac/gnn | seen 66.5
\newcommand{\lostUnassigned}{\TBD}    % HLDOC | CX lost-unassigned fraction, PU200, gnn_t
\newcommand{\unlabelledLeft}{\TBD}    % HLDOC | CX unlabelled tracks left unassigned, PU200, gnn_t
\newcommand{\eqFourBest}{\TBD}        % HLDOC | CX Eq.(4) best Clean/truth, PU200 | seen 33.5 (s_max = 1)
\newcommand{\eqFourThree}{\TBD}       % HLDOC | CX Eq.(4) s_max = 3 | seen 32.9
\newcommand{\zNear}{\TBD}             % HLDOC | CX z0_nearest Clean/truth | seen 31.9
\newcommand{\zThree}{\TBD}            % HLDOC | CX z0_3sigma Clean/truth | seen 33.0
\newcommand{\truthGuided}{\TBD}       % HLDOC | CX truth_guided Clean/truth, PU200
\newcommand{\finderCap}{\TBD}         % HLDOC | CX finder cap (0.17 mm), PU200
\newcommand{\oneTrkVtx}{\TBD}         % HLDOC | CX GNN vertices with one track, PU200 | seen 36,932 of 180,460
\newcommand{\cleanTruthDomTwo}{\TBD}  % HLDOC | CX Clean/truth, dominant n_trk >= 2 | seen 63.0
\newcommand{\cleanTruthMinTwo}{\TBD}  % HLDOC | CX Clean/truth, >= 2 tracks | seen 53.9
\newcommand{\sigmaPVmed}{\TBD}        % HLDOC | CX median peak width at PU200 [mm] | seen 0.072
\newcommand{\hsRecoDensMin}{\TBD}     % HLDOC | SUP hs_vs_density: lowest HS reconstruction in any bin [%] | seen 99.3
\newcommand{\hsSelDensLo}{\TBD}       % HLDOC | SUP hs_vs_density: smallest selection lead [pt] | seen +2.7
\newcommand{\hsSelDensHi}{\TBD}       % HLDOC | SUP largest selection lead [pt] | seen +5.0
\newcommand{\hsResDensLo}{\TBD}       % HLDOC | SUP smallest resolution gain [um] | seen 0.2
\newcommand{\hsResDensHi}{\TBD}       % HLDOC | SUP largest resolution gain [um] | seen 0.9

%======================================================================
% 5. CPU timing
%======================================================================
% quoted from the note (sections/04b_timing.tex, tab:athena-timing-breakdown): final
\newcommand{\athCPUms}{2401}
\newcommand{\athGPUms}{76}
\newcommand{\athCPUinf}{2189}
\newcommand{\athOther}{31}
\newcommand{\athGNN}{182}
\newcommand{\athAMVF}{650}
\newcommand{\athAMVFlo}{600}
\newcommand{\athAMVFhi}{750}
% 9 Nov measurement (docs/evaluation/cpu_small_model_2026-11-09.md): final
\newcommand{\spdSEightNov}{9.0}
% frozen-chain timing (HL-LHC agent): TBD
\newcommand{\cpuSEightPack}{\TBD}  % HLDOC | CPU density s8 packed ms/event | seen 174
\newcommand{\cpuVsixPad}{\TBD}     % HLDOC | CPU density v6 padded ms/event | seen 1892
\newcommand{\spdPack}{\TBD}        % HLDOC | CPU s8 packed / v6 padded speed-up | seen 10.9
\newcommand{\spdPad}{\TBD}         % HLDOC | CPU s8 padded speed-up | seen 6.2
\newcommand{\cpuGNN}{\TBD}         % HLDOC | CPU GNN v2 forward ms/event on sneezy | seen 97.9
\newcommand{\projChain}{\TBD}      % HLDOC | CPU projected Athena chain, s8 packed [ms] | seen 413
\newcommand{\projChainPad}{\TBD}   % HLDOC | CPU projected chain, s8 padded [ms] | seen 566
\newcommand{\projWorst}{\TBD}      % HLDOC | CPU worst case [ms] | seen 531
\newcommand{\projFactor}{\TBD}     % HLDOC | CPU AMVF / chain at 650 ms | seen 1.57

%======================================================================
% 6. Decision studies (committed docs, final)
%======================================================================
% s8 vs v6, docs/evaluation/model_freeze_candidates_2026-11-12.md
\newcommand{\sEightDEffA}{$-0.15$}
\newcommand{\sEightDEffB}{$-0.08$}
\newcommand{\sEightDSig}{$+1.2\pm1.1$}
\newcommand{\sEightDpse}{$+0.6$}
\newcommand{\sFiveDSig}{$+9$ to $+12$}
% after the GNN, docs/evaluation/gnn_v2_freeze_2026-11-12.md Sec. 7 (t = 0.97, inclusive)
\newcommand{\sEightDClean}{$-0.14\pm0.02$}
\newcommand{\sEightDHSsel}{$-0.10\pm0.06$}
\newcommand{\sEightDHSres}{$+0.04\pm0.04$}
% GNN v2 checkpoint and t options (v6 peaks, test r16443), gnn_v2_freeze_2026-11-12.md
\newcommand{\gnnValLoss}{0.255807}
\newcommand{\optAclean}{73.83}
\newcommand{\optDclean}{73.23}
\newcommand{\optBclean}{68.70}
\newcommand{\optCclean}{65.84}
\newcommand{\optEclean}{56.79}
\newcommand{\optAcleanPU}{67.73}
\newcommand{\optDcleanPU}{66.99}
\newcommand{\optAcleanLow}{84.22}
\newcommand{\optDcleanLow}{86.45}
\newcommand{\optAFone}{72.09}
\newcommand{\optDFone}{74.45}
\newcommand{\optAMVFFone}{74.43}
\newcommand{\optDHSsel}{97.04}
\newcommand{\optAHSsel}{96.79}
\newcommand{\optDHSres}{9.53}
\newcommand{\optAHSres}{9.70}
% HS selection rule, docs/evaluation/hs_selection_identical_rule_2026-11-11.md
\newcommand{\chiCut}{18.42}
\newcommand{\ruleAgreeAll}{96.2}
\newcommand{\ruleAgreePU}{94.8}
\newcommand{\plainAgreeAll}{86.8}
% HS resolution study (GNN v4, v6 peaks), docs/evaluation/itk_hs_resolution_2026-11-09.md
\newcommand{\hsResPeak}{12.38}
\newcommand{\hsResTruthTrk}{9.09}
\newcommand{\hsResAMVFlist}{9.50}
\newcommand{\hsResNoGNN}{9.28}
\newcommand{\hsResOneD}{11.99}
% target readout, docs/evaluation/run2_vs_run4_resolution_2026-11-09.md (2026-11-11)
\newcommand{\tgtSigRunTwo}{0.203}
\newcommand{\pvfSigRunTwo}{0.410}
\newcommand{\amvfSigRunTwo}{0.789}
\newcommand{\tgtSigHL}{0.145}
\newcommand{\tgtSigHLwone}{0.137}
\newcommand{\tgtSigHLfine}{0.124}
\newcommand{\pvfSigHLvsix}{0.218}
\newcommand{\pvfSigHLrunTwoCfg}{0.241}
\newcommand{\tgtHSum}{0.77}
\newcommand{\pvfHSum}{12.4}
\newcommand{\idealPURunTwo}{75.6}
\newcommand{\idealPUHL}{44.4}
\newcommand{\idealHSHL}{6.0}

%======================================================================
% 7. Low pile-up (LPDOC; Run 2 MC 2,550 events, data 2,500 each)
%======================================================================
\newcommand{\lpEffPVF}{\TBD}       % LPDOC R1 | VF mc/pvf/recovery | seen 81.55
\newcommand{\lpEffPVFErr}{\TBD}    % LPDOC R1 | VF mc/pvf/recovery_err | seen 0.14
\newcommand{\lpEffAMVF}{\TBD}      % LPDOC R1 | VF mc/amvf/recovery | seen 75.74
\newcommand{\lpEffDiff}{\TBD}      % LPDOC R1 | VF mc/delta_pvf_minus_amvf/recovery/delta | seen +5.81
\newcommand{\lpFakeFracPVF}{\TBD}  % LPDOC R1 | VF mc/pvf/fake_fraction/value | seen 9.86
\newcommand{\lpFakeFracAMVF}{\TBD} % LPDOC R1 | VF mc/amvf/fake_fraction/value | seen 11.06
\newcommand{\lpFakeEvtPVF}{\TBD}   % LPDOC R1 | VF mc/pvf/no_ref_per_evt | seen 2.363
\newcommand{\lpFakeEvtAMVF}{\TBD}  % LPDOC R1 | VF mc/amvf/no_ref_per_evt | seen 2.487
\newcommand{\lpSigPVF}{\TBD}       % LPDOC R1 | VF mc/pvf/sigma_vtx_vtx_mm/value | seen 0.392
\newcommand{\lpSigPVFErr}{\TBD}    % LPDOC R1 | VF mc/pvf/sigma_vtx_vtx_mm/err | seen 0.009
\newcommand{\lpSigAMVF}{\TBD}      % LPDOC R1 | VF mc/amvf/sigma_vtx_vtx_mm/value | seen 0.789
\newcommand{\lpSigAMVFErr}{\TBD}   % LPDOC R1 | VF mc/amvf/sigma_vtx_vtx_mm/err | seen 0.015
\newcommand{\lpPsePVF}{\TBD}       % LPDOC R1 | VF mc/pvf/p68_residual_um | seen 72.5
\newcommand{\lpPseAMVF}{\TBD}      % LPDOC R1 | VF mc/amvf/p68_residual_um | seen 66.8
\newcommand{\lpWinRatio}{\TBD}     % LPDOC R1 | VF mc/pvf/window_over_p68 | seen 2.35
\newcommand{\lpHSPVF}{\TBD}        % LPDOC R1 | VF mc/pvf/hs_eff/eff | seen 98.94
\newcommand{\lpHSAMVF}{\TBD}       % LPDOC R1 | VF mc/amvf/hs_eff/eff | seen 99.29
\newcommand{\lpRecRunTwo}{\TBD}    % LPDOC R3 | VF run2/pvf/recovery | seen 94.39
\newcommand{\lpRecRunThree}{\TBD}  % LPDOC R3 | VF run3/pvf/recovery | seen 94.03
\newcommand{\lpNoRefRunTwo}{\TBD}  % LPDOC R3 | VF run2/pvf/no_ref_per_evt | seen 3.006
\newcommand{\lpNoRefRunThree}{\TBD}% LPDOC R3 | VF run3/pvf/no_ref_per_evt | seen 3.192
\newcommand{\lpRecMCref}{\TBD}     % LPDOC R3 | VF mc/pvf_vs_amvf_reference/recovery | seen 93.09
\newcommand{\lpSigRunTwo}{\TBD}    % LPDOC R3 | VF run2/pvf/sigma_vtx_vtx_mm/value | seen 0.356
\newcommand{\lpSigRunThree}{\TBD}  % LPDOC R3 | VF run3/pvf/sigma_vtx_vtx_mm/value | seen 0.333
\newcommand{\lpSigAMVFRunTwo}{\TBD}  % LPDOC R3 | VF run2/amvf_ref/sigma_vtx_vtx_mm/value | seen 0.808
\newcommand{\lpSigAMVFRunThree}{\TBD}% LPDOC R3 | VF run3/amvf_ref/sigma_vtx_vtx_mm/value | seen 0.813
\newcommand{\lpCleanTruth}{\TBD}   % LPDOC R2 | GM t0.97/clean_per_truth/value | seen 72.84
\newcommand{\lpCleanTruthAMVF}{\TBD}% LPDOC R2 | GM amvf/clean_per_truth/value | seen 60.91
\newcommand{\lpClean}{\TBD}        % LPDOC R2 | GM t0.97/clean_per_reco/value | seen 86.90
\newcommand{\lpMerged}{\TBD}       % LPDOC R2 | GM t0.97/merged_per_reco/value | seen 11.05
\newcommand{\lpSplit}{\TBD}        % LPDOC R2 | GM t0.97/split_per_reco/value | seen 2.04
\newcommand{\lpFake}{\TBD}         % LPDOC R2 | GM t0.97/fake_per_reco/value | seen 0.003
\newcommand{\lpCleanAMVF}{\TBD}    % LPDOC R2 | GM amvf/clean_per_reco/value | seen 77.23
\newcommand{\lpMergedAMVF}{\TBD}   % LPDOC R2 | GM amvf/merged_per_reco/value | seen 20.19
\newcommand{\lpSplitAMVF}{\TBD}    % LPDOC R2 | GM amvf/split_per_reco/value | seen 2.53
\newcommand{\lpFakeAMVF}{\TBD}     % LPDOC R2 | GM amvf/fake_per_reco/value | seen 0.054
\newcommand{\lpFOne}{\TBD}         % LPDOC R2 | GM t0.97/trk_f1/value | seen 86.95
\newcommand{\lpFOneAMVF}{\TBD}     % LPDOC R2 | GM amvf/trk_f1/value | seen 84.89
\newcommand{\lpHSID}{\TBD}         % LPDOC R2 | GM t0.97/hs_id/value | seen 99.06
\newcommand{\lpHSIDAMVF}{\TBD}     % LPDOC R2 | GM amvf/hs_id/value | seen 98.75
\newcommand{\lpCleanTruthNN}{\TBD} % LPDOC Sec. 3 | t = 0.99 Clean/truth | seen 74.35
\newcommand{\lpEmpty}{\TBD}        % LPDOC R2 | GM t0.97/totals/n_empty | seen 593
\newcommand{\lpRthreeAssigned}{\TBD}% LPDOC R4 | R3A gnn_assigned_fraction | seen 92.38
\newcommand{\lpRthreeAgree}{\TBD}  % LPDOC R4 | R3A agreement/1.0/vs_both_assigned | seen 96.03
\newcommand{\lpRthreeAgreeW}{\TBD} % LPDOC R4 | R3A agreement/0.17/vs_both_assigned | seen 95.09
\newcommand{\lpRthreeLead}{\TBD}   % LPDOC R4 | R3A leading_vertex_agreement | seen 91.20
\newcommand{\lpCentroidShift}{\TBD}% LPDOC Sec. 4 | L/summary/diag_centroid_vs_mean.json eff difference [pt] | seen -0.04

% low-mu behaviour of GNN v2 (gnn_v2_freeze_2026-11-12.md Sec. 6, v6 peaks, validation): final
\newcommand{\lmTopOneLow}{97}
\newcommand{\lmTopOneHigh}{71}
\newcommand{\lmBelowNN}{24--34}
\newcommand{\lmSmallClean}{92}
\newcommand{\lmSmallCleanNN}{61}
\newcommand{\lmSmallBelowLow}{47}
\newcommand{\lmSmallBelowMid}{25}
\newcommand{\lmSmallBelowHigh}{41}
\newcommand{\lmTrackShare}{2}
\newcommand{\lmVoneBelow}{47}
\newcommand{\lmVtwoBelow}{33}
\newcommand{\lmSplitHS}{60}
\newcommand{\lmSplitOne}{63}
\newcommand{\lmSplitDist}{0.8}
\newcommand{\lmFakeMax}{0.4}
\newcommand{\lmCleanPeakT}{\text{0.90--0.92}}
% chain_extras finding (HL-LHC agent): GNN vs AMVF track F1 at mu 0-10
\newcommand{\lmFOneGNN}{\TBD}  % HLDOC | CX trk_f1 vs mu, bin 0-10, gnn_t | seen 87.3
\newcommand{\lmFOneAMVF}{\TBD} % HLDOC | CX trk_f1 vs mu, bin 0-10, amvf | seen 96.6
% window audit (matching_window_audit_2026-11-09.md): final
\newcommand{\greedyBias}{10}
% resolution vs mu study (v6, resolution_vs_mu_2026-11-11.md): final
\newcommand{\rvmChiLo}{4}
\newcommand{\rvmChiHi}{57}
\newcommand{\rvmExcessPVF}{1.61}
\newcommand{\rvmExcessAMVF}{1.60}
\newcommand{\rvmExcessTruth}{1.01}
\newcommand{\optAMVFcleanLow}{77.06}
\newcommand{\optAMVFHSres}{9.93}
\newcommand{\optDvsAclean}{0.6}
\newcommand{\optDvsAlow}{2.2}
\newcommand{\optDvsAFone}{2.4}
\newcommand{\optDvsAres}{0.17}
\newcommand{\hsResAMVFstudy}{10.25}

\imgsrc{mu_training}{outputs/\allowbreak{}11\_\allowbreak{}11\_\allowbreak{}2026\_\allowbreak{}output/\allowbreak{}datasets/\allowbreak{}mu\_\allowbreak{}training.pdf}{pdf}
\imgsrc{mu_validation}{outputs/\allowbreak{}11\_\allowbreak{}11\_\allowbreak{}2026\_\allowbreak{}output/\allowbreak{}datasets/\allowbreak{}mu\_\allowbreak{}validation.pdf}{pdf}
\imgsrc{mu_test}{outputs/\allowbreak{}11\_\allowbreak{}11\_\allowbreak{}2026\_\allowbreak{}output/\allowbreak{}datasets/\allowbreak{}mu\_\allowbreak{}test.pdf}{pdf}
\imgsrc{mu_r16443}{outputs/\allowbreak{}11\_\allowbreak{}11\_\allowbreak{}2026\_\allowbreak{}output/\allowbreak{}datasets/\allowbreak{}mu\_\allowbreak{}held-out\_\allowbreak{}601229\_\allowbreak{}r16443.pdf}{pdf}
\imgsrc{mu_r16638}{outputs/\allowbreak{}11\_\allowbreak{}11\_\allowbreak{}2026\_\allowbreak{}output/\allowbreak{}datasets/\allowbreak{}mu\_\allowbreak{}held-out\_\allowbreak{}601229\_\allowbreak{}r16638.pdf}{pdf}
\imgsrc{lowpu_mu}{outputs/\allowbreak{}11\_\allowbreak{}12\_\allowbreak{}2026\_\allowbreak{}output/\allowbreak{}paper\_\allowbreak{}deck/\allowbreak{}lowpu\_\allowbreak{}mu/\allowbreak{}lowpu\_\allowbreak{}mu\_\allowbreak{}distributions.pdf}{pdf}
\imgsrc{window_decision}{presentations/\allowbreak{}mattia/\allowbreak{}11\_\allowbreak{}13\_\allowbreak{}2026\_\allowbreak{}EB/\allowbreak{}images/\allowbreak{}window\_\allowbreak{}residuals.png}{png}
\imgsrc{eff_vs_mu}{outputs/\allowbreak{}11\_\allowbreak{}12\_\allowbreak{}2026\_\allowbreak{}output/\allowbreak{}hllhc\_\allowbreak{}final/\allowbreak{}chain\_\allowbreak{}extras/\allowbreak{}plots/\allowbreak{}eff\_\allowbreak{}pregnn\_\allowbreak{}vs\_\allowbreak{}mu.pdf}{pdf}
\imgsrc{fake_fraction_vs_mu}{outputs/\allowbreak{}11\_\allowbreak{}12\_\allowbreak{}2026\_\allowbreak{}output/\allowbreak{}hllhc\_\allowbreak{}final/\allowbreak{}h1\_\allowbreak{}vertex\_\allowbreak{}finding/\allowbreak{}plots/\allowbreak{}fake\_\allowbreak{}fraction\_\allowbreak{}vs\_\allowbreak{}mu\_\allowbreak{}r16443}{}
\imgsrc{fakes_per_event_vs_mu}{outputs/\allowbreak{}11\_\allowbreak{}12\_\allowbreak{}2026\_\allowbreak{}output/\allowbreak{}hllhc\_\allowbreak{}final/\allowbreak{}h1\_\allowbreak{}vertex\_\allowbreak{}finding/\allowbreak{}plots/\allowbreak{}fakes\_\allowbreak{}per\_\allowbreak{}event\_\allowbreak{}vs\_\allowbreak{}mu\_\allowbreak{}r16443}{}
\imgsrc{eff_one_to_one_vs_mu}{outputs/\allowbreak{}11\_\allowbreak{}12\_\allowbreak{}2026\_\allowbreak{}output/\allowbreak{}hllhc\_\allowbreak{}final/\allowbreak{}h1\_\allowbreak{}vertex\_\allowbreak{}finding/\allowbreak{}plots/\allowbreak{}eff\_\allowbreak{}one\_\allowbreak{}to\_\allowbreak{}one\_\allowbreak{}vs\_\allowbreak{}mu\_\allowbreak{}r16443}{}
\imgsrc{splits_per_event_vs_mu}{outputs/\allowbreak{}11\_\allowbreak{}12\_\allowbreak{}2026\_\allowbreak{}output/\allowbreak{}hllhc\_\allowbreak{}final/\allowbreak{}h1\_\allowbreak{}vertex\_\allowbreak{}finding/\allowbreak{}plots/\allowbreak{}splits\_\allowbreak{}per\_\allowbreak{}event\_\allowbreak{}vs\_\allowbreak{}mu\_\allowbreak{}r16443}{}
\imgsrc{yield_vs_mu}{outputs/\allowbreak{}11\_\allowbreak{}12\_\allowbreak{}2026\_\allowbreak{}output/\allowbreak{}hllhc\_\allowbreak{}final/\allowbreak{}h1\_\allowbreak{}vertex\_\allowbreak{}finding/\allowbreak{}plots/\allowbreak{}fig6\_\allowbreak{}vertices\_\allowbreak{}vs\_\allowbreak{}mu\_\allowbreak{}r16443}{}
\imgsrc{amvf_pos_categories_vs_mu}{outputs/\allowbreak{}11\_\allowbreak{}12\_\allowbreak{}2026\_\allowbreak{}output/\allowbreak{}hllhc\_\allowbreak{}final/\allowbreak{}h1\_\allowbreak{}vertex\_\allowbreak{}finding/\allowbreak{}plots/\allowbreak{}fig6\_\allowbreak{}categories\_\allowbreak{}vs\_\allowbreak{}mu\_\allowbreak{}amvf\_\allowbreak{}r16443}{}
\imgsrc{deltaz_pu200}{outputs/\allowbreak{}11\_\allowbreak{}12\_\allowbreak{}2026\_\allowbreak{}output/\allowbreak{}hllhc\_\allowbreak{}final/\allowbreak{}h1\_\allowbreak{}vertex\_\allowbreak{}finding/\allowbreak{}plots/\allowbreak{}paper\_\allowbreak{}figures\_\allowbreak{}r16443/\allowbreak{}fig4\_\allowbreak{}hllhc\_\allowbreak{}deltaz}{}
\imgsrc{pos_categories_pu200}{outputs/\allowbreak{}11\_\allowbreak{}12\_\allowbreak{}2026\_\allowbreak{}output/\allowbreak{}hllhc\_\allowbreak{}final/\allowbreak{}h1\_\allowbreak{}vertex\_\allowbreak{}finding/\allowbreak{}plots/\allowbreak{}paper\_\allowbreak{}figures\_\allowbreak{}r16443/\allowbreak{}fig3a\_\allowbreak{}hllhc\_\allowbreak{}categories\_\allowbreak{}pu200}{}
\imgsrc{residuals_pu200}{outputs/\allowbreak{}11\_\allowbreak{}12\_\allowbreak{}2026\_\allowbreak{}output/\allowbreak{}hllhc\_\allowbreak{}final/\allowbreak{}h1\_\allowbreak{}vertex\_\allowbreak{}finding/\allowbreak{}plots/\allowbreak{}fig7b\_\allowbreak{}residuals\_\allowbreak{}r16443}{}
\imgsrc{window_justification}{outputs/\allowbreak{}11\_\allowbreak{}12\_\allowbreak{}2026\_\allowbreak{}output/\allowbreak{}hllhc\_\allowbreak{}final/\allowbreak{}h1\_\allowbreak{}vertex\_\allowbreak{}finding/\allowbreak{}plots/\allowbreak{}window\_\allowbreak{}justification\_\allowbreak{}r16443}{}
\imgsrc{window_scan}{outputs/\allowbreak{}11\_\allowbreak{}12\_\allowbreak{}2026\_\allowbreak{}output/\allowbreak{}hllhc\_\allowbreak{}final/\allowbreak{}h1\_\allowbreak{}vertex\_\allowbreak{}finding/\allowbreak{}plots/\allowbreak{}window\_\allowbreak{}scan\_\allowbreak{}r16443}{}
\imgsrc{p68_vs_mu}{outputs/\allowbreak{}11\_\allowbreak{}12\_\allowbreak{}2026\_\allowbreak{}output/\allowbreak{}hllhc\_\allowbreak{}final/\allowbreak{}resolution\_\allowbreak{}vs\_\allowbreak{}mu/\allowbreak{}resolution\_\allowbreak{}vs\_\allowbreak{}mu\_\allowbreak{}p68}{}
\imgsrc{resolution_vs_mu}{outputs/\allowbreak{}11\_\allowbreak{}12\_\allowbreak{}2026\_\allowbreak{}output/\allowbreak{}hllhc\_\allowbreak{}final/\allowbreak{}resolution\_\allowbreak{}vs\_\allowbreak{}mu/\allowbreak{}resolution\_\allowbreak{}vs\_\allowbreak{}mu}{}
\imgsrc{hs_reco_vs_mu}{outputs/\allowbreak{}11\_\allowbreak{}12\_\allowbreak{}2026\_\allowbreak{}output/\allowbreak{}hllhc\_\allowbreak{}final/\allowbreak{}chain/\allowbreak{}plots/\allowbreak{}hs\_\allowbreak{}reco\_\allowbreak{}eff\_\allowbreak{}vs\_\allowbreak{}mu.pdf}{pdf}
\imgsrc{eff_ntrk_pre_pu200}{outputs/\allowbreak{}11\_\allowbreak{}12\_\allowbreak{}2026\_\allowbreak{}output/\allowbreak{}hllhc\_\allowbreak{}final/\allowbreak{}chain\_\allowbreak{}extras/\allowbreak{}plots/\allowbreak{}eff\_\allowbreak{}vs\_\allowbreak{}n\_\allowbreak{}trk\_\allowbreak{}pregnn\_\allowbreak{}pu200.pdf}{pdf}
\imgsrc{eff_density_pre_pu200}{outputs/\allowbreak{}11\_\allowbreak{}12\_\allowbreak{}2026\_\allowbreak{}output/\allowbreak{}hllhc\_\allowbreak{}final/\allowbreak{}chain\_\allowbreak{}extras/\allowbreak{}plots/\allowbreak{}eff\_\allowbreak{}vs\_\allowbreak{}density\_\allowbreak{}pregnn\_\allowbreak{}pu200.pdf}{pdf}
\imgsrc{eff_dz_pre_pu200}{outputs/\allowbreak{}11\_\allowbreak{}12\_\allowbreak{}2026\_\allowbreak{}output/\allowbreak{}hllhc\_\allowbreak{}final/\allowbreak{}chain\_\allowbreak{}extras/\allowbreak{}plots/\allowbreak{}eff\_\allowbreak{}vs\_\allowbreak{}dz\_\allowbreak{}nearest\_\allowbreak{}pregnn\_\allowbreak{}pu200.pdf}{pdf}
\imgsrc{classes_pu200}{outputs/\allowbreak{}11\_\allowbreak{}12\_\allowbreak{}2026\_\allowbreak{}output/\allowbreak{}hllhc\_\allowbreak{}final/\allowbreak{}supplementary/\allowbreak{}plots/\allowbreak{}main/\allowbreak{}classes\_\allowbreak{}per\_\allowbreak{}event\_\allowbreak{}pu200\_\allowbreak{}primary.pdf}{pdf}
\imgsrc{clean_evt_vs_mu}{outputs/\allowbreak{}11\_\allowbreak{}12\_\allowbreak{}2026\_\allowbreak{}output/\allowbreak{}hllhc\_\allowbreak{}final/\allowbreak{}supplementary/\allowbreak{}plots/\allowbreak{}main/\allowbreak{}clean\_\allowbreak{}per\_\allowbreak{}event\_\allowbreak{}vs\_\allowbreak{}mu\_\allowbreak{}primary.pdf}{pdf}
\imgsrc{merged_evt_vs_mu}{outputs/\allowbreak{}11\_\allowbreak{}12\_\allowbreak{}2026\_\allowbreak{}output/\allowbreak{}hllhc\_\allowbreak{}final/\allowbreak{}supplementary/\allowbreak{}plots/\allowbreak{}main/\allowbreak{}merged\_\allowbreak{}per\_\allowbreak{}event\_\allowbreak{}vs\_\allowbreak{}mu\_\allowbreak{}primary.pdf}{pdf}
\imgsrc{split_evt_vs_mu}{outputs/\allowbreak{}11\_\allowbreak{}12\_\allowbreak{}2026\_\allowbreak{}output/\allowbreak{}hllhc\_\allowbreak{}final/\allowbreak{}supplementary/\allowbreak{}plots/\allowbreak{}main/\allowbreak{}split\_\allowbreak{}per\_\allowbreak{}event\_\allowbreak{}vs\_\allowbreak{}mu\_\allowbreak{}primary.pdf}{pdf}
\imgsrc{fake_evt_vs_mu}{outputs/\allowbreak{}11\_\allowbreak{}12\_\allowbreak{}2026\_\allowbreak{}output/\allowbreak{}hllhc\_\allowbreak{}final/\allowbreak{}supplementary/\allowbreak{}plots/\allowbreak{}main/\allowbreak{}fake\_\allowbreak{}per\_\allowbreak{}event\_\allowbreak{}vs\_\allowbreak{}mu\_\allowbreak{}primary.pdf}{pdf}
\imgsrc{clean_truth_vs_mu}{outputs/\allowbreak{}11\_\allowbreak{}12\_\allowbreak{}2026\_\allowbreak{}output/\allowbreak{}hllhc\_\allowbreak{}final/\allowbreak{}chain/\allowbreak{}plots/\allowbreak{}categories\_\allowbreak{}clean\_\allowbreak{}per\_\allowbreak{}truth\_\allowbreak{}vs\_\allowbreak{}mu.pdf}{pdf}
\imgsrc{hs_sel_vs_mu}{outputs/\allowbreak{}11\_\allowbreak{}12\_\allowbreak{}2026\_\allowbreak{}output/\allowbreak{}hllhc\_\allowbreak{}final/\allowbreak{}chain/\allowbreak{}plots/\allowbreak{}hs\_\allowbreak{}selection\_\allowbreak{}vs\_\allowbreak{}mu.pdf}{pdf}
\imgsrc{hs_res_vs_mu}{outputs/\allowbreak{}11\_\allowbreak{}12\_\allowbreak{}2026\_\allowbreak{}output/\allowbreak{}hllhc\_\allowbreak{}final/\allowbreak{}chain/\allowbreak{}plots/\allowbreak{}hs\_\allowbreak{}resolution\_\allowbreak{}vs\_\allowbreak{}mu.pdf}{pdf}
\imgsrc{hs_res_vs_ntrk}{outputs/\allowbreak{}11\_\allowbreak{}12\_\allowbreak{}2026\_\allowbreak{}output/\allowbreak{}hllhc\_\allowbreak{}final/\allowbreak{}chain/\allowbreak{}plots/\allowbreak{}hs\_\allowbreak{}resolution\_\allowbreak{}vs\_\allowbreak{}ntrk}{}
\imgsrc{hs_reco_vs_density}{outputs/\allowbreak{}11\_\allowbreak{}12\_\allowbreak{}2026\_\allowbreak{}output/\allowbreak{}hllhc\_\allowbreak{}final/\allowbreak{}supplementary/\allowbreak{}plots/\allowbreak{}main/\allowbreak{}hs\_\allowbreak{}reco\_\allowbreak{}eff\_\allowbreak{}vs\_\allowbreak{}density\_\allowbreak{}itk.pdf}{pdf}
\imgsrc{hs_sel_vs_density}{outputs/\allowbreak{}11\_\allowbreak{}12\_\allowbreak{}2026\_\allowbreak{}output/\allowbreak{}hllhc\_\allowbreak{}final/\allowbreak{}supplementary/\allowbreak{}plots/\allowbreak{}main/\allowbreak{}hs\_\allowbreak{}sel\_\allowbreak{}eff\_\allowbreak{}vs\_\allowbreak{}density\_\allowbreak{}itk.pdf}{pdf}
\imgsrc{hs_res_vs_density}{outputs/\allowbreak{}11\_\allowbreak{}12\_\allowbreak{}2026\_\allowbreak{}output/\allowbreak{}hllhc\_\allowbreak{}final/\allowbreak{}supplementary/\allowbreak{}plots/\allowbreak{}main/\allowbreak{}hs\_\allowbreak{}resolution\_\allowbreak{}vs\_\allowbreak{}density\_\allowbreak{}itk.pdf}{pdf}
\imgsrc{trk_f1_vs_mu}{outputs/\allowbreak{}11\_\allowbreak{}12\_\allowbreak{}2026\_\allowbreak{}output/\allowbreak{}hllhc\_\allowbreak{}final/\allowbreak{}chain/\allowbreak{}plots/\allowbreak{}track\_\allowbreak{}f1\_\allowbreak{}vs\_\allowbreak{}mu.pdf}{pdf}
\imgsrc{trk_eff_vs_mu}{outputs/\allowbreak{}11\_\allowbreak{}12\_\allowbreak{}2026\_\allowbreak{}output/\allowbreak{}hllhc\_\allowbreak{}final/\allowbreak{}chain/\allowbreak{}plots/\allowbreak{}track\_\allowbreak{}efficiency\_\allowbreak{}vs\_\allowbreak{}mu.pdf}{pdf}
\imgsrc{trk_pur_vs_mu}{outputs/\allowbreak{}11\_\allowbreak{}12\_\allowbreak{}2026\_\allowbreak{}output/\allowbreak{}hllhc\_\allowbreak{}final/\allowbreak{}chain/\allowbreak{}plots/\allowbreak{}track\_\allowbreak{}purity\_\allowbreak{}vs\_\allowbreak{}mu.pdf}{pdf}
\imgsrc{eff_ntrk_post_pu200}{outputs/\allowbreak{}11\_\allowbreak{}12\_\allowbreak{}2026\_\allowbreak{}output/\allowbreak{}hllhc\_\allowbreak{}final/\allowbreak{}chain\_\allowbreak{}extras/\allowbreak{}plots/\allowbreak{}eff\_\allowbreak{}vs\_\allowbreak{}n\_\allowbreak{}trk\_\allowbreak{}postgnn\_\allowbreak{}pu200.pdf}{pdf}
\imgsrc{eff_density_post_pu200}{outputs/\allowbreak{}11\_\allowbreak{}12\_\allowbreak{}2026\_\allowbreak{}output/\allowbreak{}hllhc\_\allowbreak{}final/\allowbreak{}chain\_\allowbreak{}extras/\allowbreak{}plots/\allowbreak{}eff\_\allowbreak{}vs\_\allowbreak{}density\_\allowbreak{}postgnn\_\allowbreak{}pu200.pdf}{pdf}
\imgsrc{eff_dz_post_pu200}{outputs/\allowbreak{}11\_\allowbreak{}12\_\allowbreak{}2026\_\allowbreak{}output/\allowbreak{}hllhc\_\allowbreak{}final/\allowbreak{}chain\_\allowbreak{}extras/\allowbreak{}plots/\allowbreak{}eff\_\allowbreak{}vs\_\allowbreak{}dz\_\allowbreak{}nearest\_\allowbreak{}postgnn\_\allowbreak{}pu200.pdf}{pdf}
\imgsrc{hs_clean_vs_mu}{outputs/\allowbreak{}11\_\allowbreak{}12\_\allowbreak{}2026\_\allowbreak{}output/\allowbreak{}hllhc\_\allowbreak{}final/\allowbreak{}chain/\allowbreak{}plots/\allowbreak{}hs\_\allowbreak{}clean\_\allowbreak{}fraction\_\allowbreak{}vs\_\allowbreak{}mu.pdf}{pdf}
\imgsrc{hs_trkfrac_vs_mu}{outputs/\allowbreak{}11\_\allowbreak{}12\_\allowbreak{}2026\_\allowbreak{}output/\allowbreak{}hllhc\_\allowbreak{}final/\allowbreak{}chain/\allowbreak{}plots/\allowbreak{}hs\_\allowbreak{}track\_\allowbreak{}fraction\_\allowbreak{}vs\_\allowbreak{}mu.pdf}{pdf}
\imgsrc{classes_variants}{outputs/\allowbreak{}11\_\allowbreak{}12\_\allowbreak{}2026\_\allowbreak{}output/\allowbreak{}hllhc\_\allowbreak{}final/\allowbreak{}supplementary/\allowbreak{}plots/\allowbreak{}backup/\allowbreak{}classes\_\allowbreak{}per\_\allowbreak{}event\_\allowbreak{}pu200\_\allowbreak{}all\_\allowbreak{}variants.pdf}{pdf}
\imgsrc{amvf_fit_weights}{outputs/\allowbreak{}11\_\allowbreak{}12\_\allowbreak{}2026\_\allowbreak{}output/\allowbreak{}hllhc\_\allowbreak{}final/\allowbreak{}supplementary/\allowbreak{}plots/\allowbreak{}backup/\allowbreak{}amvf\_\allowbreak{}fit\_\allowbreak{}weight\_\allowbreak{}distribution.pdf}{pdf}
\imgsrc{density_definitions}{outputs/\allowbreak{}11\_\allowbreak{}12\_\allowbreak{}2026\_\allowbreak{}output/\allowbreak{}hllhc\_\allowbreak{}final/\allowbreak{}supplementary/\allowbreak{}plots/\allowbreak{}backup/\allowbreak{}local\_\allowbreak{}density\_\allowbreak{}definitions.pdf}{pdf}
\imgsrc{eq4_vs_mu}{outputs/\allowbreak{}11\_\allowbreak{}12\_\allowbreak{}2026\_\allowbreak{}output/\allowbreak{}hllhc\_\allowbreak{}final/\allowbreak{}chain\_\allowbreak{}extras/\allowbreak{}plots/\allowbreak{}eq4\_\allowbreak{}clean\_\allowbreak{}vs\_\allowbreak{}mu.pdf}{pdf}
\imgsrc{z0_baselines_vs_mu}{outputs/\allowbreak{}11\_\allowbreak{}12\_\allowbreak{}2026\_\allowbreak{}output/\allowbreak{}hllhc\_\allowbreak{}final/\allowbreak{}chain\_\allowbreak{}extras/\allowbreak{}plots/\allowbreak{}z0\_\allowbreak{}baselines\_\allowbreak{}clean\_\allowbreak{}vs\_\allowbreak{}mu.pdf}{pdf}
\imgsrc{truth_guided_vs_mu}{outputs/\allowbreak{}11\_\allowbreak{}12\_\allowbreak{}2026\_\allowbreak{}output/\allowbreak{}hllhc\_\allowbreak{}final/\allowbreak{}chain\_\allowbreak{}extras/\allowbreak{}plots/\allowbreak{}truth\_\allowbreak{}guided\_\allowbreak{}clean\_\allowbreak{}vs\_\allowbreak{}mu.pdf}{pdf}
\imgsrc{finder_cap_vs_mu}{outputs/\allowbreak{}11\_\allowbreak{}12\_\allowbreak{}2026\_\allowbreak{}output/\allowbreak{}hllhc\_\allowbreak{}final/\allowbreak{}chain\_\allowbreak{}extras/\allowbreak{}plots/\allowbreak{}finder\_\allowbreak{}cap\_\allowbreak{}vs\_\allowbreak{}mu.pdf}{pdf}
\imgsrc{stages_pu200}{outputs/\allowbreak{}11\_\allowbreak{}12\_\allowbreak{}2026\_\allowbreak{}output/\allowbreak{}hllhc\_\allowbreak{}final/\allowbreak{}chain\_\allowbreak{}extras/\allowbreak{}plots/\allowbreak{}clean\_\allowbreak{}yield\_\allowbreak{}stages\_\allowbreak{}pu200.pdf}{pdf}
\imgsrc{wrong_rate_vs_mu}{outputs/\allowbreak{}11\_\allowbreak{}12\_\allowbreak{}2026\_\allowbreak{}output/\allowbreak{}hllhc\_\allowbreak{}final/\allowbreak{}chain\_\allowbreak{}extras/\allowbreak{}plots/\allowbreak{}wrong\_\allowbreak{}rate\_\allowbreak{}vs\_\allowbreak{}mu.pdf}{pdf}
\imgsrc{lost_unassigned_vs_mu}{outputs/\allowbreak{}11\_\allowbreak{}12\_\allowbreak{}2026\_\allowbreak{}output/\allowbreak{}hllhc\_\allowbreak{}final/\allowbreak{}chain\_\allowbreak{}extras/\allowbreak{}plots/\allowbreak{}lost\_\allowbreak{}unassigned\_\allowbreak{}frac\_\allowbreak{}vs\_\allowbreak{}mu.pdf}{pdf}
\imgsrc{unlabelled_unassigned_vs_mu}{outputs/\allowbreak{}11\_\allowbreak{}12\_\allowbreak{}2026\_\allowbreak{}output/\allowbreak{}hllhc\_\allowbreak{}final/\allowbreak{}chain\_\allowbreak{}extras/\allowbreak{}plots/\allowbreak{}unlabelled\_\allowbreak{}left\_\allowbreak{}unassigned\_\allowbreak{}vs\_\allowbreak{}mu.pdf}{pdf}
\imgsrc{sigma_pv_pu200}{outputs/\allowbreak{}11\_\allowbreak{}12\_\allowbreak{}2026\_\allowbreak{}output/\allowbreak{}hllhc\_\allowbreak{}final/\allowbreak{}chain\_\allowbreak{}extras/\allowbreak{}plots/\allowbreak{}sigma\_\allowbreak{}pv\_\allowbreak{}peakwidth\_\allowbreak{}vs\_\allowbreak{}param\_\allowbreak{}pu200.pdf}{pdf}
\imgsrc{t_scan_val_s8}{outputs/\allowbreak{}11\_\allowbreak{}12\_\allowbreak{}2026\_\allowbreak{}output/\allowbreak{}gnn\_\allowbreak{}v2\_\allowbreak{}freeze/\allowbreak{}figures/\allowbreak{}t\_\allowbreak{}scan\_\allowbreak{}val\_\allowbreak{}s8\_\allowbreak{}v2e527.pdf}{pdf}
\imgsrc{vs_mu_test_s8}{outputs/\allowbreak{}11\_\allowbreak{}12\_\allowbreak{}2026\_\allowbreak{}output/\allowbreak{}gnn\_\allowbreak{}v2\_\allowbreak{}freeze/\allowbreak{}figures/\allowbreak{}vs\_\allowbreak{}mu\_\allowbreak{}test\_\allowbreak{}s8\_\allowbreak{}v2e527.pdf}{pdf}
\imgsrc{v6_vs_s8_gnn}{outputs/\allowbreak{}11\_\allowbreak{}12\_\allowbreak{}2026\_\allowbreak{}output/\allowbreak{}gnn\_\allowbreak{}v2\_\allowbreak{}freeze/\allowbreak{}figures/\allowbreak{}v6\_\allowbreak{}vs\_\allowbreak{}s8\_\allowbreak{}after\_\allowbreak{}gnn\_\allowbreak{}vs\_\allowbreak{}mu.pdf}{pdf}
\imgsrc{lowmu_displays}{outputs/\allowbreak{}11\_\allowbreak{}12\_\allowbreak{}2026\_\allowbreak{}output/\allowbreak{}gnn\_\allowbreak{}v2\_\allowbreak{}freeze/\allowbreak{}lowmu/\allowbreak{}lowmu\_\allowbreak{}displays\_\allowbreak{}val\_\allowbreak{}v6\_\allowbreak{}v2e527.pdf}{pdf}
\imgsrc{gnn_learning_curve}{outputs/\allowbreak{}11\_\allowbreak{}12\_\allowbreak{}2026\_\allowbreak{}output/\allowbreak{}gnn\_\allowbreak{}v2\_\allowbreak{}freeze/\allowbreak{}checkpoint/\allowbreak{}learning\_\allowbreak{}curve\_\allowbreak{}v1\_\allowbreak{}v2.pdf}{pdf}
\imgsrc{lp_fig8a}{outputs/\allowbreak{}11\_\allowbreak{}12\_\allowbreak{}2026\_\allowbreak{}output/\allowbreak{}lowpu\_\allowbreak{}final/\allowbreak{}summary/\allowbreak{}figures/\allowbreak{}fig8a\_\allowbreak{}run2mc\_\allowbreak{}deltaz.pdf}{pdf}
\imgsrc{lp_fig8b}{outputs/\allowbreak{}11\_\allowbreak{}12\_\allowbreak{}2026\_\allowbreak{}output/\allowbreak{}lowpu\_\allowbreak{}final/\allowbreak{}summary/\allowbreak{}figures/\allowbreak{}fig8b\_\allowbreak{}run2mc\_\allowbreak{}yield\_\allowbreak{}vs\_\allowbreak{}mu.pdf}{pdf}
\imgsrc{lp_fig9a}{outputs/\allowbreak{}11\_\allowbreak{}12\_\allowbreak{}2026\_\allowbreak{}output/\allowbreak{}lowpu\_\allowbreak{}final/\allowbreak{}gnn\_\allowbreak{}mc/\allowbreak{}figures/\allowbreak{}fig9a\_\allowbreak{}run2mc\_\allowbreak{}eff\_\allowbreak{}vs\_\allowbreak{}ntrk.pdf}{pdf}
\imgsrc{lp_fig9b}{outputs/\allowbreak{}11\_\allowbreak{}12\_\allowbreak{}2026\_\allowbreak{}output/\allowbreak{}lowpu\_\allowbreak{}final/\allowbreak{}gnn\_\allowbreak{}mc/\allowbreak{}figures/\allowbreak{}fig9b\_\allowbreak{}run2mc\_\allowbreak{}categories.pdf}{pdf}
\imgsrc{lp_fig10a}{outputs/\allowbreak{}11\_\allowbreak{}12\_\allowbreak{}2026\_\allowbreak{}output/\allowbreak{}lowpu\_\allowbreak{}final/\allowbreak{}summary/\allowbreak{}figures/\allowbreak{}fig10a\_\allowbreak{}run2\_\allowbreak{}deltaz.pdf}{pdf}
\imgsrc{lp_fig10b}{outputs/\allowbreak{}11\_\allowbreak{}12\_\allowbreak{}2026\_\allowbreak{}output/\allowbreak{}lowpu\_\allowbreak{}final/\allowbreak{}summary/\allowbreak{}figures/\allowbreak{}fig10b\_\allowbreak{}run2\_\allowbreak{}yield\_\allowbreak{}vs\_\allowbreak{}mu.pdf}{pdf}
\imgsrc{lp_fig10c}{outputs/\allowbreak{}11\_\allowbreak{}12\_\allowbreak{}2026\_\allowbreak{}output/\allowbreak{}lowpu\_\allowbreak{}final/\allowbreak{}summary/\allowbreak{}figures/\allowbreak{}fig10c\_\allowbreak{}run3\_\allowbreak{}deltaz.pdf}{pdf}
\imgsrc{lp_fig10d}{outputs/\allowbreak{}11\_\allowbreak{}12\_\allowbreak{}2026\_\allowbreak{}output/\allowbreak{}lowpu\_\allowbreak{}final/\allowbreak{}summary/\allowbreak{}figures/\allowbreak{}fig10d\_\allowbreak{}run3\_\allowbreak{}yield\_\allowbreak{}vs\_\allowbreak{}mu.pdf}{pdf}
\imgsrc{lp_t_scan}{outputs/\allowbreak{}11\_\allowbreak{}12\_\allowbreak{}2026\_\allowbreak{}output/\allowbreak{}lowpu\_\allowbreak{}final/\allowbreak{}gnn\_\allowbreak{}mc/\allowbreak{}figures/\allowbreak{}gnn\_\allowbreak{}t\_\allowbreak{}scan\_\allowbreak{}run2mc.pdf}{pdf}
\imgsrc{athena_timing}{presentations/\allowbreak{}mattia/\allowbreak{}11\_\allowbreak{}13\_\allowbreak{}2026\_\allowbreak{}EB/\allowbreak{}images/\allowbreak{}athena\_\allowbreak{}cpu\_\allowbreak{}gpu\_\allowbreak{}amvf.png}{png}

\hypersetup{pdftitle={PV-Finder, paper deck, November 13, 2026}}

\newcommand{\twofig}[4]{% left name, left caption, right name, right caption
\begin{columns}[T]
\begin{column}{.49\linewidth}\centering\plot{#1}{.58}\fcap{#2}\end{column}
\begin{column}{.49\linewidth}\centering\plot{#3}{.58}\fcap{#4}\end{column}
\end{columns}}
\newcommand{\threefig}[6]{%
\begin{columns}[T]
\begin{column}{.32\linewidth}\centering\plot{#1}{.42}\fcap{#2}\end{column}
\begin{column}{.32\linewidth}\centering\plot{#3}{.42}\fcap{#4}\end{column}
\begin{column}{.32\linewidth}\centering\plot{#5}{.42}\fcap{#6}\end{column}
\end{columns}}
\newcommand{\fourfig}[8]{%
\begin{columns}[T]
\begin{column}{.49\linewidth}\centering\plot{#1}{.33}\fcap{#2}\plot{#5}{.33}\fcap{#6}\end{column}
\begin{column}{.49\linewidth}\centering\plot{#3}{.33}\fcap{#4}\plot{#7}{.33}\fcap{#8}\end{column}
\end{columns}}

\begin{document}
\framelogo{}

\begin{frame}[plain]\titlepage\end{frame}

%======================================================================
\begin{frame}{This Deck}
\small
\begin{columns}[T]
\begin{column}{.52\linewidth}
\textbf{For the colleague finishing the paper draft} (\texttt{PVFinderPaper\_11\_11\_2026.pdf}).
\begin{enumerate}\setlength\itemsep{.3em}
\item Frozen configuration, one table per regime
\item Datasets and which file is used for what
\item Metric definitions, with the code that computes each one
\item Answers to the 29 questions in the draft, with page and line
\item Answers to Mattia's own questions
\item Results: every paper figure and number, main and backup
\item Proposed internal-note edits
\end{enumerate}
\end{column}
\begin{column}{.45\linewidth}
\textbf{How to read it}
\begin{itemize}\setlength\itemsep{.3em}
\item Every number is a macro in \texttt{numbers.tex}.
\item \TBD{} marks a number from the final recompute that is not in yet. Its source file and key are in the comment next to the macro.
\item A framed red box marks a figure not produced yet, with its expected path.
\item \mainbadge{} figures are proposed for the paper; \backupbadge{} figures go to the note or the backup.
\item PV-Finder and AMVF are always compared on the same events with the same code.
\end{itemize}
\end{column}
\end{columns}
\end{frame}

%======================================================================
\partdivider{1}{Frozen Configuration}{One configuration per regime, used for every result}

\begin{frame}{HL-LHC Configuration ($0\le\mu\le200$)}
\footnotesize
\begin{columns}[T]
\begin{column}{.6\linewidth}
\scriptsize\renewcommand{\arraystretch}{1.07}
\begin{tabular}{@{}p{.25\linewidth}p{.69\linewidth}@{}}
\toprule
Density model & \hlModel: U-Net v2, \hlUnetCh{} channels, \hlLatent{} latent channel \\
 & MLP $7\to$\hlHidden$\to1000$, \hlParams{} parameters \\
Checkpoint & \texttt{\hlCkpt} \\
\midrule
Peak amplitude $\tau$ & \hlAmp \\
Minimum width & \hlWidth{} bins (0.04 mm each) \\
Minimum integral & \hlIntegral \\
Minimum height & \hlHeight \\
Position & local centroid, $\pm$\hlCentroid{} bins about the maximum \\
\midrule
GNN & full-$\mu$ v2, epoch \hlGnnEpoch, \gnnParams{} parameters \\
Graph & $k=\gnnK$ nearest peaks per track \\
Score cut $t$ & \hlT{} (top-1 edge per track kept if score $\ge t$) \\
Output vertices & candidates with $\ge1$ assigned track \\
\midrule
Matching window & \winMM{} mm (all vertex results, both algorithms) \\
Hard-scatter window & \hsWinMM{} mm \\
$\svtx$ fit & counted peaks, \pairBins{} bins over $\pm6$ mm \\
\bottomrule
\end{tabular}
\end{column}
\begin{column}{.37\linewidth}
\begin{itemize}\setlength\itemsep{.25em}
\item One line in code: \texttt{get\_regime("hllhc")} in \texttt{src/\allowbreak pv\_finder/\allowbreak utils/\allowbreak working\_points.py}.
\item The integral and height cuts are the v6 cuts (0.30 / 0.03) scaled to give \hlModel{} the v6 fake rate.
\item The $\svtx$ fit uses the same peaks as the counts. There is no second threshold.
\item Clean requires $\ge\purityCut$ of a vertex's tracks from one truth vertex.
\item Single-vertex residuals use an optimal one-to-one assignment within \assignMM{} mm. This is a measurement radius, not a matching window.
\end{itemize}
\end{column}
\end{columns}
\end{frame}

\begin{frame}{Low Pile-up Configuration (Run 2 MC, Run 2 and Run 3 Data)}
\footnotesize
\begin{columns}[T]
\begin{column}{.6\linewidth}
\scriptsize\renewcommand{\arraystretch}{1.07}
\begin{tabular}{@{}p{.25\linewidth}p{.69\linewidth}@{}}
\toprule
Density model & \lpModel: U-Net v1, \lpUnetCh{} channels, 1 latent channel \\
 & MLP $7\to$\lpHidden$\to1000$, \lpParams{} parameters \\
Checkpoint & \texttt{\lpCkpt} \\
Training & 50 epochs MLP to histogram, 400 end to end; no KDE \\
\midrule
Peak amplitude $\tau$ & \hlAmp \\
Minimum width & \hlWidth{} bins \\
Minimum integral & \lpIntegral \\
Minimum height & none \\
Position & local centroid, $\pm$\hlCentroid{} bins \\
\midrule
GNN & \texttt{\lpGnnCkpt} \\
Score cut $t$ & \lpT{} (same as HL-LHC) \\
Output vertices & candidates with $\ge1$ assigned track \\
\midrule
Matching window & \winMM{} mm \\
Hard-scatter window & \hsWinMM{} mm \\
$\svtx$ fit & counted peaks, \pairBins{} bins over $\pm6$ mm \\
\bottomrule
\end{tabular}
\end{column}
\begin{column}{.37\linewidth}
\begin{itemize}\setlength\itemsep{.25em}
\item \texttt{get\_regime("lowpu")}, same file.
\item One density model for every low pile-up result. The ep300 reproduction model and the two-stage KDE model are no longer used.
\item The window is the HL-LHC value. At low pile-up it is \lpWinRatio{} times the PV-Finder 68\% residual of Run 2 MC.
\item $t=\lpT$: no other $t$ is better in Clean/truth, HS identification and track $F_1$ together on Run 2 MC.
\item There is no Run 3 MC. The simulation is called Run 2 MC.
\end{itemize}
\end{column}
\end{columns}
\end{frame}

\begin{frame}{What Changed with Respect to the 11 November Draft}
\scriptsize
\renewcommand{\arraystretch}{1.2}
\begin{tabular}{@{}p{.19\linewidth}p{.41\linewidth}p{.33\linewidth}@{}}
\toprule
Item & Draft of 11 November & Final \\
\midrule
HL-LHC density model & v6 (280 channels, 4 latent, 3.55 M parameters) & \hlModel{} (\hlUnetCh{} channels, 1 latent, \hlParams) \\
HL-LHC peak cuts & integral 0.30, height 0.03 & \hlIntegral, \hlHeight{} \\
$\svtx$ threshold & separate peak finding at integral 0.50 & the counted peaks \\
Matching window & 0.1694 mm; 0.5 mm in the chain; fitted $\svtx$ at low $\mu$ & \winMM{} mm everywhere \\
HL-LHC GNN & v4 (PU200 only), epoch 107 & full-$\mu$ v2, epoch \hlGnnEpoch \\
GNN score cut & 0.98 (categories), 0.95 (HS) & \hlT{} for every use \\
Low pile-up density model & ep300 (KDE-initialised), PVF-e400, KDE two-stage & \lpModel{} only \\
Low pile-up GNN peaks & integral 0.2, full-region mean & integral \lpIntegral, local centroid \\
Low pile-up $t$ & 0.99, 0.95 ($F_1$), 0.5 (HS), 0.98 (Run 3) & \lpT{} for every use \\
HS selection & plain $\sum p_\mathrm{T}^2$ & Athena 24.0 rule, same code for both \\
HS position & histogram peak & 3D adaptive fit with beam spot \\
Categories & AMVF TTVA links vs GNN tracks & same, plus weighted cross-checks \\
\bottomrule
\end{tabular}
\end{frame}

%======================================================================
\partdivider{2}{Datasets}{Training, validation and test samples}

\begin{frame}{HL-LHC Files}
\scriptsize
\renewcommand{\arraystretch}{1.05}
\begin{tabular}{@{}lrlp{.36\linewidth}@{}}
\toprule
File (mc21 14 TeV $t\bar t$, e8481\_s4494, ITk) & Events & Pile-up & Role \\
\midrule
SingleLep 601229 r16438 & 99{,}800 & 190--210, $\langle\mu\rangle=200.0$ & density training pool \\
SingleLep 601229 r16633 & 99{,}800 & 190--210 & density training pool \\
all-had 601237 r16633, part 1 & 1{,}421{,}400 & 190--210 & density training pool \\
all-had 601237 r16633, part 2 & 1{,}498{,}600 & 190--210 & density training pool \\
all-had 601237 r16633, part 3 & 1{,}200{,}000 & 190--210 & density training pool \\
all-had 601237 r16633, part 4 & 854{,}000 & 190--210 & density training pool \\
\textbf{PU200 pool} & \textbf{\nPoolEvt} & integers 190--210 & \\
\midrule
SingleLep 601229 r16443 & 100{,}000 & flat 0--200, $\langle\mu\rangle=99.6$ & entries 0--24,999: \textbf{test} of every result \\
 & & & entries $\ge25{,}000$: GNN v2 training and validation \\
SingleLep 601229 r16638 & 99{,}800 & flat 0--200, $\langle\mu\rangle=99.6$ & GNN v2 training (\gnnTrainB) \\
 & & & second held-out file for the density model \\
\bottomrule
\end{tabular}
\par\vspace{.6em}
\begin{itemize}\setlength\itemsep{.25em}
\item The density model never saw r16443 or r16638. The GNN never saw r16443 entries 0--24,999.
\item Track selection: ITk paper Table 6. In-time pile-up only. Pile-up is \texttt{ActualNumOfInt}.
\item PU200 = the \nEvtPU{} r16443 test events with $185\le\mu\le215$ (actual 185--200, $\langle\mu\rangle=\meanMuPU$), \nTruthPU{} truth vertices with $n_\mathrm{trk}\ge2$.
\end{itemize}
\src{docs/\allowbreak{}additional\_\allowbreak{}studies/\allowbreak{}paper\_\allowbreak{}questions\_\allowbreak{}2026-11-11.md Part 3; \allowbreak{}outputs/\allowbreak{}11\_\allowbreak{}11\_\allowbreak{}2026\_\allowbreak{}output/\allowbreak{}datasets/\allowbreak{}hllhc\_\allowbreak{}dataset\_\allowbreak{}inventory.json}
\end{frame}

\begin{frame}{HL-LHC Density Model: Training, Validation and Test}
\footnotesize
\begin{columns}[T]
\begin{column}{.46\linewidth}
{\scriptsize\renewcommand{\arraystretch}{1.15}
\begin{tabular}{@{}lrr@{}}
\toprule
Role & 40 mm intervals & Event equiv. \\
\midrule
Training (96\%) & \nTrainSub & \nTrainEq \\
Validation (3\%) & \nValSub & \nValEq \\
Test (1\%, unused) & \nTestSub & \nTestEq \\
\bottomrule
\end{tabular}\par}
\par\vspace{.5em}
\begin{itemize}\setlength\itemsep{.25em}
\item Seed-42 random permutation of the intervals of the PU200 pool. The roles therefore share its pile-up shape: uniform over 190--210, with $\mu=196$ doubled and $\mu=197$ nearly empty (production feature).
\item Results never use the test intervals. They use the independent files r16443 and r16638.
\end{itemize}
\end{column}
\begin{column}{.52\linewidth}\centering
\plot{mu_training}{.3}\par\vspace{.2em}
\begin{minipage}{.48\linewidth}\centering\plot{mu_validation}{.26}\end{minipage}\hfill
\begin{minipage}{.48\linewidth}\centering\plot{mu_test}{.26}\end{minipage}%
\end{column}
\end{columns}
\end{frame}

\begin{frame}{HL-LHC Held-out Files and the GNN v2 Samples}
\footnotesize
\begin{columns}[T]
\begin{column}{.55\linewidth}
\scriptsize\renewcommand{\arraystretch}{1.08}
\begin{tabular}{@{}lp{.4\linewidth}r@{}}
\toprule
GNN v2 set & Source & Events \\
\midrule
Training & r16443 entries $\ge25{,}000$ & \gnnTrainA \\
 & r16638 (stratified in 10-unit $\mu$ bins) & \gnnTrainB \\
 & total & \gnnTrain \\
Validation ($t$ choice) & r16443 entries $\ge25{,}000$ & \gnnVal \\
\quad monitor (checkpoint) & subset of validation & \gnnMon \\
Reserved test & r16443 entries $\ge25{,}000$ & \gnnResTest \\
Reserved test & r16638 & \gnnResTestB \\
\midrule
Paper test & r16443 entries 0--24,999 & 25{,}000 \\
\bottomrule
\end{tabular}
\par\vspace{.4em}
\begin{itemize}\setlength\itemsep{.2em}
\item Graphs use truth-vertex candidates, so no density model enters the GNN training data.
\item No r16638 truth-$z$ or HS $(x,y,z)$ fingerprint occurs in r16443.
\item Resolution vs $\mu$ (before the GNN): \nEvtResMu{} events of both files, fixed per-bin counts chosen before looking at results.
\end{itemize}
\end{column}
\begin{column}{.42\linewidth}\centering
\plot{mu_r16443}{.33}\par\vspace{.2em}
\plot{mu_r16638}{.33}
\end{column}
\end{columns}
\src{outputs/\allowbreak{}11\_\allowbreak{}11\_\allowbreak{}2026\_\allowbreak{}output/\allowbreak{}gnn\_\allowbreak{}fullmu\_\allowbreak{}v2/\allowbreak{}split\_\allowbreak{}manifest.json; \allowbreak{}docs/\allowbreak{}evaluation/\allowbreak{}gnn\_\allowbreak{}full\_\allowbreak{}pileup\_\allowbreak{}training\_\allowbreak{}v2\_\allowbreak{}2026-11-11.md}
\end{frame}

\begin{frame}{Low Pile-up Samples}
\footnotesize
\begin{columns}[T]
\begin{column}{.54\linewidth}
\scriptsize\renewcommand{\arraystretch}{1.08}
\begin{tabular}{@{}lrl@{}}
\toprule
Sample & Events & Use \\
\midrule
Run 2 MC ($t\bar t$, $\mu$ 0--80) & \lpMCTotal & \\
\quad training (70\%) & \lpMCTrain & density model \\
\quad validation (25\%) & \lpMCVal & density model \\
\quad test (5\%) & \lpMCTest & \textbf{every MC result} \\
\quad non-test (train + val) & \lpMCNonTest & GNN training \\
Run 2 data (ZeroBias, run 364076) & 2{,}500 & data \\
Run 3 data (\texttt{file\_3.root}) & 2{,}500 & data \\
\bottomrule
\end{tabular}
\par\vspace{.4em}
\begin{itemize}\setlength\itemsep{.2em}
\item MC test: h5 events 48,450--50,999 (the last 5\%), \lpTruthMC{} truth vertices with $n_\mathrm{trk}\ge2$, $\langle\mu\rangle=\lpMuMC$.
\item Run 2 data: entries 0--2,499. Run 3 data: the first 2,500 accepted entries of 0--3,158. Vertex finding and track association use the same events.
\item Data have no truth: AMVF vertices ($n_\mathrm{trk}\ge2$, beam-spot corrected) are the reference.
\end{itemize}
\end{column}
\begin{column}{.43\linewidth}\centering
\plot{lowpu_mu}{.5}
\fcap{Pile-up of the three evaluation samples.}
\end{column}
\end{columns}
\end{frame}

%======================================================================
\partdivider{3}{Metric Definitions}{How each number is computed, and the code that computes it}

\newcommand{\mtab}[1]{\scriptsize\renewcommand{\arraystretch}{1.22}%
\begin{tabular}{@{}p{.14\linewidth}p{.54\linewidth}>{\raggedright\arraybackslash\tiny}p{.25\linewidth}@{}}
\toprule
Quantity & Definition & Code \\
\midrule
#1
\bottomrule
\end{tabular}}

\begin{frame}{Before the GNN: Efficiency, Fakes, Splits}
\mtab{%
Efficiency (Eq.~1) & truth vertices ($n_\mathrm{trk}\ge2$) with $\ge1$ reco vertex within \winMM{} mm, over all truth vertices; ratio of sums over events. Equals clean + merged truth of the closest-first matcher (merged credit), event by event. & \texttt{compare\_res\_reco} in \texttt{efficiency\_res\_\allowbreak optimized\_atlas.py}; identity test \texttt{tests/\allowbreak test\_amvf\_fair\_\allowbreak matching.py} \\
One-to-one efficiency & truth vertices that win a dedicated reco vertex in the closest-first assignment & same \\
Fake fraction & reco vertices with no truth vertex within \winMM{} mm, over all reco vertices & same (\texttt{fake}) \\
Fakes per event & the same count over the number of events & same \\
Split & reco vertex with a truth vertex in its window that a closer reco vertex already took & same (\texttt{split}) \\
HS reconstruction & events with a reco vertex within \hsWinMM{} mm of the true HS ($z$ only, ITk Sec.~5.5), over all events & \texttt{gnn/\allowbreak evaluation/\allowbreak hs\_metrics.py} \\
}
\par\vspace{.4em}{\footnotesize Truth: \texttt{TruthVertex} with $\ge2$ selected tracks. AMVF: stored vertices with $n_\mathrm{trk}\ge2$. Errors: event bootstrap (2,000 replicas); PV-Finder minus AMVF is paired on the same resampled events.}
\end{frame}

\begin{frame}{Before the GNN: Resolution and the Window}
\mtab{%
$\svtx$ & all within-event pairs of reco vertices; $\Delta z$ in \pairBins{} bins over $\pm6$ mm; fit $a/(1+e^{b(r-|x|)})+c$; $\svtx=|r|$. Same peaks as the counts. Quoted at PU200, and per $\mu$ bin only for $\mu\ge100$. & \texttt{fit\_dz} in \texttt{run\_eval\_\allowbreak pvf\_run3.py}; \texttt{amvf\_fair\_\allowbreak comparison/\allowbreak scan.py} \\
68\% residual & 68th percentile of $|z_\mathrm{reco}-z_\mathrm{true}|$ after the optimal one-to-one assignment within \assignMM{} mm (most matches, then least total distance). The resolution versus $\mu$. & \texttt{matched\_residuals\_\allowbreak optimal}, \texttt{amvf\_fair\_\allowbreak comparison/\allowbreak matching.py} \\
Window & $3\times$ the larger 68\% residual of the two algorithms at PU200, rounded: $3\times\pseAMVF~\mu$m $=0.1694$ mm $\to$ \winMM{} mm & \texttt{working\_points.\allowbreak MATCH\_WINDOW\_MM} \\
Hard-scatter window & \hsWinMM{} mm, ITk performance paper & \texttt{working\_points.\allowbreak HS\_MATCH\_\allowbreak WINDOW\_MM} \\
}
\end{frame}

\begin{frame}{After the GNN: Vertex Categories}
\mtab{%
Assignment & each track keeps its top-1 GNN edge if the score is $\ge t=\hlT$; candidates with no track are dropped from the output & \texttt{gnn/evaluation/\allowbreak chain\_inference.py} \\
Track lists & PV-Finder: GNN-assigned tracks. AMVF: its TTVA links (\texttt{Nonprompt\_Hard\_MaxWeight}, \texttt{RecoVertex\_assocTracks}). One rule for both. & \\
Clean, Merged, Fake & per vertex, the plurality truth vertex of its tracks. Fake if the plurality is unlabelled tracks; Clean if its share is $\ge\purityCut$; Merged otherwise & \texttt{classify\_\allowbreak assignments}, \texttt{gnn/evaluation/\allowbreak classification.py} \\
Split & among vertices with the same plurality truth vertex, all but the largest $\sum p_\mathrm{T}^2$ & same \\
Per event & Clean, Merged, Split, Fake vertices per event (bar plot at PU200, versus $\mu$) & \texttt{hs\_metrics.\allowbreak classify\_vertices} \\
Clean/truth & Clean vertices over truth vertices with $n_\mathrm{trk}\ge2$ & \texttt{chain\_t\_summary} \\
Cross-checks & AMVF with its fit tracks and fit weights; PV-Finder with GNN scores as weights & \texttt{weighted\_\allowbreak classification.py} \\
Bounds & finder cap: peaks matched one to one within \winMM{} mm. Truth-guided: each labelled track to the peak nearest its own truth vertex & \texttt{hs\_metrics}, \texttt{chain\_gap\_\allowbreak decomposition.py} \\
}
\par\vspace{.3em}{\footnotesize Errors: event bootstrap (200 replicas), paired. Every category count is cross-checked by independent code.}
\end{frame}

\begin{frame}{After the GNN: Hard Scatter}
\mtab{%
HS selection & score $=\sum p_\mathrm{T}^2$ over a vertex's tracks with compatibility $\chi^2\le\chiCut$ (Athena 24.0 \texttt{SumPt2Sorting}); selected vertex = argmax. Reconstructed and selected: the selected vertex is the HS match (most HS tracks) and lies within \hsWinMM{} mm of the true HS. AMVF reference: its stored primary. & \texttt{gnn/evaluation/\allowbreak pv\_selection.py}; \texttt{hs\_selection\_\allowbreak identical\_rule.py} \\
HS resolution & events where both algorithms have a vertex within \hsWinMM{} mm of the true HS; $z_\mathrm{reco}-z_\mathrm{HS}$ of the closest; width = iterative RMS within $\pm3$ RMS until stable (Athena \texttt{iterRMS\_convergence}). PV-Finder position: 3D adaptive fit with beam-spot constraint of the GNN-assigned tracks. & \texttt{itk\_resolution.\allowbreak iterative\_rms}; \texttt{hs\_position\_\allowbreak estimators.\allowbreak adaptive\_\allowbreak beamline\_fit} \\
Clean-HS fraction & events with a Clean vertex whose plurality truth is the HS & \texttt{hs\_metrics} \\
HS-track fraction & HS tracks on the HS-matched vertex over all HS tracks & same \\
Local density & truth vertices within $\pm2$ mm of the true HS, HS excluded, over 4 mm (all truth vertices, ITk definition) & \texttt{itk\_resolution.\allowbreak local\_pu\_density} \\
}
\end{frame}

\begin{frame}{After the GNN: Track Association}
\mtab{%
Correct track & assigned to a vertex whose plurality truth vertex is the track's own & \texttt{chain\_t\_metrics} \\
Track efficiency & correct tracks / truth-labelled tracks; unassigned tracks stay in the denominator & same \\
Track purity & correct tracks / assigned tracks with a truth label & same \\
Track $F_1$ & $2\epsilon p/(\epsilon+p)$ & same \\
Wrong rate & $1-$ purity & same \\
Unassigned fraction & $1-$ (assigned labelled tracks / labelled tracks); efficiency $=$ purity $\times(1-$ unassigned fraction$)$ & \texttt{hllhc\_final\_\allowbreak chain\_extras.py} \\
}
\par\vspace{.4em}{\footnotesize Tracks without a truth label enter neither efficiency nor purity, only the assigned fraction.}
\end{frame}

%======================================================================
\partdivider{4}{The 29 Questions}{Answers with page and line of the 11 November draft}

% Answers to the 29 questions of the 11 November draft.
% Source: docs/additional_studies/paper_questions_2026-11-11.md, updated to the
% frozen configuration of 12 November. Page and line refer to
% final_paper/current_paper/paper/PVFinderPaper_11_11_2026.pdf.

\begin{frame}{Q1. Do the 40 mm Intervals Overlap?}
\draftref{p.~7, l.~190}
\qtext{Do the intervals overlap slightly at their boundaries? If so, by how much, and how is a track in an overlap region assigned?}
\small
\begin{itemize}\setlength\itemsep{.45em}
\item \textbf{No overlap.} Interval $k$ covers $[-240+40k,\,-240+40(k+1))$ mm, $k=0,\dots,11$.
\item Each track goes to the one interval that contains its $z_0$; the lower edge is included.
\item Tracks outside $[-240,240)$ mm are dropped. The ntuple selection already requires $|z_0|\le200$ mm.
\item One sentence worth adding: the target histogram is built on the full 12,000-bin grid and then cut into twelve pieces, so a truth vertex near a boundary can put target content into the neighbouring interval, which holds none of its tracks.
\end{itemize}
\vfill\src{src/\allowbreak{}pv\_\allowbreak{}finder/\allowbreak{}data/\allowbreak{}root\_\allowbreak{}to\_\allowbreak{}h5.py:34-36, 67-104, 108-140}
\end{frame}

\begin{frame}{Q2. Is Every Element of Fig.~3 Right, and Does It Cover Run 2/3?}
\draftref{p.~8, Fig.~3 caption}
\small
\begin{columns}[T]
\begin{column}{.52\linewidth}
\textbf{Corrections for the final HL-LHC configuration}
\begin{itemize}\setlength\itemsep{.35em}
\item \textbf{Latent map: $1\times1000$}, not $4\times1000$. The frozen model \hlModel{} has one latent channel (v6 had four).
\item \textbf{Inputs}: seven features per track, $d_0$, $z_0$, $\sigma(d_0)$, $\sigma(z_0)$, $\mathrm{cov}(d_0,z_0)$, interval start, interval end. The helix parameters $(\phi,\theta,q/p)$ of l.~183 do not enter the density network; $\theta$, $\phi$, $p_\mathrm{T}$ enter only the GNN.
\item \textbf{GNN update order}: track-to-vertex and vertex-to-track updates run in parallel from the previous layer (one \texttt{HeteroConv}), not vertices first (l.~307--308).
\end{itemize}
\end{column}
\begin{column}{.45\linewidth}
\textbf{Run 2/3 differences}
\begin{itemize}\setlength\itemsep{.35em}
\item U-Net v1 instead of v2 (Q10).
\item \lpUnetCh{} channels instead of \hlUnetCh.
\item Fully connected track-vertex graphs instead of $k=\gnnK$.
\end{itemize}
\par\vspace{.6em}
\final{one figure covers both regimes if the caption lists these three differences. The latent map is $1\times1000$ in both.}
\end{column}
\end{columns}
\vfill\src{root\_\allowbreak{}to\_\allowbreak{}h5.py:37; \allowbreak{}src/\allowbreak{}gnn/\allowbreak{}models/\allowbreak{}ttva\_\allowbreak{}gat.py:67-121; \allowbreak{}working\_\allowbreak{}points.py}
\end{frame}

\begin{frame}{Q3. Target Widths Come from AMVF}
\draftref{p.~9, l.~205}
\qtext{The current target-width model is derived from AMVF vertex resolutions. Has anything been done to remove this dependence?}
\small
\begin{itemize}\setlength\itemsep{.4em}
\item \textbf{Not for this paper.} The final models are trained on the AMVF-derived widths.
\item Width model: $\sigma(n_\mathrm{trk})=A\,n_\mathrm{trk}^{-B}+C$, floor 2~$\mu$m; truth vertices with fewer than two tracks get one bin (0.04 mm).
\item HL-LHC preset \texttt{hllhc\_alleta}: $A=0.14318$ mm, $B=0.362$, $C=-0.01797$ mm, a background-corrected fit of AMVF-truth residuals (Q29).
\item Only the widths come from AMVF. The target positions are the truth vertex positions.
\item What removing it takes: a preset fitted on the track-only precision (inverse-variance $z_0$ of the truth tracks, already computed as a diagnostic), a rebuild of the \nPoolEvt-event pool and a retrain (\hlModel: 13.7 h on one A100).
\end{itemize}
\vfill\src{root\_\allowbreak{}to\_\allowbreak{}h5.py:82-88; \allowbreak{}src/\allowbreak{}pv\_\allowbreak{}finder/\allowbreak{}data/\allowbreak{}resolution\_\allowbreak{}presets.py}
\end{frame}

\begin{frame}{Q4. U-Net v2 Details}
\draftref{p.~9, l.~223}
\footnotesize
\begin{columns}[T]
\begin{column}{.55\linewidth}
\textbf{Yes, the Appendix E design}, with \hlUnetCh{} channels in the frozen model (\hlModel):
\begin{itemize}\setlength\itemsep{.25em}
\item \textbf{Encoder}: ConvBNReLU($1\to\hlUnetCh$, kernel 15) at 1000 bins; residual block (kernel 5) at 1000; max-pool 2, residual block at 500; max-pool 2, residual block at 250 (bottleneck, factor 4).
\item \textbf{Decoder}, twice: nearest-neighbour upsample $\times2$ + Conv(kernel 3) + BN + ReLU, additive skip, residual block (kernel 5). 250 $\to$ 500 $\to$ 1000.
\item \textbf{Head}: ConvBNReLU($80\to80$, k5), ConvBNReLU($80\to40$, k3), Conv($40\to1$, k1), softplus.
\item Residual block: Conv + BN + ReLU + Dropout, plus identity.
\end{itemize}
\end{column}
\begin{column}{.42\linewidth}
\begin{itemize}\setlength\itemsep{.25em}
\item \textbf{Normalisation}: BatchNorm after every convolution except the last $1\times1$. ReLU throughout.
\item \textbf{Dropout} 0.25 after the ReLU of every convolution block; none in the MLP.
\item \textbf{MLP}: $7\to128\times5\to1000$, LeakyReLU (0.01), softplus; padded tracks masked, output summed over tracks and scaled by $10^{-3}$.
\item \textbf{Parameters}: \hlParams{} = MLP \hlParamsMLP{} + U-Net \hlParamsUnet.
\end{itemize}
\par\vspace{.4em}
\final{Table B.1 (280 channels) and the 3.55 M parameters describe v6. Replace by \hlUnetCh{} channels and \hlParams.}
\end{column}
\end{columns}
\vfill\src{src/\allowbreak{}pv\_\allowbreak{}finder/\allowbreak{}models/\allowbreak{}unet\_\allowbreak{}v2.py; \allowbreak{}autoencoder\_\allowbreak{}models.py:214-352; \allowbreak{}configs/\allowbreak{}vertex\_\allowbreak{}finding/\allowbreak{}config\_\allowbreak{}hllhc\_\allowbreak{}alleta\_\allowbreak{}v6\_\allowbreak{}small\_\allowbreak{}s8.yml}
\end{frame}

\begin{frame}{Q5. Concatenation or Smoothing? One Batch?}
\draftref{p.~9, l.~232}
\small
\begin{itemize}\setlength\itemsep{.5em}
\item \textbf{Plain concatenation in $z$ order.} No smoothing, normalisation or boundary treatment. The intervals do not overlap (Q1).
\item The twelve intervals run as \textbf{one batch of 12}, padded to the widest interval.
\item In Athena the input tensor is $(12, 7, N_\mathrm{trk})$, also one batch.
\item Packed-track inference (CPU timing, M6) runs the MLP only on real tracks and scatters each into its interval. The output equals the padded one to $4\times10^{-6}$; it is the same model.
\end{itemize}
\vfill\src{run\_\allowbreak{}eval\_\allowbreak{}pvf\_\allowbreak{}run3.py:195-215; \allowbreak{}src/\allowbreak{}pv\_\allowbreak{}finder/\allowbreak{}models/\allowbreak{}packed\_\allowbreak{}inference.py}
\end{frame}

\begin{frame}{Q6. Phase 1: U-Net Absent or Frozen?}
\draftref{p.~10, l.~244}
\small
\begin{itemize}\setlength\itemsep{.5em}
\item \textbf{Absent.} Phase 1 runs only the MLP; its latent output is compared directly with the target histogram.
\item The U-Net parameters are also set to \texttt{requires\_grad = False} in Phase 1 and switched on for Phase 2. They have no effect, because the U-Net is not in the forward pass.
\item \final{the paper says ``absent''.}
\end{itemize}
\vfill\src{hllhc\_\allowbreak{}helpers.py:185-195 (phase1\_\allowbreak{}forward); \allowbreak{}train\_\allowbreak{}hllhc\_\allowbreak{}e2e.py:81-90, 128}
\end{frame}

\begin{frame}{Q7. Datasets, Checkpoint, Warm-up, Time, Hardware}
\draftref{p.~10, l.~250}
\footnotesize
\begin{columns}[T]
\begin{column}{.5\linewidth}
\begin{itemize}\setlength\itemsep{.3em}
\item \textbf{Sample}: the six PU200 files of Part 2, \nPoolEvt{} events.
\item \textbf{Split}: seed-42 permutation of 62,083,200 intervals: \nTrainSub{} / \nValSub{} / \nTestSub{} (96/3/1\%). Results use r16443 and r16638 only.
\item \textbf{Schedule of \hlModel}: Phase 1 one epoch, Phase 2 one epoch (v6: 2 + 2).
\item \textbf{Checkpoint}: Phase 2 step 450,000 of 465,624 (97\%). The run stopped on a host-memory error after the last step, before the end-of-epoch checkpoint; step 450,000 is the last saved state. The learning rate there is $1.3\times10^{-6}$, at the cosine floor.
\end{itemize}
\end{column}
\begin{column}{.47\linewidth}
\begin{itemize}\setlength\itemsep{.3em}
\item \textbf{Learning rate}: Phase 1 constant $10^{-3}$. Phase 2 per-step linear warm-up from $10^{-6}$ to $10^{-4}$ over 3,000 steps, then per-step cosine decay to $10^{-6}$.
\item \textbf{Other}: Adam $(0.9, 0.999)$, batch 128, gradient-norm clip 1.0, MSE loss.
\item \textbf{Time and hardware}: one NVIDIA A100-PCIE-40GB. Phase 1 9.1 h, Phase 2 4.6 h (resumed once after a cuDNN initialisation crash on the shared GPU). 13.7 h in total.
\end{itemize}
\end{column}
\end{columns}
\vfill\src{outputs/\allowbreak{}logs/\allowbreak{}hllhc\_\allowbreak{}alleta\_\allowbreak{}v6\_\allowbreak{}s8\_\allowbreak{}1p1ep\_\allowbreak{}20261109\_\allowbreak{}211919.log, \_\allowbreak{}20261110\_\allowbreak{}211930.log; \allowbreak{}build\_\allowbreak{}phase2\_\allowbreak{}scheduler, hllhc\_\allowbreak{}helpers.py:100-160}
\end{frame}

\begin{frame}{Q8. Peak Width, $\sigma_\mathrm{PV}$ and Height}
\draftref{p.~10, l.~262}
\small
\begin{itemize}\setlength\itemsep{.4em}
\item \textbf{Width}: square root of the amplitude-weighted variance of the bin index over the whole above-threshold region (up to a split point), times 0.04 mm.
\item \textbf{$\sigma_\mathrm{PV}$}: this width, used directly in the first edge feature $(z_0-z_\mathrm{PV})/\sqrt{\sigma(z_0)^2+\sigma_\mathrm{PV}^2}$. Median at PU200: \sigmaPVmed{} mm.
\item \textbf{Height}: the maximum bin value of the region, the raw predicted density without normalisation. Second vertex-node feature.
\item \textbf{Training graphs differ}: truth $z$, $\sigma_\mathrm{PV}$ from the \texttt{hllhc\_alleta} $n_\mathrm{trk}$ parameterisation, height from the target recipe. The text near l.~441 (Eq.~4) describes $\sigma_\mathrm{PV}$ as the $n_\mathrm{trk}$ parameterisation; at inference it is the peak width.
\end{itemize}
\vfill\src{efficiency\_\allowbreak{}res\_\allowbreak{}optimized\_\allowbreak{}atlas.py:16-190; \allowbreak{}graph\_\allowbreak{}construction.py:105-168; \allowbreak{}chain\_\allowbreak{}inference.py}
\end{frame}

\begin{frame}{Q9. Final Peak-Finding Configurations; Local Centroid at Low Pile-up?}
\draftref{p.~10, l.~269}
\small
\begin{columns}[T]
\begin{column}{.55\linewidth}
\renewcommand{\arraystretch}{1.2}
\begin{tabular}{@{}lll@{}}
\toprule
 & HL-LHC & Low pile-up \\
\midrule
Amplitude & \hlAmp & \hlAmp \\
Minimum width & \hlWidth{} bins & \hlWidth{} bins \\
Integral & \hlIntegral & \lpIntegral \\
Minimum height & \hlHeight & none \\
Position & centroid $\pm$\hlCentroid & centroid $\pm$\hlCentroid \\
$\svtx$ peaks & counted peaks & counted peaks \\
\bottomrule
\end{tabular}
\end{column}
\begin{column}{.42\linewidth}
\begin{itemize}\setlength\itemsep{.35em}
\item \textbf{Yes}, the local centroid is now used at low pile-up too. On Run 2 MC it changes the efficiency by \lpCentroidShift{} points.
\item The HL-LHC cuts are those of \hlModel. The draft's 0.30 / 0.03 are the v6 values.
\item The note's ``0.50 for the HL-LHC resolution fit'' described the old code. It is gone: the fit uses the counted peaks.
\end{itemize}
\end{column}
\end{columns}
\vfill\src{working\_\allowbreak{}points.py; \allowbreak{}lowpu\_\allowbreak{}final\_\allowbreak{}numbers\_\allowbreak{}2026-11-12.md Sec.~4}
\end{frame}

\begin{frame}{Q10. Run 2/3 Density Model}
\draftref{p.~11, l.~281}
\footnotesize
\begin{columns}[T]
\begin{column}{.52\linewidth}
\textbf{\lpModel{}, every low pile-up result.}
\begin{itemize}\setlength\itemsep{.25em}
\item \textbf{MLP}: $7\to100\times5\to1000$, LeakyReLU (0.01), softplus, 1 latent channel.
\item \textbf{U-Net v1}, \lpUnetCh{} channels, the original design: encoder ConvBNReLU kernels 25, 7, 5, 5 with max-pool, $1000\to500\to250\to125$ (bottleneck 125); decoder three ConvTranspose(2, 2) + ConvBNReLU(k5) with concatenated skips; output Conv($128\to64$, k5), Conv($64\to1$, k5), softplus; dropout 0.25 in every convolution block.
\item \textbf{Parameters}: \lpParams{} trainable.
\item \textbf{Inputs}: the same seven track features as HL-LHC.
\end{itemize}
\end{column}
\begin{column}{.45\linewidth}
\begin{itemize}\setlength\itemsep{.25em}
\item \textbf{Training}: Phase 1, 50 epochs MLP to histogram; Phase 2, 400 epochs end to end; Adam, constant $10^{-3}$ in both phases, batch 128, MSE. Checkpoint: last epoch (400).
\item \textbf{Data}: \lpMCTotal{} Run 2 MC events, 70 / 25 / 5\% (Part 2).
\item \textbf{No KDE in any step}: two-phase, the same procedure as HL-LHC.
\item \textbf{Target preset}: Qi Bin's file; upstream constants (preset \texttt{run3}) $A=0.23817$ mm, $B=0.49491$, $C=-0.00079$ mm.
\end{itemize}
\end{column}
\end{columns}
\vfill\src{configs/\allowbreak{}vertex\_\allowbreak{}finding/\allowbreak{}config\_\allowbreak{}mlp\_\allowbreak{}hist\_\allowbreak{}e2e.yml; \allowbreak{}autoencoder\_\allowbreak{}models.py:801-1000; \allowbreak{}train\_\allowbreak{}mlp\_\allowbreak{}hist\_\allowbreak{}then\_\allowbreak{}e2e.py}
\end{frame}

\begin{frame}{Q11. GNN Details}
\draftref{p.~12, l.~314}
\footnotesize
\begin{columns}[T]
\begin{column}{.5\linewidth}
\begin{itemize}\setlength\itemsep{.3em}
\item \textbf{Layers}: PyG \texttt{GATConv} (original GAT, not GATv2), 4 heads averaged (\texttt{concat=False}), edge features in the attention (\texttt{edge\_dim=3}), no self-loops. Two \texttt{HeteroConv} layers.
\item \textbf{Activations}: LeakyReLU (0.01) after the two linear encoders and the two hidden edge-MLP layers; ReLU after each residual sum; sigmoid score.
\item \textbf{Normalisation}: none. \textbf{Dropout}: none. The stored \texttt{dropout=0.25} is not used in \texttt{forward}, so the trained model has no dropout.
\end{itemize}
\end{column}
\begin{column}{.47\linewidth}
\begin{itemize}\setlength\itemsep{.3em}
\item \textbf{Parameters}: \gnnParams{} (v2 epoch \hlGnnEpoch).
\item \textbf{$d_0$}: as stored in the ntuple. In the HL-LHC simulation the beam spot is at $(0,0,0)$ with $\sigma_{xy}=12~\mu$m, so ``relative to the beam spot'' and ``relative to the origin'' are the same.
\item \textbf{Tracks}: exactly the tracks of the density network; the chain builds both inputs from the same arrays.
\item \textbf{Track features}: $d_0$, $z_0$, $\sigma(d_0)$, $\sigma(z_0)$, cov, $\theta$, $\phi$, $p_\mathrm{T}/50$ GeV.
\end{itemize}
\end{column}
\end{columns}
\vfill\src{src/\allowbreak{}gnn/\allowbreak{}models/\allowbreak{}ttva\_\allowbreak{}gat.py:54-137; \allowbreak{}pu200\_\allowbreak{}chain\_\allowbreak{}graphs.py:189-217; \allowbreak{}constants.py:10}
\end{frame}

\begin{frame}{Q12. Eq.~(1) and the Fake Fraction}
\draftref{p.~14, l.~361}
\small
\begin{itemize}\setlength\itemsep{.45em}
\item \textbf{Eq.~(1) equals the note's efficiency exactly}: truth with $\ge1$ reco vertex within the window = clean + merged truth after the closest-first assignment. Checked event by event on every model and on AMVF. No recalculation of the definition is needed; the numbers are recomputed for the final configuration.
\item \textbf{Fake fraction} (reco vertices without truth within \winMM{} mm, over reco vertices), PU200, \nEvtPU{} events:
\end{itemize}
\begin{center}
\renewcommand{\arraystretch}{1.2}
\begin{tabular}{@{}lccc@{}}
\toprule
 & Fake / reco [\%] & Fakes / event & Splits / event \\
\midrule
PV-Finder (\hlModel) & \fakeFracPVF & \fakeEvtPVF & \splitEvtPVF \\
AMVF & \fakeFracAMVF & \fakeEvtAMVF & \splitEvtAMVF \\
\bottomrule
\end{tabular}
\end{center}
\vfill\src{metric\_\allowbreak{}definitions.md Sec.~3.3a; \allowbreak{}model\_\allowbreak{}freeze\_\allowbreak{}candidates\_\allowbreak{}2026-11-12.md; \allowbreak{}HLDOC}
\end{frame}

\begin{frame}{Q13. One Window for Everything?}
\draftref{p.~14, l.~367}
\qtext{We think we agreed on a single window of 0.1 mm for all results. Is that correct?}
\small
\begin{itemize}\setlength\itemsep{.45em}
\item \textbf{One window, \winMM{} mm}, for every vertex result, both algorithms, both regimes. \textbf{\hsWinMM{} mm for hard-scatter results only}, as in the ITk performance paper.
\item Not 0.1 mm for all vertices: the 68\% single-vertex residual at PU200 is 55--56~$\mu$m for both algorithms, and 0.1 mm is less than twice that.
\item Every other window in the draft is replaced: 0.1694 mm, the 0.5 mm one-to-one match of the chain, the fitted 0.416 mm (Run 2 MC) and the fitted data $\svtx$ are all \winMM{} mm in the final numbers.
\item The 1 mm of the single-vertex residuals is the assignment radius used to measure the residual that sets the window. It is not a matching window and must stay wider than it.
\end{itemize}
\vfill\src{working\_\allowbreak{}points.py; \allowbreak{}matching\_\allowbreak{}window\_\allowbreak{}audit\_\allowbreak{}2026-11-09.md}
\end{frame}

\begin{frame}{Q14. Hard-Scatter Vertex Resolution}
\draftref{p.~15, l.~392}
\small
\begin{columns}[T]
\begin{column}{.55\linewidth}
\textbf{Include it.} Definition (ITk Fig.~28, Athena \texttt{iterRMS\_convergence}):
\begin{itemize}\setlength\itemsep{.3em}
\item events where both algorithms have a vertex within \hsWinMM{} mm of the true HS;
\item residual $z_\mathrm{reco}-z_\mathrm{HS}$ of the closest such vertex;
\item width = iterative RMS: restrict to the mean $\pm3$ RMS and recompute until the selection is stable.
\item PV-Finder position: 3D adaptive vertex fit with a beam-spot constraint of the tracks the GNN assigns to the HS candidate at $t=\hlT$.
\end{itemize}
\end{column}
\begin{column}{.42\linewidth}
\renewcommand{\arraystretch}{1.25}
\begin{tabular}{@{}lc@{}}
\toprule
PU200, \hsResN{} events & Resolution [$\mu$m] \\
\midrule
PV-Finder + GNN & $\hsResGNN\pm\hsResGNNErr$ \\
AMVF & $\hsResAMVF\pm\hsResAMVFErr$ \\
\bottomrule
\end{tabular}
\par\vspace{.6em}
In the note, say only how it is computed (M9).
\end{column}
\end{columns}
\vfill\src{itk\_\allowbreak{}hs\_\allowbreak{}resolution\_\allowbreak{}2026-11-09.md; \allowbreak{}CH hs\_\allowbreak{}resolution\_\allowbreak{}um}
\end{frame}

\begin{frame}{Q15. Local Pile-up Density}
\draftref{p.~15, l.~396}
\small
\begin{itemize}\setlength\itemsep{.45em}
\item \textbf{Include it.} Density as in the ITk paper: truth vertices within $\pm2$ mm of the true HS, HS excluded, over 4 mm. All truth vertices are counted, from Trey's side ntuple (same events, all truth vertices).
\item HS reconstruction at \hsWinMM{} mm is at least \hsRecoDensMin\% in every density bin for both algorithms.
\item HS selection (same Athena rule): PV-Finder + GNN leads AMVF by \hsSelDensLo{} to \hsSelDensHi{} points in every bin.
\item HS resolution: PV-Finder + GNN narrower by \hsResDensLo{} to \hsResDensHi{} $\mu$m in every bin, flat in density.
\item Also available (backup): vertex efficiency before and after the GNN versus local density.
\end{itemize}
\vfill\src{hllhc\_\allowbreak{}supplementary\_\allowbreak{}2026-11-12.md; \allowbreak{}SUP hs\_\allowbreak{}vs\_\allowbreak{}density.json}
\end{frame}

\begin{frame}{Q16. Track Efficiency and Purity Entering $F_1$}
\draftref{p.~16, l.~428}
\small
\begin{itemize}\setlength\itemsep{.45em}
\item \textbf{Correct track}: assigned to a vertex whose plurality truth vertex is the track's own truth vertex.
\item \textbf{Efficiency} $\epsilon$ = correct tracks / all tracks with a truth vertex. Unassigned tracks count in the denominator.
\item \textbf{Purity} $p$ = correct tracks / assigned tracks that have a truth vertex.
\item \textbf{Tracks without truth} are excluded from both; they enter only the assigned fraction.
\item $F_1=2\epsilon p/(\epsilon+p)$.
\end{itemize}
\begin{center}
\renewcommand{\arraystretch}{1.2}
\begin{tabular}{@{}lccc@{}}
\toprule
PU200 & Efficiency [\%] & Purity [\%] & $F_1$ [\%] \\
\midrule
PV-Finder + GNN & \trkEffGNN & \trkPurGNN & \trkFOneGNN \\
AMVF & \trkEffAMVF & \trkPurAMVF & \trkFOneAMVF \\
\bottomrule
\end{tabular}
\end{center}
\vfill\src{threshold\_\allowbreak{}scan.py:149-164, 240-252; \allowbreak{}chain\_\allowbreak{}scan.py:96-130}
\end{frame}

\begin{frame}{Q17. Wrong-Assignment and Unassigned-Track Metrics}
\draftref{p.~16, l.~434}
\small
\begin{itemize}\setlength\itemsep{.4em}
\item \textbf{Include them} (backup figures in Part 6). With $N_\mathrm{lab}$ labelled tracks, $N_\mathrm{at}$ assigned labelled tracks, $N_c$ correct:
\begin{itemize}
\item wrong-assignment rate $=(N_\mathrm{at}-N_c)/N_\mathrm{at}=1-p$: PV-Finder + GNN \wrongGNN\%, AMVF \wrongAMVF\% (PU200);
\item unassigned fraction $u=1-N_\mathrm{at}/N_\mathrm{lab}$;
\item lost unassigned: labelled tracks left unassigned whose truth vertex has no candidate within \winMM{} mm: \lostUnassigned;
\item unlabelled tracks correctly left unassigned / all unlabelled tracks: \unlabelledLeft.
\end{itemize}
\item \textbf{Relation to $F_1$}: $\epsilon=p\,(1-u)$. The wrong rate is the complement of the purity; the unassigned fraction is what separates efficiency from purity.
\end{itemize}
\vfill\src{hllhc\_\allowbreak{}final/\allowbreak{}chain\_\allowbreak{}extras/\allowbreak{}README.md; \allowbreak{}CX}
\end{frame}

\begin{frame}{Q18. Significance-Based Association, Eq.~(4)}
\draftref{p.~16, l.~445}
\small
\begin{columns}[T]
\begin{column}{.52\linewidth}
\begin{itemize}\setlength\itemsep{.35em}
\item \textbf{Eq.~(4) confirmed}: $s=|z_0-z_\mathrm{PV}|/\sqrt{\sigma(z_0)^2+\sigma_\mathrm{PV}^2}$, the same combination as the GNN edge feature.
\item Run on the same graph: each track's $k=\gnnK$ nearest peaks; the track goes to the smallest-$s$ peak if $s<s_\mathrm{max}$. $\sigma_\mathrm{PV}$ = peak width.
\item The best $s_\mathrm{max}$ is chosen on the reported events, which favours the baseline.
\end{itemize}
\end{column}
\begin{column}{.45\linewidth}
\renewcommand{\arraystretch}{1.2}
\begin{tabular}{@{}lc@{}}
\toprule
PU200 & Clean/truth [\%] \\
\midrule
PV-Finder + GNN & \cleanTruthGNN \\
Eq.~(4), best $s_\mathrm{max}$ & \eqFourBest \\
Eq.~(4), $s_\mathrm{max}=3$ & \eqFourThree \\
nearest $z_0$ & \zNear \\
$|z_0-z_\mathrm{PV}|<3\sigma(z_0)$ & \zThree \\
AMVF & \cleanTruthAMVF \\
\bottomrule
\end{tabular}
\end{column}
\end{columns}
\vfill\src{graph\_\allowbreak{}construction.py:128; \allowbreak{}chain\_\allowbreak{}extras README; \allowbreak{}CX}
\end{frame}

\begin{frame}{Q19. Why 25,000 Events and the 1,920 PU200 Events?}
\draftref{p.~17, l.~453}
\small
\begin{itemize}\setlength\itemsep{.45em}
\item \textbf{Compute budget}, not physics. The 25,000 are the first entries of r16443 (100,000 events, flat $\mu$ 0--200, about 500 per integer $\mu$).
\item The \nEvtPU{} are those with $185\le\mu\le215$ (actual 185--200, $\langle\mu\rangle=\meanMuPU$). The whole file has 7,680 events in that range.
\item \textbf{The full chain cannot use more of these files}: r16443 entries $\ge25{,}000$ and r16638 are GNN v2 training and validation data.
\item Before the GNN, more events are used where it helps: r16638 is the second held-out file (\nEvtPUb{} PU200 events) and the resolution versus $\mu$ uses \nEvtResMu{} events of both files.
\end{itemize}
\vfill\src{paper\_\allowbreak{}questions\_\allowbreak{}2026-11-11.md Q19; \allowbreak{}gnn\_\allowbreak{}full\_\allowbreak{}pileup\_\allowbreak{}training\_\allowbreak{}v2\_\allowbreak{}2026-11-11.md}
\end{frame}

\begin{frame}{Q20. Chance Matches, Window Scan, Efficiency vs $\mu$, $n_\mathrm{trk}$, Separation}
\draftref{p.~17, l.~470}
\small
\begin{itemize}\setlength\itemsep{.45em}
\item \textbf{Window scan}: show it as an auxiliary figure. PV-Finder minus AMVF is between \scanLo{} and \scanHi{} points for windows of 0.10--0.20 mm.
\item \textbf{3 mm-shift chance match} at \winMM{} mm: PV-Finder \chancePVF\%, AMVF \chanceAMVF\%. Keep it in the note.
\item \textbf{Efficiency vs $\mu$}: available, proposed for the paper.
\item \textbf{Efficiency vs truth $n_\mathrm{trk}$} (PUB-note style): available before and after the GNN.
\item \textbf{Efficiency vs separation}: available as efficiency versus $\Delta z$ to the nearest truth vertex, and versus local density (backup). The merged credit of Eq.~(1) credits both vertices of a shared peak, so it stays high at small separation; the one-to-one efficiency falls there.
\end{itemize}
\vfill\src{h1\_\allowbreak{}vertex\_\allowbreak{}finding/\allowbreak{}plots; \allowbreak{}chain\_\allowbreak{}extras/\allowbreak{}plots}
\end{frame}

\begin{frame}{Q21. $\svtx$ Threshold, AMVF Value and Fit}
\draftref{p.~17, l.~481}
\small
\begin{itemize}\setlength\itemsep{.45em}
\item \textbf{One threshold.} The fit uses the counted peaks (\hlModel{}: integral \hlIntegral, height \hlHeight). The draft's 0.229 mm came from a second peak finding at integral 0.50; v6 at its counted peaks gives 0.218 mm.
\item \textbf{Final}: PV-Finder $\svtx=\sigPVF\pm\sigPVFErr$ mm, AMVF $\sigAMVF\pm\sigAMVFErr$ mm, same \nEvtPU{} events.
\item \textbf{Fit}: all within-event pairs of vertices ($n_\mathrm{trk}\ge2$ for AMVF), $\Delta z$ histogram with \pairBins{} bins over $\pm6$ mm (0.04 mm), sigmoid $a/(1+e^{b(r-|x|)})+c$, $\svtx=|r|$. Same code for both. The error is the fit error.
\item The sigmoid assumes a flat plateau of pile-up pairs. It describes the data at high $\mu$ only, which is why $\svtx$ is quoted at PU200 (M10).
\end{itemize}
\vfill\src{run\_\allowbreak{}eval\_\allowbreak{}pvf\_\allowbreak{}run3.py fit\_\allowbreak{}dz; \allowbreak{}model\_\allowbreak{}freeze\_\allowbreak{}candidates\_\allowbreak{}2026-11-12.md Sec.~2}
\end{frame}

\begin{frame}{Q22. Merged and Split after the GNN at PU200; Empty Candidates}
\draftref{p.~18, l.~500}
\small
\begin{columns}[T]
\begin{column}{.5\linewidth}
\renewcommand{\arraystretch}{1.2}
\begin{tabular}{@{}lcc@{}}
\toprule
PU200, per reco [\%] & PVF + GNN & AMVF \\
\midrule
Clean & \cleanRecoGNN & \cleanRecoAMVF \\
Merged & \mergedRecoGNN & \mergedRecoAMVF \\
Split & \splitRecoGNN & \splitRecoAMVF \\
Fake & \fakeRecoGNN & \fakeRecoAMVF \\
\midrule
Clean / truth & \cleanTruthGNN & \cleanTruthAMVF \\
\bottomrule
\end{tabular}
\end{column}
\begin{column}{.47\linewidth}
\begin{itemize}\setlength\itemsep{.35em}
\item Table 4 now has the four categories of Table 5.
\item \textbf{Remove candidates without tracks from the output}, not only from the rates. A vertex without tracks cannot be written as an xAOD vertex with tracks. \emptyDropped{} candidates are dropped at PU200.
\item The paper figure is the number of vertices per event in each category (Part 6).
\end{itemize}
\end{column}
\end{columns}
\vfill\src{CH table4; \allowbreak{}SUP classes}
\end{frame}

\begin{frame}{Q23. GNN Thresholds at PU200; One Working Point?}
\draftref{p.~18, l.~514}
\small
\begin{itemize}\setlength\itemsep{.45em}
\item \textbf{One working point, $t=\hlT$}, for vertex classification, HS selection, HS resolution and track $F_1$, in both regimes. The 0.98 and 0.95 values of the draft are dropped.
\item Chosen on the \gnnVal{} GNN validation events (disjoint from the test set) as the most $\mu$-robust value; details in M4.
\item At $t=\hlT$, HS reconstructed and selected at PU200 is \hsSelGNN\% (AMVF stored primary \hsSelAMVF\%).
\end{itemize}
\vfill\src{gnn\_\allowbreak{}v2\_\allowbreak{}freeze\_\allowbreak{}2026-11-12.md Sec.~5; \allowbreak{}working\_\allowbreak{}points.py}
\end{frame}

\begin{frame}{Q24. AMVF Run 2 Resolution and Efficiency on the Same Events}
\draftref{p.~20, l.~528}
\small
\textbf{Both are now on the same \lpMCTest{} Run 2 MC events with the same code.}
\begin{center}\footnotesize
\renewcommand{\arraystretch}{1.25}
\begin{tabular}{@{}lcccc@{}}
\toprule
 & $\svtx$ [mm] & Eff., \winMM{} mm [\%] & Fake/reco [\%] & 68\% res. [$\mu$m] \\
\midrule
PV-Finder (\lpModel) & $\lpSigPVF\pm\lpSigPVFErr$ & $\lpEffPVF\pm\lpEffPVFErr$ & \lpFakeFracPVF & \lpPsePVF \\
AMVF & $\lpSigAMVF\pm\lpSigAMVFErr$ & \lpEffAMVF & \lpFakeFracAMVF & \lpPseAMVF \\
\bottomrule
\end{tabular}
\end{center}
\begin{itemize}\setlength\itemsep{.3em}
\item The draft's AMVF 0.793 mm was on 2,500 other events; the PV-Finder 0.416 mm used 60 bins and the ep300 model.
\end{itemize}
\vfill\src{LPDOC R1; \allowbreak{}VF mc}
\end{frame}

\begin{frame}{Q25. Low Pile-up GNN at Integral 0.2, Vertex Finding at 0.50}
\draftref{p.~20, l.~535}
\small
\begin{itemize}\setlength\itemsep{.45em}
\item \textbf{Final: integral \lpIntegral{} for both}, on \lpModel{} peaks with the local centroid.
\item The difference was more than the threshold: the draft's GNN numbers (77.43\% Clean/truth, $t=0.99$) came from PVF-e400 at integral 0.2 and the full-region mean; the vertex finding from the ep300 model at 0.50.
\item \textbf{Recalculated}: graphs rebuilt from the frozen model at \lpIntegral, GNN at $t=\lpT$:
\begin{itemize}
\item Clean/truth \lpCleanTruth\% (AMVF \lpCleanTruthAMVF\%);
\item Clean / Merged / Split / Fake per reco: \lpClean{} / \lpMerged{} / \lpSplit{} / \lpFake\%.
\end{itemize}
\item Integral \lpIntegral{} finds fewer candidates than 0.2, mostly small vertices; that and the lower $t$ explain the change from 77.43\%.
\end{itemize}
\vfill\src{LPDOC R2, Sec.~4; \allowbreak{}GM}
\end{frame}

\begin{frame}{Q26. GNN Thresholds in Run 2 Simulation}
\draftref{p.~21, l.~548}
\small
\begin{itemize}\setlength\itemsep{.45em}
\item \textbf{One $t=\lpT$} for vertex classification, track $F_1$ and HS identification. The 0.99 / 0.95 / 0.5 of the draft are dropped.
\item Rule: the HL-LHC value unless clearly worse at low pile-up. No $t$ beats \lpT{} in Clean/truth, HS identification and $F_1$ together by more than two paired standard errors.
\item At $t=\lpT$: Clean/truth \lpCleanTruth\%, $F_1$ \lpFOne\% (AMVF \lpFOneAMVF\%), HS identification \lpHSID\% (AMVF \lpHSIDAMVF\%).
\item $t=0.99$ gives Clean/truth \lpCleanTruthNN\%, with a lower $F_1$; its gain comes from $\mu>40$ only, and at $\mu<20$ $t=\lpT$ is better.
\item The draft's 98.98\% HS identification was obtained at $t=0.99$, not at the 0.5 the note lists.
\end{itemize}
\vfill\src{LPDOC Sec.~3; \allowbreak{}L/\allowbreak{}gnn\_\allowbreak{}mc/\allowbreak{}t\_\allowbreak{}options.md}
\end{frame}

\begin{frame}{Q27. Run 3 Data Association: $t$, Peak Threshold, Events}
\draftref{p.~22, l.~574}
\small
\begin{columns}[T]
\begin{column}{.5\linewidth}
\begin{itemize}\setlength\itemsep{.4em}
\item \textbf{Final}: $t=\lpT$, \lpModel{} at integral \lpIntegral{} with the local centroid.
\item \textbf{Events}: the same 2,500 Run 3 events and the same peaks as the vertex finding. The draft's 2,000 events were a random sample of \texttt{file\_3.root}; only 8 were in the vertex-finding set.
\item The draft's model was the two-stage KDE model at $t=0.98$; both are replaced.
\end{itemize}
\end{column}
\begin{column}{.47\linewidth}
\renewcommand{\arraystretch}{1.2}
\begin{tabular}{@{}lc@{}}
\toprule
Run 3 data, $t=\lpT$ & [\%] \\
\midrule
Tracks assigned (GNN) & \lpRthreeAssigned \\
Agreement with AMVF, 1 mm & \lpRthreeAgree \\
Agreement with AMVF, \winMM{} mm & \lpRthreeAgreeW \\
Leading-vertex agreement & \lpRthreeLead \\
\bottomrule
\end{tabular}
\par\vspace{.4em}
{\footnotesize Agreement is over tracks assigned by both. These are agreement numbers, not efficiencies.}
\end{column}
\end{columns}
\vfill\src{LPDOC R4; \allowbreak{}R3A}
\end{frame}

\begin{frame}{Q28. ``integral=0.2'' in the Fig.~10 Panel Titles}
\draftref{p.~23, Fig.~10 caption}
\small
\begin{itemize}\setlength\itemsep{.45em}
\item \textbf{0.5 was used}; the titles were stale. Re-running the old checkpoint on 25 Run 2 MC events reproduces the stored peak lists in 23 of 25 events at 0.5 and in 7 of 25 at 0.2; the other 2 differ by CPU/GPU rounding.
\item The label came from the April 2026 plotting code, which carried the then-default value.
\item \textbf{Final}: all panels are regenerated with \lpModel{} at the final configuration. Every panel states ``integral threshold 0.50''; simulation panels say ``Simulation Internal'', data panels ``Internal''; the Run 3 panels name the simulation Run 2 MC.
\end{itemize}
\vfill\src{LPDOC Sec.~6; \allowbreak{}L/\allowbreak{}summary/\allowbreak{}figures/\allowbreak{}fig10*}
\end{frame}

\begin{frame}{Q29. AMVF Dependence; Sample of \texttt{hllhc\_alleta}; Run 2/3 Parameters}
\draftref{p.~28, l.~652}
\small
\begin{itemize}\setlength\itemsep{.45em}
\item \textbf{AMVF dependence}: Q3.
\item \textbf{\texttt{hllhc\_alleta}} was fitted on the \textbf{full-$|\eta|$ sample} ($|\eta|<4$): 601229 r16633, 99,800 events, 2.31 M truth vertices with $n_\mathrm{trk}\ge2$, background-corrected fit of AMVF-truth residuals. Not the $|\eta|<2.5$ sample.
\item Parameters: $A=0.14318$ mm, $B=0.362$, $C=-0.01797$ mm.
\item \textbf{Run 2/3} (preset \texttt{run3}, Qi Bin's file): $A=0.23817443$ mm, $B=0.49491396$, $C=-0.000787436$ mm.
\end{itemize}
\vfill\src{src/\allowbreak{}pv\_\allowbreak{}finder/\allowbreak{}data/\allowbreak{}resolution\_\allowbreak{}presets.py; \allowbreak{}outputs/\allowbreak{}07\_\allowbreak{}31\_\allowbreak{}2026\_\allowbreak{}output/\allowbreak{}amvf\_\allowbreak{}resolution\_\allowbreak{}alleta/\allowbreak{}}
\end{frame}


%======================================================================
\partdivider{5}{Mattia's Questions}{Decisions behind the configuration}

% Answers to Mattia's own questions (meeting of 11 November 2026).

\begin{frame}{M1. Why a \winMM{} mm Window}
\small
\begin{columns}[T]
\begin{column}{.5\linewidth}\footnotesize
\begin{itemize}\setlength\itemsep{.3em}
\item The window asks whether a vertex was put where the truth vertex is. It is set by how precisely one vertex is placed, the same for both algorithms.
\item PU200, \nEvtPU{} events: reco and truth vertices assigned one to one within \assignMM{} mm (most matches, then least total distance). 68\% residual: PV-Finder \pseVsix{} $\mu$m, AMVF \pseAMVF{} $\mu$m.
\item Window $=3\times$ the larger $=3\times\pseAMVF~\mu\mathrm{m}=0.1694$ mm, rounded to \winMM{} mm. Not a Gaussian $3\sigma$: the residuals have tails.
\item A greedy assignment would bias the width about \greedyBias\% low.
\end{itemize}
\end{column}
\begin{column}{.47\linewidth}\centering
\plot{window_decision}{.42}
\fcap{Single-vertex residuals at PU200 (v6, 9 November). The \hlModel{} version is in Part 6.}
\par\vspace{.3em}
{\footnotesize Not $\svtx$: it measures how close two vertices can be and still be separated, and differs by 40\% between the algorithms. At low pile-up the window is \lpWinRatio{} times the PV-Finder residual.}
\end{column}
\end{columns}
\end{frame}

\begin{frame}{M2. Hard Scatter at \hsWinMM{} mm, for ITk Consistency}
\small
\begin{itemize}\setlength\itemsep{.5em}
\item The ITk performance paper (arXiv:2412.15090, Sec.~5.5, Fig.~26b) reconstructs the hard-scatter vertex with a \hsWinMM{} mm criterion and quotes only the hard scatter. The HS has tens of tracks and a resolution near 10~$\mu$m, so \hsWinMM{} mm is wide for it.
\item The all-vertex population is dominated by low-multiplicity pile-up vertices with a 68\% residual of 55--56~$\mu$m; for them \hsWinMM{} mm would mix position precision into the efficiency. They keep \winMM{} mm.
\item Every HS number uses \hsWinMM{} mm: reconstruction, selection, resolution.
\item HS reconstruction at PU200: PV-Finder \hsRecoPVF\%, AMVF \hsRecoAMVF\%.
\end{itemize}
\takeaway{Sentence for the paper: ``For consistency with ITk performance studies, which quote only the hard-scatter vertex, we also give the hard-scatter results with a 0.1~mm window.''}
\end{frame}

\begin{frame}{M3. Why \hlModel{} Was Frozen}
\footnotesize
\begin{columns}[T]
\begin{column}{.55\linewidth}
\scriptsize\renewcommand{\arraystretch}{1.05}
\begin{tabular}{@{}p{.5\linewidth}p{.44\linewidth}@{}}
\toprule
\hlModel{} minus v6, same events & \\
\midrule
Efficiency at the v6 fake rate (r16443, r16638) & \sEightDEffA{}, \sEightDEffB{} pt \\
$\svtx$ & \sEightDSig{} $\mu$m \\
68\% residual & \sEightDpse{} $\mu$m \\
Clean/truth after the GNN ($t=\hlT$) & \sEightDClean{} pt \\
HS selection after the GNN & \sEightDHSsel{} pt \\
HS resolution after the GNN & \sEightDHSres{} $\mu$m \\
\midrule
CPU speed-up, packed, vs v6 as deployed & \spdSEightNov$\times$ (9 Nov), \spdPack$\times$ (12 Nov) \\
Parameters & \hlParams{} vs 3,547,993 \\
\bottomrule
\end{tabular}
\end{column}
\begin{column}{.42\linewidth}\scriptsize
\begin{itemize}\setlength\itemsep{.25em}
\item Bar: efficiency loss below 0.2 points at the v6 fake rate, and no worse resolution. Only \hlModel{} meets it.
\item s5 (112 channels) separates close vertices worse: $\svtx$ \sFiveDSig{} $\mu$m at 6--7 standard deviations, with either schedule.
\item Cuts: v6's 0.30 / 0.03 scaled until the fakes per event match v6; the two-file mean is \hlIntegral{} / \hlHeight.
\item After the GNN the trends with $\mu$ are the same for both models.
\end{itemize}
\end{column}
\end{columns}

\vfill\src{model\_\allowbreak{}freeze\_\allowbreak{}candidates\_\allowbreak{}2026-11-12.md; \allowbreak{}gnn\_\allowbreak{}v2\_\allowbreak{}freeze\_\allowbreak{}2026-11-12.md Sec.~7; \allowbreak{}cpu\_\allowbreak{}small\_\allowbreak{}model\_\allowbreak{}2026-11-09.md}
\end{frame}

\begin{frame}{M3. \hlModel{} and v6 after the GNN}
\centering
\plot{v6_vs_s8_gnn}{.72}
\fcap{Chain metrics versus $\mu$ for both density models with GNN v2, test events (12 November freeze). The trends are the same; the differences are at most a few tenths of a point.}
\end{frame}

\begin{frame}{M4. The GNN Score Cut $t=\hlT$}
\footnotesize
\begin{columns}[T]
\begin{column}{.6\linewidth}
\scriptsize\setlength{\tabcolsep}{3pt}\renewcommand{\arraystretch}{1.1}
\begin{tabular}{@{}llccccc@{}}
\toprule
Rule (on validation) & $t$ & \multicolumn{3}{c}{Clean/truth [\%]} & $F_1$ & HS res \\
 & & incl. & PU200 & $\mu<20$ & [\%] & [$\mu$m] \\
\midrule
A: max Clean/truth & 0.98 & \optAclean & \optAcleanPU & \optAcleanLow & \optAFone & \optAHSres \\
\textbf{D: most $\mu$-robust} & \textbf{0.97} & \optDclean & \optDcleanPU & \optDcleanLow & \optDFone & \optDHSres \\
B: max HS selection & 0.94 & \optBclean & & & & \\
C: max track $F_1$ & 0.92 & \optCclean & & & & \\
E: no cut & 0 & \optEclean & & & & \\
AMVF & & \cleanTruthAMVFincl & \cleanTruthAMVF & \optAMVFcleanLow & \optAMVFFone & \optAMVFHSres \\
\bottomrule
\end{tabular}
\par\vspace{.3em}
{\tiny Test events, v6 peaks (12 November freeze). D: largest worst-case ratio, over five $\mu$ bins, of Clean/truth to its best value in the bin.}
\end{column}
\begin{column}{.37\linewidth}\scriptsize
\begin{itemize}\setlength\itemsep{.2em}
\item Checkpoint: epoch \hlGnnEpoch, the lowest validation-monitor loss (\gnnValLoss) of 540 shard epochs; plateaued at the learning-rate floor.
\item $t=\hlT$ is within \optDvsAclean{} points of the maximum Clean/truth, gains \optDvsAlow{} points at $\mu<20$, \optDvsAFone{} points of $F_1$ and \optDvsAres~$\mu$m of HS resolution.
\item $t=0.99$ is worse than 0.98 in Clean/truth, $F_1$, HS selection and HS resolution.
\item Chosen by Mattia, 12 November; also used at low pile-up.
\end{itemize}
\end{column}
\end{columns}

\vfill\src{gnn\_\allowbreak{}v2\_\allowbreak{}freeze\_\allowbreak{}2026-11-12.md Sec.~1, 5; \allowbreak{}G/\allowbreak{}options.md}
\end{frame}

\begin{frame}{M4. The $t$ Scan in $\mu$ Bins}
\centering
\plot{t_scan_val_s8}{.72}
\fcap{All chain metrics versus $t$ per $\mu$ bin, GNN validation events, \hlModel{} peaks.}
\end{frame}

\begin{frame}{M5. Behaviour of the GNN at Low $\mu$}
\footnotesize
\begin{columns}[T]
\begin{column}{.55\linewidth}\scriptsize
\textbf{Seen} at $\mu<20$: Clean/truth peaks at $t=\lmCleanPeakT$ and falls at higher $t$; above $\mu=50$ it rises up to 0.98.
\par\vspace{.3em}
\textbf{Why}, measured:
\begin{itemize}\setlength\itemsep{.15em}
\item At $\mu<10$, \lmTopOneLow\% of top-1 edges are correct (\lmTopOneHigh\% at $\mu$ 150--200). A higher $t$ then removes mostly good tracks.
\item The network is less confident on small vertices at low $\mu$: \lmSmallBelowLow\% of correct tracks of 5--10-track vertices score below 0.99 at $\mu<10$ (\lmSmallBelowMid\% at $\mu$ 20--50, \lmSmallBelowHigh\% at 150--200), on truth-candidate graphs too.
\item $\mu<20$ events hold about \lmTrackShare\% of the tracks of a flat-$\mu$ sample, so they weigh little in the loss.
\item v1 was worse (\lmVoneBelow\% vs \lmVtwoBelow\% of correct tracks below 0.99 at $\mu<10$): the behaviour Mattia saw.
\item Splits at low $\mu$ are HS satellites: \lmSplitHS\% have the HS as dominant truth, \lmSplitOne\% hold one track.
\end{itemize}
\end{column}
\begin{column}{.42\linewidth}
\begin{itemize}\setlength\itemsep{.25em}
\item Fake/reco stays at or below \lmFakeMax\%; there is no failure mode.
\item The satellite splits are weak density peaks next to the large HS peak, a larger share of the list when there are few vertices. They are not GNN errors.
\item $t=\hlT$ is the compromise: within \optDvsAclean{} points of the best Clean/truth at high $\mu$, \optDvsAlow{} points better than 0.98 at $\mu<20$.
\end{itemize}
\end{column}
\end{columns}
\par\vspace{.2em}
{\scriptsize At the working point, $\mu$ 0--10: track $F_1$ \lmFOneGNN\% (PV-Finder + GNN) vs \lmFOneAMVF\% (AMVF). No track-$F_1$ claim at low $\mu$ in the paper.}
\end{frame}

\begin{frame}{M5. Low-$\mu$ Event Displays}
\centering
\plot{lowmu_displays}{.72}
\fcap{Validation events at $\mu\approx17$: isolated 7- and 11-track vertices whose tracks all score 0.9--0.99, so the vertex vanishes at $t=0.99$.}
\end{frame}

\begin{frame}{M6. CPU Timing against AMVF}
\footnotesize
\begin{columns}[T]
\begin{column}{.5\linewidth}
\textbf{Projected Athena CPU time} (Xeon 4216, one thread), from speed-ups measured on sneezy and the note's Athena table:
\[T_\mathrm{chain}=\frac{\athCPUinf~\mathrm{ms}}{S}+\athOther~\mathrm{ms}+\athGNN~\mathrm{ms}\]
\renewcommand{\arraystretch}{1.2}
\begin{tabular}{@{}lcc@{}}
\toprule
Chain & $S$ & $T_\mathrm{chain}$ [ms] \\
\midrule
v6 padded (Athena today) & 1 & \athCPUms \\
\hlModel{} packed + GNN v2 & \spdPack & \projChain \\
\hlModel{} padded + GNN v2 & \spdPad & \projChainPad \\
AMVF (note figure, $\mu=200$) & & \athAMVF{} (\athAMVFlo--\athAMVFhi) \\
\bottomrule
\end{tabular}
\end{column}
\begin{column}{.47\linewidth}
\begin{itemize}\setlength\itemsep{.35em}
\item With packed inference the chain is \projFactor$\times$ faster than AMVF on CPU. The worst case tried, \projWorst{} ms, is still below \athAMVFlo{} ms.
\item Packed inference needs a new Athena tensor builder: a track list plus an interval index.
\item With today's padded tensor, \hlModel{} is comparable to AMVF, not faster.
\item The GNN is now the largest step (\athGNN{} ms in Athena).
\item On GPU the full chain takes \athGPUms{} ms (note).
\item AMVF cannot run on sneezy: its time is read off the note's figure. The chain time is a projection, not an Athena measurement.
\end{itemize}
\end{column}
\end{columns}
\vfill\src{H/\allowbreak{}cpu\_\allowbreak{}timing/\allowbreak{}README.md; \allowbreak{}note sections/\allowbreak{}04b\_\allowbreak{}timing.tex, tab:athena-timing-breakdown}
\end{frame}

\begin{frame}{M7. Binning: Is the Network or the 0.04 mm Grid the Limit?}
\footnotesize
\qtext{Professor's question: does the smaller Run 2 $\to$ HL-LHC gain of the AMVF/PV-Finder ratio mean that PV-Finder does not use the better HL-LHC tracks, so that finer bins are needed?}
\begin{columns}[T]
\begin{column}{.52\linewidth}
Test: the training target is what a perfect network would output. It goes through the same peak finder and the same metrics.
\par\vspace{.3em}
\renewcommand{\arraystretch}{1.15}
\begin{tabular}{@{}lcc@{}}
\toprule
$\svtx$ [mm] & Run 2 MC & HL-LHC PU200 \\
\midrule
Target, 0.04 mm grid & \tgtSigRunTwo & \tgtSigHL{} (\tgtSigHLwone$^\ast$) \\
Target, 0.02 mm grid$^\ast$ & & \tgtSigHLfine \\
PV-Finder & \pvfSigRunTwo & \pvfSigHLvsix{} (v6) \\
AMVF & \amvfSigRunTwo & \sigAMVF \\
\bottomrule
\end{tabular}
\par{\scriptsize $^\ast$1-bin minimum width, so that every HS target spike is found.}
\end{column}
\begin{column}{.45\linewidth}
\begin{itemize}\setlength\itemsep{.3em}
\item PV-Finder is 2.0 times its target in Run 2 and 1.5--1.7 times at HL-LHC: closer to its ideal at HL-LHC.
\item A 0.02 mm grid improves the ideal by 10\%, far from the gap to PV-Finder.
\item The 0.04 mm grid carries the HS position to \tgtHSum~$\mu$m; PV-Finder reaches \pvfHSum~$\mu$m.
\item Track-based ideal precision of pile-up vertices: \idealPURunTwo{} $\to$ \idealPUHL~$\mu$m (1.7$\times$). At the same peak setting PV-Finder $\svtx$ also improves 1.7$\times$ (\pvfSigRunTwo{} $\to$ \pvfSigHLrunTwoCfg{} mm).
\end{itemize}
\end{column}
\end{columns}
\takeaway{The headroom is in the network, not in the grid. Finer bins are not needed.}
\end{frame}

\begin{frame}{M8. Hard-Scatter Selection with One Identical Rule}
\footnotesize
\begin{columns}[T]
\begin{column}{.5\linewidth}
\textbf{The Athena 24.0 rule} (read from source):
\begin{itemize}\setlength\itemsep{.25em}
\item vertices sorted by $\sum p_\mathrm{T}^2$ of their linked tracks (\texttt{SumPt2Sorting}); the first is \texttt{PriVtx};
\item a track is linked if its compatibility with the vertex is $\chi^2\le\chiCut$; no weights, no $p_\mathrm{T}$ cap;
\item HS match: the vertex with the most HS tracks (\texttt{bestHardScatterMatch}).
\end{itemize}
\textbf{Validation on AMVF} (first 3,000 entries): on its stored lists our implementation agrees with the stored primary in \ruleAgreeAll\% of events (\ruleAgreePU\% at PU200); plain $\sum p_\mathrm{T}^2$ gives \plainAgreeAll\%.
\end{column}
\begin{column}{.47\linewidth}
\renewcommand{\arraystretch}{1.2}
\begin{tabular}{@{}lc@{}}
\toprule
PU200, reco and selected (\hsWinMM{} mm) & [\%] \\
\midrule
PV-Finder + GNN, Athena rule & $\hsSelGNN\pm\hsSelGNNErr$ \\
AMVF, stored primary & \hsSelAMVF \\
AMVF, same rule, our code & \hsSelAMVFrule \\
AMVF, plain $\sum p_\mathrm{T}^2$ (old note) & \hsSelAMVFplain \\
\midrule
PV-Finder + GNN $-$ AMVF stored & $\hsSelDiff\pm\hsSelDiffErr$ pt \\
\bottomrule
\end{tabular}
\par\vspace{.4em}
\begin{itemize}\setlength\itemsep{.25em}
\item Headline: GNN with the Athena rule against AMVF's own stored primary. Our code can only lose AMVF information, so this favours AMVF.
\item It replaces the note's +3.4 points, which compared different rules.
\end{itemize}
\end{column}
\end{columns}
\vfill\src{hs\_\allowbreak{}selection\_\allowbreak{}identical\_\allowbreak{}rule\_\allowbreak{}2026-11-11.md; \allowbreak{}src/\allowbreak{}gnn/\allowbreak{}evaluation/\allowbreak{}pv\_\allowbreak{}selection.py}
\end{frame}

\begin{frame}{M9. Hard-Scatter Resolution: Why PV-Finder + GNN Reaches AMVF}
\footnotesize
\begin{columns}[T]
\begin{column}{.5\linewidth}
\renewcommand{\arraystretch}{1.15}
\begin{tabular}{@{}lc@{}}
\toprule
HS position (GNN v4 study) & [$\mu$m] \\
\midrule
Histogram peak, local centroid & \hsResPeak \\
1D adaptive fit of $z_0$, GNN tracks & \hsResOneD \\
3D adaptive fit, all tracks within 1 mm & \hsResNoGNN \\
3D adaptive fit, AMVF's own HS tracks & \hsResAMVFlist \\
3D adaptive fit, true HS tracks & \hsResTruthTrk \\
AMVF stored vertex & \hsResAMVFstudy \\
\midrule
\textbf{Final}, 3D fit, GNN tracks, $t=\hlT$ & $\hsResGNN\pm\hsResGNNErr$ \\
\textbf{Final}, AMVF & $\hsResAMVF\pm\hsResAMVFErr$ \\
\bottomrule
\end{tabular}
\end{column}
\begin{column}{.47\linewidth}
\begin{itemize}\setlength\itemsep{.3em}
\item The missing piece was the transverse vertex position. ITk tracks reach $|\eta|=4$ (median $|\cot\theta|=3.2$), so a 10--15~$\mu$m offset from the beam axis moves $z_0$ by tens to hundreds of $\mu$m.
\item The fit: $(x,y,z)$ with a beam-spot constraint, $(d_0,z_0)$ covariance, adaptive weights with annealing, seeded at the histogram peak.
\item GNN tracks do as well as the true HS tracks; on AMVF's own tracks the same fit beats AMVF's stored position. The fit model, not the association, made the difference.
\item In the note: say only how it is computed.
\end{itemize}
\end{column}
\end{columns}
\vfill\src{itk\_\allowbreak{}hs\_\allowbreak{}resolution\_\allowbreak{}2026-11-09.md (2026-11-11 sections); \allowbreak{}hs\_\allowbreak{}position\_\allowbreak{}estimators.adaptive\_\allowbreak{}beamline\_\allowbreak{}fit}
\end{frame}

\begin{frame}{M10. Resolution versus $\mu$: the 68\% Residual}
\footnotesize
\begin{columns}[T]
\begin{column}{.5\linewidth}
\textbf{Choice}: the resolution versus $\mu$ in the paper is the 68\% residual $|z_\mathrm{reco}-z_\mathrm{true}|$. $\svtx$ is quoted at PU200 only, or per bin for $\mu\ge100$.
\par\vspace{.4em}
\textbf{Why}:
\begin{itemize}\setlength\itemsep{.15em}
\item Both algorithms make more close vertex pairs than truth has: close-pair rate over flat expectation \rvmExcessPVF{} (PV-Finder), \rvmExcessAMVF{} (AMVF), \rvmExcessTruth{} (truth) at $\mu<20$. The hard-scatter neighbourhood is the largest source.
\item This excess is about constant per event while random pairs grow as $\mu^2$, so it dominates at low $\mu$. Below $\mu=100$ the sigmoid fit measures a notch inside it, not the separation ($\chi^2/\mathrm{ndf}$ \rvmChiLo--\rvmChiHi).
\item $\svtx$ is not flat in $\mu$ for either algorithm, so a single all-$\mu$ value is wrong.
\item The 68\% residual is defined at every $\mu$; the two algorithms agree within 2~$\mu$m in every bin.
\end{itemize}
\end{column}
\begin{column}{.47\linewidth}\centering
\plot{resolution_vs_mu}{.53}
\fcap{68\% residual (top) and $\svtx$ (bottom; open markers below $\mu=100$), \hlModel.}
\end{column}
\end{columns}
\vfill\src{resolution\_\allowbreak{}vs\_\allowbreak{}mu\_\allowbreak{}2026-11-11.md; \allowbreak{}H/\allowbreak{}resolution\_\allowbreak{}vs\_\allowbreak{}mu}
\end{frame}


%======================================================================
\partdivider{6}{Results}{Every paper figure and number; PV-Finder and AMVF on the same events}

% Results: every paper figure and number at the frozen configuration.
% Figures through images.txt; placeholders until the recompute agents finish.

%---------------------------------------------------------------------- HL-LHC before the GNN
\begin{frame}{HL-LHC before the GNN: Table 3 at PU200}
\mainbadge\hfill{\footnotesize r16443, \nEvtPU{} events, $\langle\mu\rangle=\meanMuPU$, \nTruthPU{} truth vertices ($n_\mathrm{trk}\ge2$), window \winMM{} mm}
\par\vspace{.3em}\small
\begin{center}
\renewcommand{\arraystretch}{1.12}
\begin{tabular}{@{}lccc@{}}
\toprule
 & PV-Finder (\hlModel) & AMVF & PVF $-$ AMVF \\
\midrule
Efficiency, Eq.~(1) [\%] & $\effPVF\pm\effPVFErr$ & \effAMVF & $\effDiff\pm\effDiffErr$ pt \\
One-to-one efficiency [\%] & \effOneOnePVF & \effOneOneAMVF & \\
Fake / reco [\%] & \fakeFracPVF & \fakeFracAMVF & \\
Fakes / event & \fakeEvtPVF & \fakeEvtAMVF & \\
Splits / event & \splitEvtPVF & \splitEvtAMVF & \\
Reco vertices / event & \recoEvtPVF & \recoEvtAMVF & (truth \truthEvt) \\
$\svtx$ [mm] & $\sigPVF\pm\sigPVFErr$ & $\sigAMVF\pm\sigAMVFErr$ & \\
68\% residual [$\mu$m] & \psePVF & \pseAMVF & \\
HS reconstruction, \hsWinMM{} mm [\%] & \hsRecoPVF & \hsRecoAMVF & \\
\midrule
Efficiency on r16638 (\nEvtPUb{} events) [\%] & \effPVFb & \effAMVFb & \\
\bottomrule
\end{tabular}
\end{center}
\src{HLDOC; \allowbreak{}H1 = H/\allowbreak{}h1\_\allowbreak{}vertex\_\allowbreak{}finding/\allowbreak{}h1\_\allowbreak{}summary.json (table3.tex)}
\end{frame}

\begin{frame}{HL-LHC before the GNN: Efficiency and Fakes versus $\mu$}
\mainbadge
\twofig{eff_vs_mu}{Vertex-finding efficiency, Eq.~(1), \winMM{} mm, r16443 entries 0--24,999.}%
{fake_fraction_vs_mu}{Fake fraction: reco vertices without truth within \winMM{} mm.}
\end{frame}

\begin{frame}{HL-LHC before the GNN: Fakes and Yield versus $\mu$}
\mainbadge
\twofig{fakes_per_event_vs_mu}{Fake vertices per event.}%
{yield_vs_mu}{Reconstructed vertices per event, with truth (draft Fig.~6).}
\end{frame}

\begin{frame}{HL-LHC before the GNN: Vertex Separation and Position at PU200}
\mainbadge
\twofig{deltaz_pu200}{$\Delta z$ between reconstructed vertices with the $\svtx$ fit (draft Fig.~7a).}%
{residuals_pu200}{Single-vertex residual $z_\mathrm{reco}-z_\mathrm{true}$, one-to-one within 1 mm (draft Fig.~7b).}
\end{frame}

\begin{frame}{HL-LHC before the GNN: Resolution versus $\mu$}
\mainbadge
\begin{columns}[T]
\begin{column}{.5\linewidth}\centering
\plot{p68_vs_mu}{.6}
\fcap{68\% residual versus $\mu$: the paper's resolution versus pile-up. \nEvtResMu{} events, r16443 and r16638.}
\end{column}
\begin{column}{.47\linewidth}
\footnotesize
\renewcommand{\arraystretch}{1.2}
\begin{tabular}{@{}lcc@{}}
\toprule
 & PV-Finder & AMVF \\
\midrule
68\% residual, $\mu$ 0--10 [$\mu$m] & \pseMuLoPVF & \pseMuLoAMVF \\
68\% residual, $\mu$ 180--200 [$\mu$m] & \pseMuHiPVF & \pseMuHiAMVF \\
$\svtx$, $\mu$ 100--120 [$\mu$m] & \sigMuHundredPVF & \sigMuHundredAMVF \\
$\svtx$, $\mu$ 180--200 [$\mu$m] & \sigMuTopPVF & \sigMuTopAMVF \\
AMVF / PVF $\svtx$, $\mu\ge100$ & \multicolumn{2}{c}{\sigRatioMu} \\
\bottomrule
\end{tabular}
\par\vspace{.5em}
$\svtx$ per bin only for $\mu\ge100$ (M10). The two-panel version is on the M10 slide.
\end{column}
\end{columns}
\end{frame}

\begin{frame}{HL-LHC before the GNN: Window Justification and Scan}
\twofig{window_justification}{\mainbadge{} $|z_\mathrm{reco}-z_\mathrm{true}|$ with the 68\% points and $3\times$ the larger one: the \winMM{} mm window.}%
{window_scan}{\backupbadge{} PV-Finder minus AMVF efficiency versus window.}
\end{frame}

\begin{frame}{HL-LHC before the GNN: Hard-Scatter Reconstruction}
\mainbadge
\begin{columns}[T]
\begin{column}{.58\linewidth}\centering
\plot{hs_reco_vs_mu}{.62}
\end{column}
\begin{column}{.38\linewidth}
\small
\fcap{Fraction of events with a reconstructed vertex within \hsWinMM{} mm of the true hard scatter, versus $\mu$.}
\par\vspace{.6em}
PU200: PV-Finder \hsRecoPVF\%, AMVF \hsRecoAMVF\%.
\end{column}
\end{columns}
\end{frame}

\begin{frame}{HL-LHC before the GNN: Backup}
\backupbadge
\twofig{eff_one_to_one_vs_mu}{One-to-one efficiency versus $\mu$.}%
{splits_per_event_vs_mu}{Split vertices per event versus $\mu$.}
\end{frame}

\begin{frame}{HL-LHC before the GNN: Positional Categories (Backup)}
\backupbadge
\twofig{pos_categories_pu200}{Positional categories at PU200, both algorithms.}%
{amvf_pos_categories_vs_mu}{AMVF positional categories versus $\mu$.}
\end{frame}

\begin{frame}{HL-LHC before the GNN: Efficiency vs $n_\mathrm{trk}$, Density, Separation (Backup)}
\backupbadge
\threefig{eff_ntrk_pre_pu200}{versus truth $n_\mathrm{trk}$, PU200.}%
{eff_density_pre_pu200}{versus local density, PU200.}%
{eff_dz_pre_pu200}{versus $\Delta z$ to the nearest truth vertex, PU200.}
\end{frame}

%---------------------------------------------------------------------- HL-LHC after the GNN
\begin{frame}{HL-LHC after the GNN: Table 4 at PU200}
\mainbadge\hfill{\footnotesize \hlModel{} + GNN v2 ($t=\hlT$) vs AMVF TTVA links, \nEvtPU{} events, vertices with $\ge1$ track}
\par\vspace{.2em}\scriptsize
\begin{center}
\renewcommand{\arraystretch}{1.1}
\begin{tabular}{@{}lcc@{}}
\toprule
 & PV-Finder + GNN & AMVF \\
\midrule
Clean / Merged / Split / Fake per event & \clnEvtGNN{} / \mrgEvtGNN{} / \splEvtGNN{} / \fakEvtGNN & \clnEvtAMVF{} / \mrgEvtAMVF{} / \splEvtAMVF{} / \fakEvtAMVF \\
Clean / Merged / Split / Fake per reco [\%] & \cleanRecoGNN{} / \mergedRecoGNN{} / \splitRecoGNN{} / \fakeRecoGNN & \cleanRecoAMVF{} / \mergedRecoAMVF{} / \splitRecoAMVF{} / \fakeRecoAMVF \\
Clean / truth [\%] & $\cleanTruthGNN\pm\cleanTruthGNNErr$ & \cleanTruthAMVF \\
\midrule
HS reconstructed and selected [\%] & $\hsSelGNN\pm\hsSelGNNErr$ & \hsSelAMVF{} (stored primary) \\
HS longitudinal resolution [$\mu$m] & $\hsResGNN\pm\hsResGNNErr$ & $\hsResAMVF\pm\hsResAMVFErr$ \\
Clean-HS fraction [\%] & \hsCleanGNN & \hsCleanAMVF \\
HS-track fraction [\%] & \hsTrkFracGNN & \hsTrkFracAMVF \\
\midrule
Track efficiency / purity / $F_1$ [\%] & \trkEffGNN{} / \trkPurGNN{} / \trkFOneGNN & \trkEffAMVF{} / \trkPurAMVF{} / \trkFOneAMVF \\
Wrong-assignment rate [\%] & \wrongGNN & \wrongAMVF \\
\midrule
All $\mu$ (25,000 events): Clean/truth [\%] & \cleanTruthGNNincl & \cleanTruthAMVFincl \\
All $\mu$: HS reconstructed and selected [\%] & \hsSelGNNincl & \hsSelAMVFincl \\
\bottomrule
\end{tabular}
\end{center}
\src{HLDOC; \allowbreak{}CH, CHI, SUP}
\end{frame}

\begin{frame}{HL-LHC after the GNN: Pile-up Vertex Quality at PU200}
\mainbadge
\begin{columns}[T]
\begin{column}{.6\linewidth}\centering
\plot{classes_pu200}{.62}
\end{column}
\begin{column}{.37\linewidth}
\small
\fcap{Clean, Merged, Split and Fake vertices per event at PU200, track-content classification (70\% rule), PV-Finder + GNN tracks vs AMVF TTVA links.}
\par\vspace{.5em}
Clean per event: \clnEvtGNN{} vs \clnEvtAMVF.
\end{column}
\end{columns}
\end{frame}

\begin{frame}{HL-LHC after the GNN: Vertex Categories versus $\mu$}
\mainbadge
\fourfig{clean_evt_vs_mu}{Clean / event}{merged_evt_vs_mu}{Merged / event}%
{split_evt_vs_mu}{Split / event}{fake_evt_vs_mu}{Fake / event}
\end{frame}

\begin{frame}{HL-LHC after the GNN: Clean/truth and Efficiency versus $n_\mathrm{trk}$}
\mainbadge
\twofig{clean_truth_vs_mu}{Clean vertices over truth vertices versus $\mu$.}%
{eff_ntrk_post_pu200}{Efficiency after the GNN versus truth $n_\mathrm{trk}$, PU200 (PUB-note style).}
\end{frame}

\begin{frame}{HL-LHC after the GNN: Hard-Scatter Selection and Resolution}
\mainbadge
\twofig{hs_sel_vs_mu}{HS reconstructed and selected (\hsWinMM{} mm), Athena rule for PV-Finder + GNN; AMVF stored primary and same rule.}%
{hs_res_vs_mu}{HS longitudinal resolution (iterative 3-RMS), 3D adaptive fit on GNN tracks vs AMVF.}
\end{frame}

\begin{frame}{HL-LHC after the GNN: Hard Scatter versus Local Pile-up Density}
\mainbadge
\threefig{hs_reco_vs_density}{HS reconstruction}%
{hs_sel_vs_density}{HS selection}%
{hs_res_vs_density}{HS resolution}
\par\vspace{.3em}{\footnotesize Density: truth vertices within $\pm2$ mm of the true HS, HS excluded, per mm (ITk). All $\mu$, r16443 entries 0--24,999.}
\end{frame}

\begin{frame}{HL-LHC after the GNN: Track Association}
\threefig{trk_f1_vs_mu}{\mainbadge{} Track $F_1$}%
{trk_eff_vs_mu}{\backupbadge{} Track efficiency}%
{trk_pur_vs_mu}{\backupbadge{} Track purity}
\par\vspace{.3em}{\footnotesize At PU200 the GNN trades efficiency (\trkEffGNN{} vs \trkEffAMVF\%) for purity (\trkPurGNN{} vs \trkPurAMVF\%); $F_1$ \trkFOneGNN{} vs \trkFOneAMVF\%. At low $\mu$ the GNN $F_1$ is below AMVF (M5).}
\end{frame}

\begin{frame}{HL-LHC after the GNN: Hard-Scatter Tracks (Backup)}
\backupbadge
\twofig{hs_clean_vs_mu}{Clean-HS fraction versus $\mu$.}%
{hs_trkfrac_vs_mu}{HS tracks on the HS-matched vertex over all HS tracks.}
\end{frame}

\begin{frame}{HL-LHC after the GNN: HS Resolution vs $n_\mathrm{trk}$, Density Definition (Backup)}
\backupbadge
\twofig{hs_res_vs_ntrk}{HS resolution versus HS truth $n_\mathrm{trk}$, PU200.}%
{density_definitions}{Local density: ITk definition (all truth vertices) and the stored-truth one.}
\end{frame}

\begin{frame}{HL-LHC after the GNN: Classification Cross-checks (Backup)}
\backupbadge
\twofig{classes_variants}{Categories per event at PU200: unweighted (headline), GNN-score weights, AMVF fit tracks with fit weights.}%
{amvf_fit_weights}{AMVF fit-weight distribution.}
\par\vspace{.2em}{\footnotesize With its own fit weights AMVF has \clnEvtAMVFfit{} Clean per event; PV-Finder + GNN leads by \cleanLeadFit$\times$ instead of \cleanLeadTTVA$\times$.}
\end{frame}

\begin{frame}{HL-LHC after the GNN: Association Baselines (Backup)}
\backupbadge
\twofig{eq4_vs_mu}{Eq.~(4) significance association, Clean/truth versus $\mu$.}%
{z0_baselines_vs_mu}{Nearest $z_0$ and $3\sigma(z_0)$ baselines.}
\end{frame}

\begin{frame}{HL-LHC after the GNN: Bounds and Stages (Backup)}
\backupbadge
\threefig{truth_guided_vs_mu}{Truth-guided association on the same peaks.}%
{finder_cap_vs_mu}{Finder cap (one to one, \winMM{} mm).}%
{stages_pu200}{Clean yield by stage, PU200.}
\end{frame}

\begin{frame}{HL-LHC after the GNN: Unassigned and Wrong Tracks (Backup)}
\backupbadge
\threefig{wrong_rate_vs_mu}{Wrong-assignment rate ($1-$ purity).}%
{lost_unassigned_vs_mu}{Labelled tracks left unassigned with no candidate within \winMM{} mm.}%
{unlabelled_unassigned_vs_mu}{Unlabelled tracks correctly left unassigned.}
\end{frame}

\begin{frame}{HL-LHC after the GNN: Density, Separation, Peak Width (Backup)}
\backupbadge
\threefig{eff_density_post_pu200}{Efficiency after the GNN versus local density.}%
{eff_dz_post_pu200}{versus $\Delta z$ to the nearest truth vertex.}%
{sigma_pv_pu200}{Peak width ($\sigma_\mathrm{PV}$) vs the $n_\mathrm{trk}$ parameterisation.}
\end{frame}

\begin{frame}{GNN v2: Learning Curve and Metrics versus $\mu$ (Backup)}
\backupbadge
\twofig{gnn_learning_curve}{Validation-monitor loss, v1 and v2.}%
{vs_mu_test_s8}{All chain metrics versus $\mu$ for several $t$, \hlModel{} peaks, test events.}
\end{frame}

%---------------------------------------------------------------------- Low pile-up
\begin{frame}{Low Pile-up: Vertex Finding (Table 6)}
\mainbadge\hfill{\footnotesize \lpModel, integral \lpIntegral, window \winMM{} mm}
\par\vspace{.4em}\footnotesize
\begin{center}
\renewcommand{\arraystretch}{1.2}
\begin{tabular}{@{}lccc@{}}
\toprule
Run 2 MC, \lpMCTest{} events, \lpTruthMC{} truth vertices & PV-Finder & AMVF & PVF $-$ AMVF \\
\midrule
Efficiency [\%] & $\lpEffPVF\pm\lpEffPVFErr$ & \lpEffAMVF & \lpEffDiff{} pt \\
Fake / reco [\%] & \lpFakeFracPVF & \lpFakeFracAMVF & \\
Fakes / event & \lpFakeEvtPVF & \lpFakeEvtAMVF & \\
$\svtx$ [mm] & $\lpSigPVF\pm\lpSigPVFErr$ & $\lpSigAMVF\pm\lpSigAMVFErr$ & \\
68\% residual [$\mu$m] & \lpPsePVF & \lpPseAMVF & \\
HS reconstruction, \hsWinMM{} mm [\%] & \lpHSPVF & \lpHSAMVF & \\
\bottomrule
\end{tabular}
\par\vspace{.6em}
\begin{tabular}{@{}lccc@{}}
\toprule
Collision data, AMVF as reference & Run 2 data & Run 3 data & Run 2 MC, same convention \\
\midrule
AMVF-vertex recovery [\%] & \lpRecRunTwo & \lpRecRunThree & \lpRecMCref \\
PVF vertices without AMVF vertex / event & \lpNoRefRunTwo & \lpNoRefRunThree & \\
$\svtx$ PV-Finder / AMVF [mm] & \lpSigRunTwo{} / \lpSigAMVFRunTwo & \lpSigRunThree{} / \lpSigAMVFRunThree & \\
\bottomrule
\end{tabular}
\end{center}
\src{LPDOC R1, R3; \allowbreak{}VF}
\end{frame}

\begin{frame}{Low Pile-up: Run 2 MC Vertex Finding (Fig.~8)}
\mainbadge
\twofig{lp_fig8a}{(a) $\Delta z$ between vertices, PV-Finder and AMVF.}%
{lp_fig8b}{(b) Vertices per event versus $\mu$.}
\end{frame}

\begin{frame}{Low Pile-up: Run 2 MC after the GNN (Table 5, Fig.~9)}
\mainbadge
\begin{columns}[T]
\begin{column}{.36\linewidth}
\footnotesize
\renewcommand{\arraystretch}{1.18}
\begin{tabular}{@{}lcc@{}}
\toprule
{[\%]} & PVF + GNN & AMVF \\
\midrule
Clean & \lpClean & \lpCleanAMVF \\
Merged & \lpMerged & \lpMergedAMVF \\
Split & \lpSplit & \lpSplitAMVF \\
Fake & \lpFake & \lpFakeAMVF \\
Clean/truth & \lpCleanTruth & \lpCleanTruthAMVF \\
Track $F_1$ & \lpFOne & \lpFOneAMVF \\
HS ident. & \lpHSID & \lpHSIDAMVF \\
\bottomrule
\end{tabular}
\par\vspace{.3em}{\scriptsize $t=\lpT$; per reco vertex with $\ge1$ track; \lpEmpty{} empty candidates dropped.}
\end{column}
\begin{column}{.31\linewidth}\centering
\plot{lp_fig9a}{.5}\fcap{(a) Matched-vertex efficiency vs truth $n_\mathrm{trk}$.}
\end{column}
\begin{column}{.31\linewidth}\centering
\plot{lp_fig9b}{.5}\fcap{(b) Categories.}
\end{column}
\end{columns}
\end{frame}

\begin{frame}{Low Pile-up: Collision Data (Fig.~10)}
\mainbadge
\fourfig{lp_fig10a}{(a) Run 2 data $\Delta z$, with Run 2 MC}{lp_fig10b}{(b) Run 2 data vertices vs $\mu$}%
{lp_fig10c}{(c) Run 3 data $\Delta z$, with Run 2 MC}{lp_fig10d}{(d) Run 3 data vertices vs $\mu$}
\end{frame}

\begin{frame}{Low Pile-up: Run 3 Data Association and the $t$ Scan}
\begin{columns}[T]
\begin{column}{.45\linewidth}
\mainbadge\par\vspace{.4em}
\footnotesize
\renewcommand{\arraystretch}{1.2}
\begin{tabular}{@{}lc@{}}
\toprule
Run 3 data, 2,500 events, $t=\lpT$ & [\%] \\
\midrule
Tracks assigned by the GNN & \lpRthreeAssigned \\
Agreement with AMVF (1 mm) & \lpRthreeAgree \\
Agreement with AMVF (\winMM{} mm) & \lpRthreeAgreeW \\
Leading-vertex agreement & \lpRthreeLead \\
\bottomrule
\end{tabular}
\par\vspace{.3em}{\scriptsize Same events and peaks as the vertex finding. Agreement, not efficiency.}
\end{column}
\begin{column}{.52\linewidth}\centering
\backupbadge\par
\plot{lp_t_scan}{.55}\fcap{Run 2 MC metrics versus $t$.}
\end{column}
\end{columns}
\end{frame}

%---------------------------------------------------------------------- Timing
\begin{frame}{Timing: CPU and GPU}
\begin{columns}[T]
\begin{column}{.55\linewidth}\centering
\mainbadge\par
\plot{athena_timing}{.6}
\fcap{Athena timing versus pile-up: AMVF on CPU, PV-Finder chain on GPU (note figure).}
\end{column}
\begin{column}{.42\linewidth}
\footnotesize
\renewcommand{\arraystretch}{1.2}
\begin{tabular}{@{}lc@{}}
\toprule
$\mu=200$, per event & [ms] \\
\midrule
PV-Finder chain, GPU (Athena) & \athGPUms \\
PV-Finder chain, CPU (Athena, v6) & \athCPUms \\
\hlModel{} chain, CPU, projected & \projChain \\
AMVF, CPU (note figure) & \athAMVF \\
\bottomrule
\end{tabular}
\par\vspace{.5em}
The CPU claim in the paper needs packed inference in Athena (M6). The GPU claim stands as measured.
\end{column}
\end{columns}
\end{frame}


%======================================================================
\partdivider{7}{Proposed Internal-Note Edits}{Old text and new [MM: ...] text, with page and line}

% Proposed internal-note edits. To be filled by the note-edit agent:
% replace each \noteedit placeholder by the page, line, old text and the new
% red [MM: ...] text. The topics are those of
% docs/additional_studies/meeting_2026-11-11_plan.md ("Note edits").

\newcommand{\editbox}[2]{% label, colour
  {\setlength{\fboxsep}{5pt}\fcolorbox{#2!60!black}{#2!4}{\parbox[c][.22\textheight][t]{.93\linewidth}{%
  {\scriptsize\bfseries\textcolor{#2!60!black}{#1}}\par\vspace{.3em}{\footnotesize\textcolor{red}{TBD}}}}}}
\newcommand{\noteedit}[2]{% number, topic
\begin{frame}{Note Edit #1: #2}
\vspace*{.8em}{\small\textbf{Note page, line}: \textcolor{red}{TBD} \hfill \textbf{Section}: \textcolor{red}{TBD}}
\par\vspace{.5em}\centering
\editbox{Old text}{black}\par\vspace{.6em}
\editbox{New text, in red [MM: ...]}{red}
\end{frame}}

\begin{frame}{Note Edits: List}
\small
Internal note, GitLab main. Changes in red as \textcolor{red}{[MM: ...]} on a branch with a merge request.
\begin{enumerate}\setlength\itemsep{.3em}
\item The \winMM{} mm window with its justification plot; the hard scatter also at \hsWinMM{} mm.
\item Split-inclusive efficiency: one sentence.
\item HS longitudinal resolution from the 3D adaptive fit: one paragraph and one figure.
\item Timing claim restricted to GPU.
\item Simulation on the Fig.~6.2 panels, and the AMVF-reference bias.
\item $z_0$ baselines in Fig.~7.2.
\item AMVF HS selection \hsSelAMVF\% and the HS numbers.
\item The AMVF-weights paragraph.
\item The frozen configuration: \hlModel, GNN v2 epoch \hlGnnEpoch, $t=\hlT$; \lpModel{} at low pile-up.
\end{enumerate}
\end{frame}

\noteedit{1}{Matching Window}
\noteedit{2}{Split-Inclusive Efficiency}
\noteedit{3}{Hard-Scatter Longitudinal Resolution}
\noteedit{4}{Timing Claim}
\noteedit{5}{Fig.~6.2 and the AMVF Reference}
\noteedit{6}{$z_0$ Baselines in Fig.~7.2}
\noteedit{7}{AMVF Hard-Scatter Selection}
\noteedit{8}{AMVF Weights}
\noteedit{9}{Frozen Configuration}


\end{document}
")
%
% \readingonly{...} keeps its argument only in the reading build; \talkonly{...}
% only in the talk build.  Every number lives in exactly one place, so the two
% versions cannot drift apart.
\providecommand{\talkbuild}{0}
\newcommand{\readingonly}[1]{\ifnum\talkbuild=0\relax#1\fi}
\newcommand{\talkonly}[1]{\ifnum\talkbuild=1\relax#1\fi}
% \bnum{reading}{talk} expands to a bare number, so it can be used inside a
% length: \includegraphics[height=\bnum{0.5}{0.6}\textheight]{...}.  The talk
% build drops prose, which frees vertical room for the figure that is left.
\newcommand{\bnum}[2]{\ifnum\talkbuild=0 #1\else #2\fi}
% ---------------------------------------------------------------------------

% Part-divider helper: \partdivider{numeral}{title}{one-line subtitle}
\newcommand{\partdivider}[3]{%
    \begin{frame}[plain]
        \centering
        \vfill
        \begin{tikzpicture}
            \fill[pennred] (0,0) circle (0.72cm);
            \node[text=white, font=\Large\bfseries] at (0,0.02) {#1};
        \end{tikzpicture}\\[0.85em]
        {\LARGE\textbf{\textcolor{pennred}{#2}}}\\[0.75em]
        {\color{pennred!28}\rule{0.30\linewidth}{1.1pt}}\\[0.75em]
        {\small\textcolor{black!50}{#3}}
        \vfill
    \end{frame}%
}

%%%%%%%%%%%%%%%%%%%%%%%%%%%%%%%%%%%%%%

\usetheme{Philadelphia}
\hypersetup{colorlinks,allcolors=pennred}

\title[PV-Finder]{PV-Finder: primary vertex reconstruction and track association with neural networks}
\subtitle{Paper deck: configuration, definitions, answers and figures}

\author[Mattia Mauro]{Qi-Bin Lei, Rocky B. Garg, Lauren Tompkins,\\\textcolor{pennred}{Mattia Mauro}, Trey Warner}

\institute[]{Stanford University}

\date{November 13, 2026}

\titlegraphic{
    \hspace{5pt}
    \href{https://stanford.edu/group/stanford_atlas/}{\includegraphics[height=1cm]{logo-Stanford/Block_S_2_color.png}}
    \href{https://atlas.cern/}{\includegraphics[trim={0 11.5cm 0 11.5cm},clip,height=1cm]{logo-ATLAS/ATLAS-Logo-Ref-Pantone.pdf}}
    \href{https://iris-hep.org/}{\includegraphics[height=1cm]{logo-iris-hep/Iris-hep-4-no-long-name.png}}
}


\newcommand{\readingnote}[1]{\readingonly{\par\vspace{.35em}%
  {\footnotesize\color{black!65}#1\par}}}
\newcommand{\takeaway}[1]{\par\vspace{.65em}{\centering
  \setlength{\fboxsep}{7pt}\colorbox{pennred!5}{\parbox{.94\linewidth}{%
  \centering\small\textcolor{pennred}{#1}}}\par}}
\newcommand{\context}[1]{{\centering\footnotesize\color{black!55}#1\par}\vspace{.65em}}
\newcommand{\bigvalue}[2]{{\LARGE\bfseries\textcolor{pennred}{#1}}\\[.4em]{\small #2}\par}
\newcommand{\pubnote}{\href{https://cds.cern.ch/record/2858348}{ATL-PHYS-PUB-2023-011}}
\newcommand{\notelink}{\href{https://gitlab.cern.ch/atlas-physics-office/IDTR/ANA-IDTR-2026-01/ANA-IDTR-2026-01-INT1}{ANA-IDTR-2026-01-INT1}}
\newcommand{\treysource}{\href{https://gitlab.cern.ch/atlas-physics-office/IDTR/ANA-IDTR-2026-01/ANA-IDTR-2026-01-INT1/-/blob/4c50ed51500c897dd6a1b5bb52b0f1e3c12e7575/sections/04b_timing.tex}{Note Sec. 5.1 (18 Sep 2026)}}

\newcommand{\svtx}{\sigma_{\mathrm{vtx-vtx}}}
\setbeamertemplate{navigation symbols}{}
\hypersetup{pdftitle={PV-Finder, paper deck, November 13, 2026},
  pdfauthor={Mattia Mauro, Qi-Bin Lei, Rocky B. Garg, Lauren Tompkins, Trey Warner}}

% ---------------------------------------------------------------------------
% Paper deck helpers
\usepackage{array}
% \imgsrc{name}{expected source path}{ext}: written by build.sh into
% images_manifest.tex. \plot{name}{height} includes images/name.ext when the
% figure exists and draws a framed placeholder with the expected path if not.
\makeatletter
\newcommand{\imgsrc}[3]{%
  \expandafter\def\csname pd@src@#1\endcsname{#2}%
  \expandafter\def\csname pd@ext@#1\endcsname{#3}}
\newcommand{\plot}[2]{%
  \expandafter\ifx\csname pd@ext@#1\endcsname\relax
    \pd@placeholder{#1}{#2}{(not in images.txt)}%
  \else
    \edef\pd@e{\csname pd@ext@#1\endcsname}%
    \ifx\pd@e\@empty
      \pd@placeholder{#1}{#2}{\csname pd@src@#1\endcsname.\{pdf,png\}}%
    \else
      \includegraphics[width=\linewidth,height=#2\textheight,keepaspectratio]{images/#1.\pd@e}%
    \fi
  \fi}
\newcommand{\pd@placeholder}[3]{%
  {\setlength{\fboxsep}{6pt}\setlength{\fboxrule}{0.8pt}%
   \fcolorbox{red!60!black}{red!3}{\parbox[c][#2\textheight][c]{.9\linewidth}{\centering
   {\small\bfseries\textcolor{red!60!black}{Figure not produced yet}}\\[.5em]
   {\scriptsize\ttfamily #3}\\[.4em]
   {\tiny\textcolor{black!55}{image \texttt{\detokenize{#1}}; rebuild with build.sh when it exists}}}}}}
\makeatother
% Role badges for the results section
\newcommand{\mainbadge}{{\setlength{\fboxsep}{2pt}\colorbox{mygreen!15}{\scriptsize\bfseries\textcolor{mygreen!70!black}{PAPER MAIN}}}}
\newcommand{\backupbadge}{{\setlength{\fboxsep}{2pt}\colorbox{black!8}{\scriptsize\bfseries\textcolor{black!60}{BACKUP / NOTE}}}}
\newcommand{\fcap}[1]{\par\vspace{.25em}{\footnotesize\color{black!70}#1\par}}
\newcommand{\src}[1]{{\raggedright\tiny\color{black!50}\texttt{#1}\par}}
\newcommand{\draftref}[1]{{\small\textcolor{pennred}{\textbf{Draft #1}}}}
\newcommand{\qtext}[1]{{\footnotesize\itshape\color{black!65}#1\par}\vspace{.4em}}
\newcommand{\final}[1]{\textcolor{mydarkblue}{\textbf{Final:}} #1}

% numbers.tex: every number of the paper deck, one macro each.
%
% Rule: \TBD (red) marks a number that is not final yet. Each TBD line names
% the doc and the JSON key it is filled from. "seen" gives the value present in
% that file on 12-13 Nov 2026, before the recompute agents finished: it is for
% checking the fill only, never for quoting.
%
% Sources:
%   HLDOC = docs/evaluation/hllhc_final_numbers_2026-11-12.md (HL-LHC agent)
%   LPDOC = docs/evaluation/lowpu_final_numbers_2026-11-12.md (low pile-up agent)
%   H = outputs/11_12_2026_output/hllhc_final    L = outputs/11_12_2026_output/lowpu_final
%   CH  = H/chain/chain_headline.json :: headline/PU200 (185-215)/...
%   CHI = H/chain/chain_headline.json :: headline/inclusive/...
%   SUP = H/supplementary/supplementary_summary.json
%   CX  = H/chain_extras/chain_extras_summary.json
%   H1  = H/h1_vertex_finding/h1_summary.json (r16443, PU200, 1,920 events)
%   RVM = H/resolution_vs_mu/results.json
%   CPU = H/cpu_timing/cpu_timing.json (and H/cpu_timing/README.md)
%   VF  = L/summary/vf_results.json   GM = L/gnn_mc/summary.json :: bins/inclusive/...
%   R3A = L/run3_assoc/summary.json :: by_t/0.97/...
% Final numbers without a TBD come from committed code or committed docs.
\providecommand{\TBD}{\textcolor{red}{\textnormal{TBD}}}

%======================================================================
% 1. Frozen configuration (src/pv_finder/utils/working_points.py, final)
%======================================================================
\newcommand{\winMM}{0.17}
\newcommand{\hsWinMM}{0.1}
\newcommand{\assignMM}{1}
\newcommand{\pairBins}{300}
\newcommand{\purityCut}{70\%}
% HL-LHC density model s8
\newcommand{\hlModel}{s8}
\newcommand{\hlCkpt}{hllhc\_alleta\_v6\_s8\_\allowbreak 1p1ep\_phase2\_\allowbreak step\_450000}
\newcommand{\hlUnetCh}{80}
\newcommand{\hlLatent}{1}
\newcommand{\hlHidden}{$5\times128$}
% training log outputs/logs/hllhc_alleta_v6_s8_1p1ep_20261109_211919.log
\newcommand{\hlParams}{439{,}593}
% derived: MLP 7-128x5-1000 = 196,072; U-Net = 439,593 - 196,072
\newcommand{\hlParamsMLP}{196{,}072}
\newcommand{\hlParamsUnet}{243{,}521}
\newcommand{\hlAmp}{0.01}
\newcommand{\hlWidth}{3}
\newcommand{\hlIntegral}{0.184}
\newcommand{\hlHeight}{0.0184}
\newcommand{\hlCentroid}{3}
% HL-LHC GNN
\newcommand{\hlGnnEpoch}{527}
\newcommand{\hlT}{0.97}
\newcommand{\gnnParams}{39{,}521}
\newcommand{\gnnK}{20}
% low pile-up
\newcommand{\lpModel}{PVF-e400}
\newcommand{\lpCkpt}{e2e\_mlpHist50\_e2e400\_\allowbreak 1latent\_mse\_\allowbreak phase2\_epoch\_400}
\newcommand{\lpParams}{358{,}713}
\newcommand{\lpUnetCh}{64}
\newcommand{\lpHidden}{$5\times100$}
\newcommand{\lpIntegral}{0.50}
\newcommand{\lpGnnCkpt}{ttva\_gat\_mu60\_v2\_aug epoch 115}
% working_points.LOWPU.gnn.t (frozen 12 Nov; LPDOC Sec. 3)
\newcommand{\lpT}{0.97}

%======================================================================
% 2. Datasets (final)
%======================================================================
% docs/additional_studies/paper_questions_2026-11-11.md Part 3;
% outputs/11_11_2026_output/datasets/hllhc_dataset_inventory.json
\newcommand{\nPoolEvt}{5{,}173{,}600}
\newcommand{\nTrainSub}{59{,}599{,}872}
\newcommand{\nValSub}{1{,}862{,}496}
\newcommand{\nTestSub}{620{,}832}
\newcommand{\nTrainEq}{4{,}966{,}656}
\newcommand{\nValEq}{155{,}208}
\newcommand{\nTestEq}{51{,}736}
\newcommand{\nEvtPU}{1{,}920}
\newcommand{\nEvtPUb}{1{,}888}
\newcommand{\nTruthPU}{213{,}738}
\newcommand{\truthEvt}{111.3}
\newcommand{\meanMuPU}{192.5}
% GNN v2: outputs/11_11_2026_output/gnn_fullmu_v2/split_manifest.json
\newcommand{\gnnTrainA}{60{,}000}
\newcommand{\gnnTrainB}{89{,}820}
\newcommand{\gnnTrain}{149{,}820}
\newcommand{\gnnVal}{7{,}500}
\newcommand{\gnnMon}{1{,}500}
\newcommand{\gnnResTest}{7{,}500}
\newcommand{\gnnResTestB}{9{,}980}
% resolution vs mu sample (docs/evaluation/resolution_vs_mu_2026-11-11.md)
\newcommand{\nEvtResMu}{44{,}999}
% low pile-up: outputs/11_12_2026_output/paper_deck/lowpu_mu/lowpu_mu_distributions.json
\newcommand{\lpMCTotal}{51{,}000}
\newcommand{\lpMCTrain}{35{,}700}
\newcommand{\lpMCVal}{12{,}750}
\newcommand{\lpMCTest}{2{,}550}
\newcommand{\lpMCNonTest}{48{,}450}
\newcommand{\lpMuMC}{38.9}
\newcommand{\lpMuRunTwo}{42.5}
\newcommand{\lpMuRunThree}{50.0}
\newcommand{\lpTruthMC}{72{,}189}

%======================================================================
% 3. HL-LHC before the GNN, PU200 (r16443, 1,920 events)
%======================================================================
% AMVF numbers do not depend on the density model: final, from
% docs/evaluation/model_freeze_candidates_2026-11-12.md (0.17 mm).
\newcommand{\effAMVF}{79.85}
\newcommand{\effAMVFb}{80.13}
\newcommand{\fakeFracAMVF}{20.08}
\newcommand{\fakeEvtAMVF}{19.65}
\newcommand{\splitEvtAMVF}{0.184}
\newcommand{\recoEvtAMVF}{97.9}
\newcommand{\sigAMVF}{0.305}
\newcommand{\sigAMVFErr}{0.006}
\newcommand{\pseAMVF}{56.45}
\newcommand{\hsRecoAMVF}{99.79}
% window derivation (v6 and AMVF, docs/evaluation/matching_window_audit_2026-11-09.md)
\newcommand{\pseVsix}{55.05}
% PV-Finder s8 at the frozen cuts: TBD
\newcommand{\effPVF}{\TBD}      % HLDOC | H1 Eq.(1) eff r16443 | seen: freeze doc 83.05 (v6 83.20 - 0.15)
\newcommand{\effPVFErr}{\TBD}   % HLDOC | H1 bootstrap error
\newcommand{\effDiff}{\TBD}     % HLDOC | H1 PVF - AMVF, paired [pt] | seen: about +3.2
\newcommand{\effDiffErr}{\TBD}  % HLDOC | H1 paired bootstrap error
\newcommand{\effPVFb}{\TBD}     % HLDOC | H1 Eq.(1) eff r16638 (1,888 events) | seen: 83.21
\newcommand{\effOneOnePVF}{\TBD}  % HLDOC | H1 one-to-one efficiency, PVF
\newcommand{\effOneOneAMVF}{\TBD} % HLDOC | H1 one-to-one efficiency, AMVF (0.1694 mm: 70.07)
\newcommand{\fakeFracPVF}{\TBD} % HLDOC | H1 fake / reco [%] | seen: about 18.5
\newcommand{\fakeEvtPVF}{\TBD}  % HLDOC | H1 fakes / event | seen: 18.64
\newcommand{\splitEvtPVF}{\TBD} % HLDOC | H1 splits / event
\newcommand{\recoEvtPVF}{\TBD}  % HLDOC | H1 reco / event
\newcommand{\sigPVF}{\TBD}      % HLDOC | H1 sigma_vtx-vtx [mm], counted peaks | seen: 0.217 at 0.182 cuts
\newcommand{\sigPVFErr}{\TBD}   % HLDOC | H1 fit error
\newcommand{\psePVF}{\TBD}      % HLDOC | H1 68% residual [um] | seen: 55.60 at 0.182 cuts
\newcommand{\hsRecoPVF}{\TBD}   % HLDOC | H1 / CH HS reco at 0.1 mm [%] | seen: 99.90 (1,918)
\newcommand{\chancePVF}{\TBD}   % HLDOC | H1 3 mm-shift chance match [%] (v6: 18.6)
\newcommand{\chanceAMVF}{\TBD}  % HLDOC | H1 3 mm-shift chance match [%] (0.1694 mm: 18.4)
\newcommand{\scanLo}{\TBD}      % HLDOC | H1 window scan: smallest PVF-AMVF over 0.10-0.20 mm (v6: +3.1)
\newcommand{\scanHi}{\TBD}      % HLDOC | H1 window scan: largest PVF-AMVF over 0.10-0.20 mm (v6: +3.4)
% resolution vs mu (s8 rerun of the 11 Nov study; AMVF final from that study)
\newcommand{\pseMuLoAMVF}{32.1}
\newcommand{\pseMuHiAMVF}{56.5}
\newcommand{\sigMuHundredAMVF}{275}
\newcommand{\sigMuTopAMVF}{304}
\newcommand{\pseMuLoPVF}{\TBD}      % HLDOC | RVM bin 0-10 p68 [um] (v6: 33.7)
\newcommand{\pseMuHiPVF}{\TBD}      % HLDOC | RVM bin 180-200 p68 [um] (v6: 54.9)
\newcommand{\sigMuHundredPVF}{\TBD} % HLDOC | RVM bin 100-120 sigma [um] (v6: 194.8)
\newcommand{\sigMuTopPVF}{\TBD}     % HLDOC | RVM bin 180-200 sigma [um] (v6: 216.2)
\newcommand{\sigRatioMu}{\TBD}      % HLDOC | RVM AMVF/PVF sigma range for mu >= 100 (v6: 1.41-1.43)

%======================================================================
% 4. HL-LHC after the GNN (s8 + GNN v2 e527, t = 0.97)
%======================================================================
% AMVF, model independent, final (gnn_v2_freeze_2026-11-12.md test r16443;
% metric_definitions.md Sec. 1; supplementary counts 66,127 / 80,104 / 41,594 / 135 over 1,920)
\newcommand{\cleanTruthAMVF}{30.94}
\newcommand{\cleanTruthAMVFincl}{40.49}
\newcommand{\cleanRecoAMVF}{35.2}
\newcommand{\mergedRecoAMVF}{42.6}
\newcommand{\splitRecoAMVF}{22.1}
\newcommand{\fakeRecoAMVF}{0.07}
\newcommand{\clnEvtAMVF}{34.4}
\newcommand{\mrgEvtAMVF}{41.7}
\newcommand{\splEvtAMVF}{21.7}
\newcommand{\fakEvtAMVF}{0.07}
\newcommand{\hsSelAMVF}{92.92}
\newcommand{\hsSelAMVFrule}{89.11}
\newcommand{\hsSelAMVFplain}{78.80}
\newcommand{\hsSelAMVFincl}{95.14}
\newcommand{\hsResAMVF}{10.24}
\newcommand{\hsResAMVFErr}{0.27}
\newcommand{\hsCleanAMVF}{92.5}
\newcommand{\hsTrkFracAMVF}{74.1}
% PV-Finder + GNN: TBD
\newcommand{\cleanTruthGNN}{\TBD}     % HLDOC | CH clean_per_truth/gnn | seen 66.72
\newcommand{\cleanTruthGNNErr}{\TBD}  % HLDOC | CH clean_per_truth/gnn_err | seen 0.13
\newcommand{\cleanTruthGNNincl}{\TBD} % HLDOC | CHI clean_per_truth/gnn | seen 73.09
\newcommand{\cleanRecoGNN}{\TBD}      % HLDOC | CH clean_per_reco/gnn | seen 79.03
\newcommand{\mergedRecoGNN}{\TBD}     % HLDOC | CH merged_per_reco/gnn | seen 13.55
\newcommand{\splitRecoGNN}{\TBD}      % HLDOC | CH split_per_reco/gnn | seen 7.24
\newcommand{\fakeRecoGNN}{\TBD}       % HLDOC | CH fake_per_reco/gnn | seen 0.18
\newcommand{\clnEvtGNN}{\TBD}         % HLDOC | SUP classes, PU200, gnn unweighted, Clean/event | seen 74.28
\newcommand{\mrgEvtGNN}{\TBD}         % HLDOC | SUP Merged/event | seen 12.73
\newcommand{\splEvtGNN}{\TBD}         % HLDOC | SUP Split/event | seen 6.81
\newcommand{\fakEvtGNN}{\TBD}         % HLDOC | SUP Fake/event | seen 0.17
\newcommand{\clnEvtGNNw}{\TBD}        % HLDOC | SUP gnn, GNN-score weights, Clean/event | seen 74.34
\newcommand{\clnEvtAMVFfit}{\TBD}     % HLDOC | SUP amvf fit tracks + fit weights, Clean/event | seen 46.26
\newcommand{\mrgEvtAMVFfit}{\TBD}     % HLDOC | SUP | seen 31.98
\newcommand{\splEvtAMVFfit}{\TBD}     % HLDOC | SUP | seen 19.56
\newcommand{\cleanLeadFit}{\TBD}      % HLDOC | SUP Clean ratio GNN/AMVF with AMVF fit weights | seen 1.61
\newcommand{\cleanLeadTTVA}{\TBD}     % HLDOC | SUP Clean ratio GNN/AMVF with TTVA links | seen 2.16
\newcommand{\emptyDropped}{\TBD}      % HLDOC | CH table4/empty_candidates_dropped | seen 13,204
\newcommand{\hsSelGNN}{\TBD}          % HLDOC | CH hs_selection/gnn_rule | seen 94.95
\newcommand{\hsSelGNNErr}{\TBD}       % HLDOC | CH hs_selection/gnn_rule_err | seen 0.53
\newcommand{\hsSelDiff}{\TBD}         % HLDOC | CH gnn_rule - amvf_stored [pt] | seen +2.03
\newcommand{\hsSelDiffErr}{\TBD}      % HLDOC | CH hs_selection/gnn_minus_stored_err | seen 0.56
\newcommand{\hsSelGNNincl}{\TBD}      % HLDOC | CHI hs_selection/gnn_rule | seen 96.94
\newcommand{\hsResGNN}{\TBD}          % HLDOC | CH hs_resolution_um/gnn | seen 9.48
\newcommand{\hsResGNNErr}{\TBD}       % HLDOC | CH hs_resolution_um/gnn_err | seen 0.28
\newcommand{\hsResN}{\TBD}            % HLDOC | CH hs_resolution_um/n_common | seen 1,915
\newcommand{\hsResGNNincl}{\TBD}      % HLDOC | CHI hs_resolution_um/gnn | seen 9.57
\newcommand{\hsCleanGNN}{\TBD}        % HLDOC | CH hs_clean/gnn | seen 94.7
\newcommand{\hsTrkFracGNN}{\TBD}      % HLDOC | CH hs_track_frac/gnn | seen 80.9
\newcommand{\trkEffGNN}{\TBD}         % HLDOC | CH trk_eff/gnn | seen 57.45
\newcommand{\trkPurGNN}{\TBD}         % HLDOC | CH trk_pur/gnn | seen 85.87
\newcommand{\trkFOneGNN}{\TBD}        % HLDOC | CH trk_f1/gnn | seen 68.84
\newcommand{\trkEffAMVF}{\TBD}        % HLDOC | CH trk_eff/amvf | seen 68.32
\newcommand{\trkPurAMVF}{\TBD}        % HLDOC | CH trk_pur/amvf | seen 68.65
\newcommand{\trkFOneAMVF}{\TBD}       % HLDOC | CH trk_f1/amvf | seen 68.48
\newcommand{\wrongGNN}{\TBD}          % HLDOC | CH trk_wrong_rate/gnn | seen 14.1
\newcommand{\wrongAMVF}{\TBD}         % HLDOC | CH trk_wrong_rate/amvf | seen 31.4
\newcommand{\assignedGNN}{\TBD}       % HLDOC | CH assigned_frac/gnn | seen 66.5
\newcommand{\lostUnassigned}{\TBD}    % HLDOC | CX lost-unassigned fraction, PU200, gnn_t
\newcommand{\unlabelledLeft}{\TBD}    % HLDOC | CX unlabelled tracks left unassigned, PU200, gnn_t
\newcommand{\eqFourBest}{\TBD}        % HLDOC | CX Eq.(4) best Clean/truth, PU200 | seen 33.5 (s_max = 1)
\newcommand{\eqFourThree}{\TBD}       % HLDOC | CX Eq.(4) s_max = 3 | seen 32.9
\newcommand{\zNear}{\TBD}             % HLDOC | CX z0_nearest Clean/truth | seen 31.9
\newcommand{\zThree}{\TBD}            % HLDOC | CX z0_3sigma Clean/truth | seen 33.0
\newcommand{\truthGuided}{\TBD}       % HLDOC | CX truth_guided Clean/truth, PU200
\newcommand{\finderCap}{\TBD}         % HLDOC | CX finder cap (0.17 mm), PU200
\newcommand{\oneTrkVtx}{\TBD}         % HLDOC | CX GNN vertices with one track, PU200 | seen 36,932 of 180,460
\newcommand{\cleanTruthDomTwo}{\TBD}  % HLDOC | CX Clean/truth, dominant n_trk >= 2 | seen 63.0
\newcommand{\cleanTruthMinTwo}{\TBD}  % HLDOC | CX Clean/truth, >= 2 tracks | seen 53.9
\newcommand{\sigmaPVmed}{\TBD}        % HLDOC | CX median peak width at PU200 [mm] | seen 0.072
\newcommand{\hsRecoDensMin}{\TBD}     % HLDOC | SUP hs_vs_density: lowest HS reconstruction in any bin [%] | seen 99.3
\newcommand{\hsSelDensLo}{\TBD}       % HLDOC | SUP hs_vs_density: smallest selection lead [pt] | seen +2.7
\newcommand{\hsSelDensHi}{\TBD}       % HLDOC | SUP largest selection lead [pt] | seen +5.0
\newcommand{\hsResDensLo}{\TBD}       % HLDOC | SUP smallest resolution gain [um] | seen 0.2
\newcommand{\hsResDensHi}{\TBD}       % HLDOC | SUP largest resolution gain [um] | seen 0.9

%======================================================================
% 5. CPU timing
%======================================================================
% quoted from the note (sections/04b_timing.tex, tab:athena-timing-breakdown): final
\newcommand{\athCPUms}{2401}
\newcommand{\athGPUms}{76}
\newcommand{\athCPUinf}{2189}
\newcommand{\athOther}{31}
\newcommand{\athGNN}{182}
\newcommand{\athAMVF}{650}
\newcommand{\athAMVFlo}{600}
\newcommand{\athAMVFhi}{750}
% 9 Nov measurement (docs/evaluation/cpu_small_model_2026-11-09.md): final
\newcommand{\spdSEightNov}{9.0}
% frozen-chain timing (HL-LHC agent): TBD
\newcommand{\cpuSEightPack}{\TBD}  % HLDOC | CPU density s8 packed ms/event | seen 174
\newcommand{\cpuVsixPad}{\TBD}     % HLDOC | CPU density v6 padded ms/event | seen 1892
\newcommand{\spdPack}{\TBD}        % HLDOC | CPU s8 packed / v6 padded speed-up | seen 10.9
\newcommand{\spdPad}{\TBD}         % HLDOC | CPU s8 padded speed-up | seen 6.2
\newcommand{\cpuGNN}{\TBD}         % HLDOC | CPU GNN v2 forward ms/event on sneezy | seen 97.9
\newcommand{\projChain}{\TBD}      % HLDOC | CPU projected Athena chain, s8 packed [ms] | seen 413
\newcommand{\projChainPad}{\TBD}   % HLDOC | CPU projected chain, s8 padded [ms] | seen 566
\newcommand{\projWorst}{\TBD}      % HLDOC | CPU worst case [ms] | seen 531
\newcommand{\projFactor}{\TBD}     % HLDOC | CPU AMVF / chain at 650 ms | seen 1.57

%======================================================================
% 6. Decision studies (committed docs, final)
%======================================================================
% s8 vs v6, docs/evaluation/model_freeze_candidates_2026-11-12.md
\newcommand{\sEightDEffA}{$-0.15$}
\newcommand{\sEightDEffB}{$-0.08$}
\newcommand{\sEightDSig}{$+1.2\pm1.1$}
\newcommand{\sEightDpse}{$+0.6$}
\newcommand{\sFiveDSig}{$+9$ to $+12$}
% after the GNN, docs/evaluation/gnn_v2_freeze_2026-11-12.md Sec. 7 (t = 0.97, inclusive)
\newcommand{\sEightDClean}{$-0.14\pm0.02$}
\newcommand{\sEightDHSsel}{$-0.10\pm0.06$}
\newcommand{\sEightDHSres}{$+0.04\pm0.04$}
% GNN v2 checkpoint and t options (v6 peaks, test r16443), gnn_v2_freeze_2026-11-12.md
\newcommand{\gnnValLoss}{0.255807}
\newcommand{\optAclean}{73.83}
\newcommand{\optDclean}{73.23}
\newcommand{\optBclean}{68.70}
\newcommand{\optCclean}{65.84}
\newcommand{\optEclean}{56.79}
\newcommand{\optAcleanPU}{67.73}
\newcommand{\optDcleanPU}{66.99}
\newcommand{\optAcleanLow}{84.22}
\newcommand{\optDcleanLow}{86.45}
\newcommand{\optAFone}{72.09}
\newcommand{\optDFone}{74.45}
\newcommand{\optAMVFFone}{74.43}
\newcommand{\optDHSsel}{97.04}
\newcommand{\optAHSsel}{96.79}
\newcommand{\optDHSres}{9.53}
\newcommand{\optAHSres}{9.70}
% HS selection rule, docs/evaluation/hs_selection_identical_rule_2026-11-11.md
\newcommand{\chiCut}{18.42}
\newcommand{\ruleAgreeAll}{96.2}
\newcommand{\ruleAgreePU}{94.8}
\newcommand{\plainAgreeAll}{86.8}
% HS resolution study (GNN v4, v6 peaks), docs/evaluation/itk_hs_resolution_2026-11-09.md
\newcommand{\hsResPeak}{12.38}
\newcommand{\hsResTruthTrk}{9.09}
\newcommand{\hsResAMVFlist}{9.50}
\newcommand{\hsResNoGNN}{9.28}
\newcommand{\hsResOneD}{11.99}
% target readout, docs/evaluation/run2_vs_run4_resolution_2026-11-09.md (2026-11-11)
\newcommand{\tgtSigRunTwo}{0.203}
\newcommand{\pvfSigRunTwo}{0.410}
\newcommand{\amvfSigRunTwo}{0.789}
\newcommand{\tgtSigHL}{0.145}
\newcommand{\tgtSigHLwone}{0.137}
\newcommand{\tgtSigHLfine}{0.124}
\newcommand{\pvfSigHLvsix}{0.218}
\newcommand{\pvfSigHLrunTwoCfg}{0.241}
\newcommand{\tgtHSum}{0.77}
\newcommand{\pvfHSum}{12.4}
\newcommand{\idealPURunTwo}{75.6}
\newcommand{\idealPUHL}{44.4}
\newcommand{\idealHSHL}{6.0}

%======================================================================
% 7. Low pile-up (LPDOC; Run 2 MC 2,550 events, data 2,500 each)
%======================================================================
\newcommand{\lpEffPVF}{\TBD}       % LPDOC R1 | VF mc/pvf/recovery | seen 81.55
\newcommand{\lpEffPVFErr}{\TBD}    % LPDOC R1 | VF mc/pvf/recovery_err | seen 0.14
\newcommand{\lpEffAMVF}{\TBD}      % LPDOC R1 | VF mc/amvf/recovery | seen 75.74
\newcommand{\lpEffDiff}{\TBD}      % LPDOC R1 | VF mc/delta_pvf_minus_amvf/recovery/delta | seen +5.81
\newcommand{\lpFakeFracPVF}{\TBD}  % LPDOC R1 | VF mc/pvf/fake_fraction/value | seen 9.86
\newcommand{\lpFakeFracAMVF}{\TBD} % LPDOC R1 | VF mc/amvf/fake_fraction/value | seen 11.06
\newcommand{\lpFakeEvtPVF}{\TBD}   % LPDOC R1 | VF mc/pvf/no_ref_per_evt | seen 2.363
\newcommand{\lpFakeEvtAMVF}{\TBD}  % LPDOC R1 | VF mc/amvf/no_ref_per_evt | seen 2.487
\newcommand{\lpSigPVF}{\TBD}       % LPDOC R1 | VF mc/pvf/sigma_vtx_vtx_mm/value | seen 0.392
\newcommand{\lpSigPVFErr}{\TBD}    % LPDOC R1 | VF mc/pvf/sigma_vtx_vtx_mm/err | seen 0.009
\newcommand{\lpSigAMVF}{\TBD}      % LPDOC R1 | VF mc/amvf/sigma_vtx_vtx_mm/value | seen 0.789
\newcommand{\lpSigAMVFErr}{\TBD}   % LPDOC R1 | VF mc/amvf/sigma_vtx_vtx_mm/err | seen 0.015
\newcommand{\lpPsePVF}{\TBD}       % LPDOC R1 | VF mc/pvf/p68_residual_um | seen 72.5
\newcommand{\lpPseAMVF}{\TBD}      % LPDOC R1 | VF mc/amvf/p68_residual_um | seen 66.8
\newcommand{\lpWinRatio}{\TBD}     % LPDOC R1 | VF mc/pvf/window_over_p68 | seen 2.35
\newcommand{\lpHSPVF}{\TBD}        % LPDOC R1 | VF mc/pvf/hs_eff/eff | seen 98.94
\newcommand{\lpHSAMVF}{\TBD}       % LPDOC R1 | VF mc/amvf/hs_eff/eff | seen 99.29
\newcommand{\lpRecRunTwo}{\TBD}    % LPDOC R3 | VF run2/pvf/recovery | seen 94.39
\newcommand{\lpRecRunThree}{\TBD}  % LPDOC R3 | VF run3/pvf/recovery | seen 94.03
\newcommand{\lpNoRefRunTwo}{\TBD}  % LPDOC R3 | VF run2/pvf/no_ref_per_evt | seen 3.006
\newcommand{\lpNoRefRunThree}{\TBD}% LPDOC R3 | VF run3/pvf/no_ref_per_evt | seen 3.192
\newcommand{\lpRecMCref}{\TBD}     % LPDOC R3 | VF mc/pvf_vs_amvf_reference/recovery | seen 93.09
\newcommand{\lpSigRunTwo}{\TBD}    % LPDOC R3 | VF run2/pvf/sigma_vtx_vtx_mm/value | seen 0.356
\newcommand{\lpSigRunThree}{\TBD}  % LPDOC R3 | VF run3/pvf/sigma_vtx_vtx_mm/value | seen 0.333
\newcommand{\lpSigAMVFRunTwo}{\TBD}  % LPDOC R3 | VF run2/amvf_ref/sigma_vtx_vtx_mm/value | seen 0.808
\newcommand{\lpSigAMVFRunThree}{\TBD}% LPDOC R3 | VF run3/amvf_ref/sigma_vtx_vtx_mm/value | seen 0.813
\newcommand{\lpCleanTruth}{\TBD}   % LPDOC R2 | GM t0.97/clean_per_truth/value | seen 72.84
\newcommand{\lpCleanTruthAMVF}{\TBD}% LPDOC R2 | GM amvf/clean_per_truth/value | seen 60.91
\newcommand{\lpClean}{\TBD}        % LPDOC R2 | GM t0.97/clean_per_reco/value | seen 86.90
\newcommand{\lpMerged}{\TBD}       % LPDOC R2 | GM t0.97/merged_per_reco/value | seen 11.05
\newcommand{\lpSplit}{\TBD}        % LPDOC R2 | GM t0.97/split_per_reco/value | seen 2.04
\newcommand{\lpFake}{\TBD}         % LPDOC R2 | GM t0.97/fake_per_reco/value | seen 0.003
\newcommand{\lpCleanAMVF}{\TBD}    % LPDOC R2 | GM amvf/clean_per_reco/value | seen 77.23
\newcommand{\lpMergedAMVF}{\TBD}   % LPDOC R2 | GM amvf/merged_per_reco/value | seen 20.19
\newcommand{\lpSplitAMVF}{\TBD}    % LPDOC R2 | GM amvf/split_per_reco/value | seen 2.53
\newcommand{\lpFakeAMVF}{\TBD}     % LPDOC R2 | GM amvf/fake_per_reco/value | seen 0.054
\newcommand{\lpFOne}{\TBD}         % LPDOC R2 | GM t0.97/trk_f1/value | seen 86.95
\newcommand{\lpFOneAMVF}{\TBD}     % LPDOC R2 | GM amvf/trk_f1/value | seen 84.89
\newcommand{\lpHSID}{\TBD}         % LPDOC R2 | GM t0.97/hs_id/value | seen 99.06
\newcommand{\lpHSIDAMVF}{\TBD}     % LPDOC R2 | GM amvf/hs_id/value | seen 98.75
\newcommand{\lpCleanTruthNN}{\TBD} % LPDOC Sec. 3 | t = 0.99 Clean/truth | seen 74.35
\newcommand{\lpEmpty}{\TBD}        % LPDOC R2 | GM t0.97/totals/n_empty | seen 593
\newcommand{\lpRthreeAssigned}{\TBD}% LPDOC R4 | R3A gnn_assigned_fraction | seen 92.38
\newcommand{\lpRthreeAgree}{\TBD}  % LPDOC R4 | R3A agreement/1.0/vs_both_assigned | seen 96.03
\newcommand{\lpRthreeAgreeW}{\TBD} % LPDOC R4 | R3A agreement/0.17/vs_both_assigned | seen 95.09
\newcommand{\lpRthreeLead}{\TBD}   % LPDOC R4 | R3A leading_vertex_agreement | seen 91.20
\newcommand{\lpCentroidShift}{\TBD}% LPDOC Sec. 4 | L/summary/diag_centroid_vs_mean.json eff difference [pt] | seen -0.04

% low-mu behaviour of GNN v2 (gnn_v2_freeze_2026-11-12.md Sec. 6, v6 peaks, validation): final
\newcommand{\lmTopOneLow}{97}
\newcommand{\lmTopOneHigh}{71}
\newcommand{\lmBelowNN}{24--34}
\newcommand{\lmSmallClean}{92}
\newcommand{\lmSmallCleanNN}{61}
\newcommand{\lmSmallBelowLow}{47}
\newcommand{\lmSmallBelowMid}{25}
\newcommand{\lmSmallBelowHigh}{41}
\newcommand{\lmTrackShare}{2}
\newcommand{\lmVoneBelow}{47}
\newcommand{\lmVtwoBelow}{33}
\newcommand{\lmSplitHS}{60}
\newcommand{\lmSplitOne}{63}
\newcommand{\lmSplitDist}{0.8}
\newcommand{\lmFakeMax}{0.4}
\newcommand{\lmCleanPeakT}{\text{0.90--0.92}}
% chain_extras finding (HL-LHC agent): GNN vs AMVF track F1 at mu 0-10
\newcommand{\lmFOneGNN}{\TBD}  % HLDOC | CX trk_f1 vs mu, bin 0-10, gnn_t | seen 87.3
\newcommand{\lmFOneAMVF}{\TBD} % HLDOC | CX trk_f1 vs mu, bin 0-10, amvf | seen 96.6
% window audit (matching_window_audit_2026-11-09.md): final
\newcommand{\greedyBias}{10}
% resolution vs mu study (v6, resolution_vs_mu_2026-11-11.md): final
\newcommand{\rvmChiLo}{4}
\newcommand{\rvmChiHi}{57}
\newcommand{\rvmExcessPVF}{1.61}
\newcommand{\rvmExcessAMVF}{1.60}
\newcommand{\rvmExcessTruth}{1.01}
\newcommand{\optAMVFcleanLow}{77.06}
\newcommand{\optAMVFHSres}{9.93}
\newcommand{\optDvsAclean}{0.6}
\newcommand{\optDvsAlow}{2.2}
\newcommand{\optDvsAFone}{2.4}
\newcommand{\optDvsAres}{0.17}
\newcommand{\hsResAMVFstudy}{10.25}

\imgsrc{mu_training}{outputs/\allowbreak{}11\_\allowbreak{}11\_\allowbreak{}2026\_\allowbreak{}output/\allowbreak{}datasets/\allowbreak{}mu\_\allowbreak{}training.pdf}{pdf}
\imgsrc{mu_validation}{outputs/\allowbreak{}11\_\allowbreak{}11\_\allowbreak{}2026\_\allowbreak{}output/\allowbreak{}datasets/\allowbreak{}mu\_\allowbreak{}validation.pdf}{pdf}
\imgsrc{mu_test}{outputs/\allowbreak{}11\_\allowbreak{}11\_\allowbreak{}2026\_\allowbreak{}output/\allowbreak{}datasets/\allowbreak{}mu\_\allowbreak{}test.pdf}{pdf}
\imgsrc{mu_r16443}{outputs/\allowbreak{}11\_\allowbreak{}11\_\allowbreak{}2026\_\allowbreak{}output/\allowbreak{}datasets/\allowbreak{}mu\_\allowbreak{}held-out\_\allowbreak{}601229\_\allowbreak{}r16443.pdf}{pdf}
\imgsrc{mu_r16638}{outputs/\allowbreak{}11\_\allowbreak{}11\_\allowbreak{}2026\_\allowbreak{}output/\allowbreak{}datasets/\allowbreak{}mu\_\allowbreak{}held-out\_\allowbreak{}601229\_\allowbreak{}r16638.pdf}{pdf}
\imgsrc{lowpu_mu}{outputs/\allowbreak{}11\_\allowbreak{}12\_\allowbreak{}2026\_\allowbreak{}output/\allowbreak{}paper\_\allowbreak{}deck/\allowbreak{}lowpu\_\allowbreak{}mu/\allowbreak{}lowpu\_\allowbreak{}mu\_\allowbreak{}distributions.pdf}{pdf}
\imgsrc{window_decision}{presentations/\allowbreak{}mattia/\allowbreak{}11\_\allowbreak{}13\_\allowbreak{}2026\_\allowbreak{}EB/\allowbreak{}images/\allowbreak{}window\_\allowbreak{}residuals.png}{png}
\imgsrc{eff_vs_mu}{outputs/\allowbreak{}11\_\allowbreak{}12\_\allowbreak{}2026\_\allowbreak{}output/\allowbreak{}hllhc\_\allowbreak{}final/\allowbreak{}chain\_\allowbreak{}extras/\allowbreak{}plots/\allowbreak{}eff\_\allowbreak{}pregnn\_\allowbreak{}vs\_\allowbreak{}mu.pdf}{pdf}
\imgsrc{fake_fraction_vs_mu}{outputs/\allowbreak{}11\_\allowbreak{}12\_\allowbreak{}2026\_\allowbreak{}output/\allowbreak{}hllhc\_\allowbreak{}final/\allowbreak{}h1\_\allowbreak{}vertex\_\allowbreak{}finding/\allowbreak{}plots/\allowbreak{}fake\_\allowbreak{}fraction\_\allowbreak{}vs\_\allowbreak{}mu\_\allowbreak{}r16443}{}
\imgsrc{fakes_per_event_vs_mu}{outputs/\allowbreak{}11\_\allowbreak{}12\_\allowbreak{}2026\_\allowbreak{}output/\allowbreak{}hllhc\_\allowbreak{}final/\allowbreak{}h1\_\allowbreak{}vertex\_\allowbreak{}finding/\allowbreak{}plots/\allowbreak{}fakes\_\allowbreak{}per\_\allowbreak{}event\_\allowbreak{}vs\_\allowbreak{}mu\_\allowbreak{}r16443}{}
\imgsrc{eff_one_to_one_vs_mu}{outputs/\allowbreak{}11\_\allowbreak{}12\_\allowbreak{}2026\_\allowbreak{}output/\allowbreak{}hllhc\_\allowbreak{}final/\allowbreak{}h1\_\allowbreak{}vertex\_\allowbreak{}finding/\allowbreak{}plots/\allowbreak{}eff\_\allowbreak{}one\_\allowbreak{}to\_\allowbreak{}one\_\allowbreak{}vs\_\allowbreak{}mu\_\allowbreak{}r16443}{}
\imgsrc{splits_per_event_vs_mu}{outputs/\allowbreak{}11\_\allowbreak{}12\_\allowbreak{}2026\_\allowbreak{}output/\allowbreak{}hllhc\_\allowbreak{}final/\allowbreak{}h1\_\allowbreak{}vertex\_\allowbreak{}finding/\allowbreak{}plots/\allowbreak{}splits\_\allowbreak{}per\_\allowbreak{}event\_\allowbreak{}vs\_\allowbreak{}mu\_\allowbreak{}r16443}{}
\imgsrc{yield_vs_mu}{outputs/\allowbreak{}11\_\allowbreak{}12\_\allowbreak{}2026\_\allowbreak{}output/\allowbreak{}hllhc\_\allowbreak{}final/\allowbreak{}h1\_\allowbreak{}vertex\_\allowbreak{}finding/\allowbreak{}plots/\allowbreak{}fig6\_\allowbreak{}vertices\_\allowbreak{}vs\_\allowbreak{}mu\_\allowbreak{}r16443}{}
\imgsrc{amvf_pos_categories_vs_mu}{outputs/\allowbreak{}11\_\allowbreak{}12\_\allowbreak{}2026\_\allowbreak{}output/\allowbreak{}hllhc\_\allowbreak{}final/\allowbreak{}h1\_\allowbreak{}vertex\_\allowbreak{}finding/\allowbreak{}plots/\allowbreak{}fig6\_\allowbreak{}categories\_\allowbreak{}vs\_\allowbreak{}mu\_\allowbreak{}amvf\_\allowbreak{}r16443}{}
\imgsrc{deltaz_pu200}{outputs/\allowbreak{}11\_\allowbreak{}12\_\allowbreak{}2026\_\allowbreak{}output/\allowbreak{}hllhc\_\allowbreak{}final/\allowbreak{}h1\_\allowbreak{}vertex\_\allowbreak{}finding/\allowbreak{}plots/\allowbreak{}paper\_\allowbreak{}figures\_\allowbreak{}r16443/\allowbreak{}fig4\_\allowbreak{}hllhc\_\allowbreak{}deltaz}{}
\imgsrc{pos_categories_pu200}{outputs/\allowbreak{}11\_\allowbreak{}12\_\allowbreak{}2026\_\allowbreak{}output/\allowbreak{}hllhc\_\allowbreak{}final/\allowbreak{}h1\_\allowbreak{}vertex\_\allowbreak{}finding/\allowbreak{}plots/\allowbreak{}paper\_\allowbreak{}figures\_\allowbreak{}r16443/\allowbreak{}fig3a\_\allowbreak{}hllhc\_\allowbreak{}categories\_\allowbreak{}pu200}{}
\imgsrc{residuals_pu200}{outputs/\allowbreak{}11\_\allowbreak{}12\_\allowbreak{}2026\_\allowbreak{}output/\allowbreak{}hllhc\_\allowbreak{}final/\allowbreak{}h1\_\allowbreak{}vertex\_\allowbreak{}finding/\allowbreak{}plots/\allowbreak{}fig7b\_\allowbreak{}residuals\_\allowbreak{}r16443}{}
\imgsrc{window_justification}{outputs/\allowbreak{}11\_\allowbreak{}12\_\allowbreak{}2026\_\allowbreak{}output/\allowbreak{}hllhc\_\allowbreak{}final/\allowbreak{}h1\_\allowbreak{}vertex\_\allowbreak{}finding/\allowbreak{}plots/\allowbreak{}window\_\allowbreak{}justification\_\allowbreak{}r16443}{}
\imgsrc{window_scan}{outputs/\allowbreak{}11\_\allowbreak{}12\_\allowbreak{}2026\_\allowbreak{}output/\allowbreak{}hllhc\_\allowbreak{}final/\allowbreak{}h1\_\allowbreak{}vertex\_\allowbreak{}finding/\allowbreak{}plots/\allowbreak{}window\_\allowbreak{}scan\_\allowbreak{}r16443}{}
\imgsrc{p68_vs_mu}{outputs/\allowbreak{}11\_\allowbreak{}12\_\allowbreak{}2026\_\allowbreak{}output/\allowbreak{}hllhc\_\allowbreak{}final/\allowbreak{}resolution\_\allowbreak{}vs\_\allowbreak{}mu/\allowbreak{}resolution\_\allowbreak{}vs\_\allowbreak{}mu\_\allowbreak{}p68}{}
\imgsrc{resolution_vs_mu}{outputs/\allowbreak{}11\_\allowbreak{}12\_\allowbreak{}2026\_\allowbreak{}output/\allowbreak{}hllhc\_\allowbreak{}final/\allowbreak{}resolution\_\allowbreak{}vs\_\allowbreak{}mu/\allowbreak{}resolution\_\allowbreak{}vs\_\allowbreak{}mu}{}
\imgsrc{hs_reco_vs_mu}{outputs/\allowbreak{}11\_\allowbreak{}12\_\allowbreak{}2026\_\allowbreak{}output/\allowbreak{}hllhc\_\allowbreak{}final/\allowbreak{}chain/\allowbreak{}plots/\allowbreak{}hs\_\allowbreak{}reco\_\allowbreak{}eff\_\allowbreak{}vs\_\allowbreak{}mu.pdf}{pdf}
\imgsrc{eff_ntrk_pre_pu200}{outputs/\allowbreak{}11\_\allowbreak{}12\_\allowbreak{}2026\_\allowbreak{}output/\allowbreak{}hllhc\_\allowbreak{}final/\allowbreak{}chain\_\allowbreak{}extras/\allowbreak{}plots/\allowbreak{}eff\_\allowbreak{}vs\_\allowbreak{}n\_\allowbreak{}trk\_\allowbreak{}pregnn\_\allowbreak{}pu200.pdf}{pdf}
\imgsrc{eff_density_pre_pu200}{outputs/\allowbreak{}11\_\allowbreak{}12\_\allowbreak{}2026\_\allowbreak{}output/\allowbreak{}hllhc\_\allowbreak{}final/\allowbreak{}chain\_\allowbreak{}extras/\allowbreak{}plots/\allowbreak{}eff\_\allowbreak{}vs\_\allowbreak{}density\_\allowbreak{}pregnn\_\allowbreak{}pu200.pdf}{pdf}
\imgsrc{eff_dz_pre_pu200}{outputs/\allowbreak{}11\_\allowbreak{}12\_\allowbreak{}2026\_\allowbreak{}output/\allowbreak{}hllhc\_\allowbreak{}final/\allowbreak{}chain\_\allowbreak{}extras/\allowbreak{}plots/\allowbreak{}eff\_\allowbreak{}vs\_\allowbreak{}dz\_\allowbreak{}nearest\_\allowbreak{}pregnn\_\allowbreak{}pu200.pdf}{pdf}
\imgsrc{classes_pu200}{outputs/\allowbreak{}11\_\allowbreak{}12\_\allowbreak{}2026\_\allowbreak{}output/\allowbreak{}hllhc\_\allowbreak{}final/\allowbreak{}supplementary/\allowbreak{}plots/\allowbreak{}main/\allowbreak{}classes\_\allowbreak{}per\_\allowbreak{}event\_\allowbreak{}pu200\_\allowbreak{}primary.pdf}{pdf}
\imgsrc{clean_evt_vs_mu}{outputs/\allowbreak{}11\_\allowbreak{}12\_\allowbreak{}2026\_\allowbreak{}output/\allowbreak{}hllhc\_\allowbreak{}final/\allowbreak{}supplementary/\allowbreak{}plots/\allowbreak{}main/\allowbreak{}clean\_\allowbreak{}per\_\allowbreak{}event\_\allowbreak{}vs\_\allowbreak{}mu\_\allowbreak{}primary.pdf}{pdf}
\imgsrc{merged_evt_vs_mu}{outputs/\allowbreak{}11\_\allowbreak{}12\_\allowbreak{}2026\_\allowbreak{}output/\allowbreak{}hllhc\_\allowbreak{}final/\allowbreak{}supplementary/\allowbreak{}plots/\allowbreak{}main/\allowbreak{}merged\_\allowbreak{}per\_\allowbreak{}event\_\allowbreak{}vs\_\allowbreak{}mu\_\allowbreak{}primary.pdf}{pdf}
\imgsrc{split_evt_vs_mu}{outputs/\allowbreak{}11\_\allowbreak{}12\_\allowbreak{}2026\_\allowbreak{}output/\allowbreak{}hllhc\_\allowbreak{}final/\allowbreak{}supplementary/\allowbreak{}plots/\allowbreak{}main/\allowbreak{}split\_\allowbreak{}per\_\allowbreak{}event\_\allowbreak{}vs\_\allowbreak{}mu\_\allowbreak{}primary.pdf}{pdf}
\imgsrc{fake_evt_vs_mu}{outputs/\allowbreak{}11\_\allowbreak{}12\_\allowbreak{}2026\_\allowbreak{}output/\allowbreak{}hllhc\_\allowbreak{}final/\allowbreak{}supplementary/\allowbreak{}plots/\allowbreak{}main/\allowbreak{}fake\_\allowbreak{}per\_\allowbreak{}event\_\allowbreak{}vs\_\allowbreak{}mu\_\allowbreak{}primary.pdf}{pdf}
\imgsrc{clean_truth_vs_mu}{outputs/\allowbreak{}11\_\allowbreak{}12\_\allowbreak{}2026\_\allowbreak{}output/\allowbreak{}hllhc\_\allowbreak{}final/\allowbreak{}chain/\allowbreak{}plots/\allowbreak{}categories\_\allowbreak{}clean\_\allowbreak{}per\_\allowbreak{}truth\_\allowbreak{}vs\_\allowbreak{}mu.pdf}{pdf}
\imgsrc{hs_sel_vs_mu}{outputs/\allowbreak{}11\_\allowbreak{}12\_\allowbreak{}2026\_\allowbreak{}output/\allowbreak{}hllhc\_\allowbreak{}final/\allowbreak{}chain/\allowbreak{}plots/\allowbreak{}hs\_\allowbreak{}selection\_\allowbreak{}vs\_\allowbreak{}mu.pdf}{pdf}
\imgsrc{hs_res_vs_mu}{outputs/\allowbreak{}11\_\allowbreak{}12\_\allowbreak{}2026\_\allowbreak{}output/\allowbreak{}hllhc\_\allowbreak{}final/\allowbreak{}chain/\allowbreak{}plots/\allowbreak{}hs\_\allowbreak{}resolution\_\allowbreak{}vs\_\allowbreak{}mu.pdf}{pdf}
\imgsrc{hs_res_vs_ntrk}{outputs/\allowbreak{}11\_\allowbreak{}12\_\allowbreak{}2026\_\allowbreak{}output/\allowbreak{}hllhc\_\allowbreak{}final/\allowbreak{}chain/\allowbreak{}plots/\allowbreak{}hs\_\allowbreak{}resolution\_\allowbreak{}vs\_\allowbreak{}ntrk}{}
\imgsrc{hs_reco_vs_density}{outputs/\allowbreak{}11\_\allowbreak{}12\_\allowbreak{}2026\_\allowbreak{}output/\allowbreak{}hllhc\_\allowbreak{}final/\allowbreak{}supplementary/\allowbreak{}plots/\allowbreak{}main/\allowbreak{}hs\_\allowbreak{}reco\_\allowbreak{}eff\_\allowbreak{}vs\_\allowbreak{}density\_\allowbreak{}itk.pdf}{pdf}
\imgsrc{hs_sel_vs_density}{outputs/\allowbreak{}11\_\allowbreak{}12\_\allowbreak{}2026\_\allowbreak{}output/\allowbreak{}hllhc\_\allowbreak{}final/\allowbreak{}supplementary/\allowbreak{}plots/\allowbreak{}main/\allowbreak{}hs\_\allowbreak{}sel\_\allowbreak{}eff\_\allowbreak{}vs\_\allowbreak{}density\_\allowbreak{}itk.pdf}{pdf}
\imgsrc{hs_res_vs_density}{outputs/\allowbreak{}11\_\allowbreak{}12\_\allowbreak{}2026\_\allowbreak{}output/\allowbreak{}hllhc\_\allowbreak{}final/\allowbreak{}supplementary/\allowbreak{}plots/\allowbreak{}main/\allowbreak{}hs\_\allowbreak{}resolution\_\allowbreak{}vs\_\allowbreak{}density\_\allowbreak{}itk.pdf}{pdf}
\imgsrc{trk_f1_vs_mu}{outputs/\allowbreak{}11\_\allowbreak{}12\_\allowbreak{}2026\_\allowbreak{}output/\allowbreak{}hllhc\_\allowbreak{}final/\allowbreak{}chain/\allowbreak{}plots/\allowbreak{}track\_\allowbreak{}f1\_\allowbreak{}vs\_\allowbreak{}mu.pdf}{pdf}
\imgsrc{trk_eff_vs_mu}{outputs/\allowbreak{}11\_\allowbreak{}12\_\allowbreak{}2026\_\allowbreak{}output/\allowbreak{}hllhc\_\allowbreak{}final/\allowbreak{}chain/\allowbreak{}plots/\allowbreak{}track\_\allowbreak{}efficiency\_\allowbreak{}vs\_\allowbreak{}mu.pdf}{pdf}
\imgsrc{trk_pur_vs_mu}{outputs/\allowbreak{}11\_\allowbreak{}12\_\allowbreak{}2026\_\allowbreak{}output/\allowbreak{}hllhc\_\allowbreak{}final/\allowbreak{}chain/\allowbreak{}plots/\allowbreak{}track\_\allowbreak{}purity\_\allowbreak{}vs\_\allowbreak{}mu.pdf}{pdf}
\imgsrc{eff_ntrk_post_pu200}{outputs/\allowbreak{}11\_\allowbreak{}12\_\allowbreak{}2026\_\allowbreak{}output/\allowbreak{}hllhc\_\allowbreak{}final/\allowbreak{}chain\_\allowbreak{}extras/\allowbreak{}plots/\allowbreak{}eff\_\allowbreak{}vs\_\allowbreak{}n\_\allowbreak{}trk\_\allowbreak{}postgnn\_\allowbreak{}pu200.pdf}{pdf}
\imgsrc{eff_density_post_pu200}{outputs/\allowbreak{}11\_\allowbreak{}12\_\allowbreak{}2026\_\allowbreak{}output/\allowbreak{}hllhc\_\allowbreak{}final/\allowbreak{}chain\_\allowbreak{}extras/\allowbreak{}plots/\allowbreak{}eff\_\allowbreak{}vs\_\allowbreak{}density\_\allowbreak{}postgnn\_\allowbreak{}pu200.pdf}{pdf}
\imgsrc{eff_dz_post_pu200}{outputs/\allowbreak{}11\_\allowbreak{}12\_\allowbreak{}2026\_\allowbreak{}output/\allowbreak{}hllhc\_\allowbreak{}final/\allowbreak{}chain\_\allowbreak{}extras/\allowbreak{}plots/\allowbreak{}eff\_\allowbreak{}vs\_\allowbreak{}dz\_\allowbreak{}nearest\_\allowbreak{}postgnn\_\allowbreak{}pu200.pdf}{pdf}
\imgsrc{hs_clean_vs_mu}{outputs/\allowbreak{}11\_\allowbreak{}12\_\allowbreak{}2026\_\allowbreak{}output/\allowbreak{}hllhc\_\allowbreak{}final/\allowbreak{}chain/\allowbreak{}plots/\allowbreak{}hs\_\allowbreak{}clean\_\allowbreak{}fraction\_\allowbreak{}vs\_\allowbreak{}mu.pdf}{pdf}
\imgsrc{hs_trkfrac_vs_mu}{outputs/\allowbreak{}11\_\allowbreak{}12\_\allowbreak{}2026\_\allowbreak{}output/\allowbreak{}hllhc\_\allowbreak{}final/\allowbreak{}chain/\allowbreak{}plots/\allowbreak{}hs\_\allowbreak{}track\_\allowbreak{}fraction\_\allowbreak{}vs\_\allowbreak{}mu.pdf}{pdf}
\imgsrc{classes_variants}{outputs/\allowbreak{}11\_\allowbreak{}12\_\allowbreak{}2026\_\allowbreak{}output/\allowbreak{}hllhc\_\allowbreak{}final/\allowbreak{}supplementary/\allowbreak{}plots/\allowbreak{}backup/\allowbreak{}classes\_\allowbreak{}per\_\allowbreak{}event\_\allowbreak{}pu200\_\allowbreak{}all\_\allowbreak{}variants.pdf}{pdf}
\imgsrc{amvf_fit_weights}{outputs/\allowbreak{}11\_\allowbreak{}12\_\allowbreak{}2026\_\allowbreak{}output/\allowbreak{}hllhc\_\allowbreak{}final/\allowbreak{}supplementary/\allowbreak{}plots/\allowbreak{}backup/\allowbreak{}amvf\_\allowbreak{}fit\_\allowbreak{}weight\_\allowbreak{}distribution.pdf}{pdf}
\imgsrc{density_definitions}{outputs/\allowbreak{}11\_\allowbreak{}12\_\allowbreak{}2026\_\allowbreak{}output/\allowbreak{}hllhc\_\allowbreak{}final/\allowbreak{}supplementary/\allowbreak{}plots/\allowbreak{}backup/\allowbreak{}local\_\allowbreak{}density\_\allowbreak{}definitions.pdf}{pdf}
\imgsrc{eq4_vs_mu}{outputs/\allowbreak{}11\_\allowbreak{}12\_\allowbreak{}2026\_\allowbreak{}output/\allowbreak{}hllhc\_\allowbreak{}final/\allowbreak{}chain\_\allowbreak{}extras/\allowbreak{}plots/\allowbreak{}eq4\_\allowbreak{}clean\_\allowbreak{}vs\_\allowbreak{}mu.pdf}{pdf}
\imgsrc{z0_baselines_vs_mu}{outputs/\allowbreak{}11\_\allowbreak{}12\_\allowbreak{}2026\_\allowbreak{}output/\allowbreak{}hllhc\_\allowbreak{}final/\allowbreak{}chain\_\allowbreak{}extras/\allowbreak{}plots/\allowbreak{}z0\_\allowbreak{}baselines\_\allowbreak{}clean\_\allowbreak{}vs\_\allowbreak{}mu.pdf}{pdf}
\imgsrc{truth_guided_vs_mu}{outputs/\allowbreak{}11\_\allowbreak{}12\_\allowbreak{}2026\_\allowbreak{}output/\allowbreak{}hllhc\_\allowbreak{}final/\allowbreak{}chain\_\allowbreak{}extras/\allowbreak{}plots/\allowbreak{}truth\_\allowbreak{}guided\_\allowbreak{}clean\_\allowbreak{}vs\_\allowbreak{}mu.pdf}{pdf}
\imgsrc{finder_cap_vs_mu}{outputs/\allowbreak{}11\_\allowbreak{}12\_\allowbreak{}2026\_\allowbreak{}output/\allowbreak{}hllhc\_\allowbreak{}final/\allowbreak{}chain\_\allowbreak{}extras/\allowbreak{}plots/\allowbreak{}finder\_\allowbreak{}cap\_\allowbreak{}vs\_\allowbreak{}mu.pdf}{pdf}
\imgsrc{stages_pu200}{outputs/\allowbreak{}11\_\allowbreak{}12\_\allowbreak{}2026\_\allowbreak{}output/\allowbreak{}hllhc\_\allowbreak{}final/\allowbreak{}chain\_\allowbreak{}extras/\allowbreak{}plots/\allowbreak{}clean\_\allowbreak{}yield\_\allowbreak{}stages\_\allowbreak{}pu200.pdf}{pdf}
\imgsrc{wrong_rate_vs_mu}{outputs/\allowbreak{}11\_\allowbreak{}12\_\allowbreak{}2026\_\allowbreak{}output/\allowbreak{}hllhc\_\allowbreak{}final/\allowbreak{}chain\_\allowbreak{}extras/\allowbreak{}plots/\allowbreak{}wrong\_\allowbreak{}rate\_\allowbreak{}vs\_\allowbreak{}mu.pdf}{pdf}
\imgsrc{lost_unassigned_vs_mu}{outputs/\allowbreak{}11\_\allowbreak{}12\_\allowbreak{}2026\_\allowbreak{}output/\allowbreak{}hllhc\_\allowbreak{}final/\allowbreak{}chain\_\allowbreak{}extras/\allowbreak{}plots/\allowbreak{}lost\_\allowbreak{}unassigned\_\allowbreak{}frac\_\allowbreak{}vs\_\allowbreak{}mu.pdf}{pdf}
\imgsrc{unlabelled_unassigned_vs_mu}{outputs/\allowbreak{}11\_\allowbreak{}12\_\allowbreak{}2026\_\allowbreak{}output/\allowbreak{}hllhc\_\allowbreak{}final/\allowbreak{}chain\_\allowbreak{}extras/\allowbreak{}plots/\allowbreak{}unlabelled\_\allowbreak{}left\_\allowbreak{}unassigned\_\allowbreak{}vs\_\allowbreak{}mu.pdf}{pdf}
\imgsrc{sigma_pv_pu200}{outputs/\allowbreak{}11\_\allowbreak{}12\_\allowbreak{}2026\_\allowbreak{}output/\allowbreak{}hllhc\_\allowbreak{}final/\allowbreak{}chain\_\allowbreak{}extras/\allowbreak{}plots/\allowbreak{}sigma\_\allowbreak{}pv\_\allowbreak{}peakwidth\_\allowbreak{}vs\_\allowbreak{}param\_\allowbreak{}pu200.pdf}{pdf}
\imgsrc{t_scan_val_s8}{outputs/\allowbreak{}11\_\allowbreak{}12\_\allowbreak{}2026\_\allowbreak{}output/\allowbreak{}gnn\_\allowbreak{}v2\_\allowbreak{}freeze/\allowbreak{}figures/\allowbreak{}t\_\allowbreak{}scan\_\allowbreak{}val\_\allowbreak{}s8\_\allowbreak{}v2e527.pdf}{pdf}
\imgsrc{vs_mu_test_s8}{outputs/\allowbreak{}11\_\allowbreak{}12\_\allowbreak{}2026\_\allowbreak{}output/\allowbreak{}gnn\_\allowbreak{}v2\_\allowbreak{}freeze/\allowbreak{}figures/\allowbreak{}vs\_\allowbreak{}mu\_\allowbreak{}test\_\allowbreak{}s8\_\allowbreak{}v2e527.pdf}{pdf}
\imgsrc{v6_vs_s8_gnn}{outputs/\allowbreak{}11\_\allowbreak{}12\_\allowbreak{}2026\_\allowbreak{}output/\allowbreak{}gnn\_\allowbreak{}v2\_\allowbreak{}freeze/\allowbreak{}figures/\allowbreak{}v6\_\allowbreak{}vs\_\allowbreak{}s8\_\allowbreak{}after\_\allowbreak{}gnn\_\allowbreak{}vs\_\allowbreak{}mu.pdf}{pdf}
\imgsrc{lowmu_displays}{outputs/\allowbreak{}11\_\allowbreak{}12\_\allowbreak{}2026\_\allowbreak{}output/\allowbreak{}gnn\_\allowbreak{}v2\_\allowbreak{}freeze/\allowbreak{}lowmu/\allowbreak{}lowmu\_\allowbreak{}displays\_\allowbreak{}val\_\allowbreak{}v6\_\allowbreak{}v2e527.pdf}{pdf}
\imgsrc{gnn_learning_curve}{outputs/\allowbreak{}11\_\allowbreak{}12\_\allowbreak{}2026\_\allowbreak{}output/\allowbreak{}gnn\_\allowbreak{}v2\_\allowbreak{}freeze/\allowbreak{}checkpoint/\allowbreak{}learning\_\allowbreak{}curve\_\allowbreak{}v1\_\allowbreak{}v2.pdf}{pdf}
\imgsrc{lp_fig8a}{outputs/\allowbreak{}11\_\allowbreak{}12\_\allowbreak{}2026\_\allowbreak{}output/\allowbreak{}lowpu\_\allowbreak{}final/\allowbreak{}summary/\allowbreak{}figures/\allowbreak{}fig8a\_\allowbreak{}run2mc\_\allowbreak{}deltaz.pdf}{pdf}
\imgsrc{lp_fig8b}{outputs/\allowbreak{}11\_\allowbreak{}12\_\allowbreak{}2026\_\allowbreak{}output/\allowbreak{}lowpu\_\allowbreak{}final/\allowbreak{}summary/\allowbreak{}figures/\allowbreak{}fig8b\_\allowbreak{}run2mc\_\allowbreak{}yield\_\allowbreak{}vs\_\allowbreak{}mu.pdf}{pdf}
\imgsrc{lp_fig9a}{outputs/\allowbreak{}11\_\allowbreak{}12\_\allowbreak{}2026\_\allowbreak{}output/\allowbreak{}lowpu\_\allowbreak{}final/\allowbreak{}gnn\_\allowbreak{}mc/\allowbreak{}figures/\allowbreak{}fig9a\_\allowbreak{}run2mc\_\allowbreak{}eff\_\allowbreak{}vs\_\allowbreak{}ntrk.pdf}{pdf}
\imgsrc{lp_fig9b}{outputs/\allowbreak{}11\_\allowbreak{}12\_\allowbreak{}2026\_\allowbreak{}output/\allowbreak{}lowpu\_\allowbreak{}final/\allowbreak{}gnn\_\allowbreak{}mc/\allowbreak{}figures/\allowbreak{}fig9b\_\allowbreak{}run2mc\_\allowbreak{}categories.pdf}{pdf}
\imgsrc{lp_fig10a}{outputs/\allowbreak{}11\_\allowbreak{}12\_\allowbreak{}2026\_\allowbreak{}output/\allowbreak{}lowpu\_\allowbreak{}final/\allowbreak{}summary/\allowbreak{}figures/\allowbreak{}fig10a\_\allowbreak{}run2\_\allowbreak{}deltaz.pdf}{pdf}
\imgsrc{lp_fig10b}{outputs/\allowbreak{}11\_\allowbreak{}12\_\allowbreak{}2026\_\allowbreak{}output/\allowbreak{}lowpu\_\allowbreak{}final/\allowbreak{}summary/\allowbreak{}figures/\allowbreak{}fig10b\_\allowbreak{}run2\_\allowbreak{}yield\_\allowbreak{}vs\_\allowbreak{}mu.pdf}{pdf}
\imgsrc{lp_fig10c}{outputs/\allowbreak{}11\_\allowbreak{}12\_\allowbreak{}2026\_\allowbreak{}output/\allowbreak{}lowpu\_\allowbreak{}final/\allowbreak{}summary/\allowbreak{}figures/\allowbreak{}fig10c\_\allowbreak{}run3\_\allowbreak{}deltaz.pdf}{pdf}
\imgsrc{lp_fig10d}{outputs/\allowbreak{}11\_\allowbreak{}12\_\allowbreak{}2026\_\allowbreak{}output/\allowbreak{}lowpu\_\allowbreak{}final/\allowbreak{}summary/\allowbreak{}figures/\allowbreak{}fig10d\_\allowbreak{}run3\_\allowbreak{}yield\_\allowbreak{}vs\_\allowbreak{}mu.pdf}{pdf}
\imgsrc{lp_t_scan}{outputs/\allowbreak{}11\_\allowbreak{}12\_\allowbreak{}2026\_\allowbreak{}output/\allowbreak{}lowpu\_\allowbreak{}final/\allowbreak{}gnn\_\allowbreak{}mc/\allowbreak{}figures/\allowbreak{}gnn\_\allowbreak{}t\_\allowbreak{}scan\_\allowbreak{}run2mc.pdf}{pdf}
\imgsrc{athena_timing}{presentations/\allowbreak{}mattia/\allowbreak{}11\_\allowbreak{}13\_\allowbreak{}2026\_\allowbreak{}EB/\allowbreak{}images/\allowbreak{}athena\_\allowbreak{}cpu\_\allowbreak{}gpu\_\allowbreak{}amvf.png}{png}

\hypersetup{pdftitle={PV-Finder, paper deck, November 13, 2026}}

\newcommand{\twofig}[4]{% left name, left caption, right name, right caption
\begin{columns}[T]
\begin{column}{.49\linewidth}\centering\plot{#1}{.58}\fcap{#2}\end{column}
\begin{column}{.49\linewidth}\centering\plot{#3}{.58}\fcap{#4}\end{column}
\end{columns}}
\newcommand{\threefig}[6]{%
\begin{columns}[T]
\begin{column}{.32\linewidth}\centering\plot{#1}{.42}\fcap{#2}\end{column}
\begin{column}{.32\linewidth}\centering\plot{#3}{.42}\fcap{#4}\end{column}
\begin{column}{.32\linewidth}\centering\plot{#5}{.42}\fcap{#6}\end{column}
\end{columns}}
\newcommand{\fourfig}[8]{%
\begin{columns}[T]
\begin{column}{.49\linewidth}\centering\plot{#1}{.33}\fcap{#2}\plot{#5}{.33}\fcap{#6}\end{column}
\begin{column}{.49\linewidth}\centering\plot{#3}{.33}\fcap{#4}\plot{#7}{.33}\fcap{#8}\end{column}
\end{columns}}

\begin{document}
\framelogo{}

\begin{frame}[plain]\titlepage\end{frame}

%======================================================================
\begin{frame}{This Deck}
\small
\begin{columns}[T]
\begin{column}{.52\linewidth}
\textbf{For the colleague finishing the paper draft} (\texttt{PVFinderPaper\_11\_11\_2026.pdf}).
\begin{enumerate}\setlength\itemsep{.3em}
\item Frozen configuration, one table per regime
\item Datasets and which file is used for what
\item Metric definitions, with the code that computes each one
\item Answers to the 29 questions in the draft, with page and line
\item Answers to Mattia's own questions
\item Results: every paper figure and number, main and backup
\item Proposed internal-note edits
\end{enumerate}
\end{column}
\begin{column}{.45\linewidth}
\textbf{How to read it}
\begin{itemize}\setlength\itemsep{.3em}
\item Every number is a macro in \texttt{numbers.tex}.
\item \TBD{} marks a number from the final recompute that is not in yet. Its source file and key are in the comment next to the macro.
\item A framed red box marks a figure not produced yet, with its expected path.
\item \mainbadge{} figures are proposed for the paper; \backupbadge{} figures go to the note or the backup.
\item PV-Finder and AMVF are always compared on the same events with the same code.
\end{itemize}
\end{column}
\end{columns}
\end{frame}

%======================================================================
\partdivider{1}{Frozen Configuration}{One configuration per regime, used for every result}

\begin{frame}{HL-LHC Configuration ($0\le\mu\le200$)}
\footnotesize
\begin{columns}[T]
\begin{column}{.6\linewidth}
\scriptsize\renewcommand{\arraystretch}{1.07}
\begin{tabular}{@{}p{.25\linewidth}p{.69\linewidth}@{}}
\toprule
Density model & \hlModel: U-Net v2, \hlUnetCh{} channels, \hlLatent{} latent channel \\
 & MLP $7\to$\hlHidden$\to1000$, \hlParams{} parameters \\
Checkpoint & \texttt{\hlCkpt} \\
\midrule
Peak amplitude $\tau$ & \hlAmp \\
Minimum width & \hlWidth{} bins (0.04 mm each) \\
Minimum integral & \hlIntegral \\
Minimum height & \hlHeight \\
Position & local centroid, $\pm$\hlCentroid{} bins about the maximum \\
\midrule
GNN & full-$\mu$ v2, epoch \hlGnnEpoch, \gnnParams{} parameters \\
Graph & $k=\gnnK$ nearest peaks per track \\
Score cut $t$ & \hlT{} (top-1 edge per track kept if score $\ge t$) \\
Output vertices & candidates with $\ge1$ assigned track \\
\midrule
Matching window & \winMM{} mm (all vertex results, both algorithms) \\
Hard-scatter window & \hsWinMM{} mm \\
$\svtx$ fit & counted peaks, \pairBins{} bins over $\pm6$ mm \\
\bottomrule
\end{tabular}
\end{column}
\begin{column}{.37\linewidth}
\begin{itemize}\setlength\itemsep{.25em}
\item One line in code: \texttt{get\_regime("hllhc")} in \texttt{src/\allowbreak pv\_finder/\allowbreak utils/\allowbreak working\_points.py}.
\item The integral and height cuts are the v6 cuts (0.30 / 0.03) scaled to give \hlModel{} the v6 fake rate.
\item The $\svtx$ fit uses the same peaks as the counts. There is no second threshold.
\item Clean requires $\ge\purityCut$ of a vertex's tracks from one truth vertex.
\item Single-vertex residuals use an optimal one-to-one assignment within \assignMM{} mm. This is a measurement radius, not a matching window.
\end{itemize}
\end{column}
\end{columns}
\end{frame}

\begin{frame}{Low Pile-up Configuration (Run 2 MC, Run 2 and Run 3 Data)}
\footnotesize
\begin{columns}[T]
\begin{column}{.6\linewidth}
\scriptsize\renewcommand{\arraystretch}{1.07}
\begin{tabular}{@{}p{.25\linewidth}p{.69\linewidth}@{}}
\toprule
Density model & \lpModel: U-Net v1, \lpUnetCh{} channels, 1 latent channel \\
 & MLP $7\to$\lpHidden$\to1000$, \lpParams{} parameters \\
Checkpoint & \texttt{\lpCkpt} \\
Training & 50 epochs MLP to histogram, 400 end to end; no KDE \\
\midrule
Peak amplitude $\tau$ & \hlAmp \\
Minimum width & \hlWidth{} bins \\
Minimum integral & \lpIntegral \\
Minimum height & none \\
Position & local centroid, $\pm$\hlCentroid{} bins \\
\midrule
GNN & \texttt{\lpGnnCkpt} \\
Score cut $t$ & \lpT{} (same as HL-LHC) \\
Output vertices & candidates with $\ge1$ assigned track \\
\midrule
Matching window & \winMM{} mm \\
Hard-scatter window & \hsWinMM{} mm \\
$\svtx$ fit & counted peaks, \pairBins{} bins over $\pm6$ mm \\
\bottomrule
\end{tabular}
\end{column}
\begin{column}{.37\linewidth}
\begin{itemize}\setlength\itemsep{.25em}
\item \texttt{get\_regime("lowpu")}, same file.
\item One density model for every low pile-up result. The ep300 reproduction model and the two-stage KDE model are no longer used.
\item The window is the HL-LHC value. At low pile-up it is \lpWinRatio{} times the PV-Finder 68\% residual of Run 2 MC.
\item $t=\lpT$: no other $t$ is better in Clean/truth, HS identification and track $F_1$ together on Run 2 MC.
\item There is no Run 3 MC. The simulation is called Run 2 MC.
\end{itemize}
\end{column}
\end{columns}
\end{frame}

\begin{frame}{What Changed with Respect to the 11 November Draft}
\scriptsize
\renewcommand{\arraystretch}{1.2}
\begin{tabular}{@{}p{.19\linewidth}p{.41\linewidth}p{.33\linewidth}@{}}
\toprule
Item & Draft of 11 November & Final \\
\midrule
HL-LHC density model & v6 (280 channels, 4 latent, 3.55 M parameters) & \hlModel{} (\hlUnetCh{} channels, 1 latent, \hlParams) \\
HL-LHC peak cuts & integral 0.30, height 0.03 & \hlIntegral, \hlHeight{} \\
$\svtx$ threshold & separate peak finding at integral 0.50 & the counted peaks \\
Matching window & 0.1694 mm; 0.5 mm in the chain; fitted $\svtx$ at low $\mu$ & \winMM{} mm everywhere \\
HL-LHC GNN & v4 (PU200 only), epoch 107 & full-$\mu$ v2, epoch \hlGnnEpoch \\
GNN score cut & 0.98 (categories), 0.95 (HS) & \hlT{} for every use \\
Low pile-up density model & ep300 (KDE-initialised), PVF-e400, KDE two-stage & \lpModel{} only \\
Low pile-up GNN peaks & integral 0.2, full-region mean & integral \lpIntegral, local centroid \\
Low pile-up $t$ & 0.99, 0.95 ($F_1$), 0.5 (HS), 0.98 (Run 3) & \lpT{} for every use \\
HS selection & plain $\sum p_\mathrm{T}^2$ & Athena 24.0 rule, same code for both \\
HS position & histogram peak & 3D adaptive fit with beam spot \\
Categories & AMVF TTVA links vs GNN tracks & same, plus weighted cross-checks \\
\bottomrule
\end{tabular}
\end{frame}

%======================================================================
\partdivider{2}{Datasets}{Training, validation and test samples}

\begin{frame}{HL-LHC Files}
\scriptsize
\renewcommand{\arraystretch}{1.05}
\begin{tabular}{@{}lrlp{.36\linewidth}@{}}
\toprule
File (mc21 14 TeV $t\bar t$, e8481\_s4494, ITk) & Events & Pile-up & Role \\
\midrule
SingleLep 601229 r16438 & 99{,}800 & 190--210, $\langle\mu\rangle=200.0$ & density training pool \\
SingleLep 601229 r16633 & 99{,}800 & 190--210 & density training pool \\
all-had 601237 r16633, part 1 & 1{,}421{,}400 & 190--210 & density training pool \\
all-had 601237 r16633, part 2 & 1{,}498{,}600 & 190--210 & density training pool \\
all-had 601237 r16633, part 3 & 1{,}200{,}000 & 190--210 & density training pool \\
all-had 601237 r16633, part 4 & 854{,}000 & 190--210 & density training pool \\
\textbf{PU200 pool} & \textbf{\nPoolEvt} & integers 190--210 & \\
\midrule
SingleLep 601229 r16443 & 100{,}000 & flat 0--200, $\langle\mu\rangle=99.6$ & entries 0--24,999: \textbf{test} of every result \\
 & & & entries $\ge25{,}000$: GNN v2 training and validation \\
SingleLep 601229 r16638 & 99{,}800 & flat 0--200, $\langle\mu\rangle=99.6$ & GNN v2 training (\gnnTrainB) \\
 & & & second held-out file for the density model \\
\bottomrule
\end{tabular}
\par\vspace{.6em}
\begin{itemize}\setlength\itemsep{.25em}
\item The density model never saw r16443 or r16638. The GNN never saw r16443 entries 0--24,999.
\item Track selection: ITk paper Table 6. In-time pile-up only. Pile-up is \texttt{ActualNumOfInt}.
\item PU200 = the \nEvtPU{} r16443 test events with $185\le\mu\le215$ (actual 185--200, $\langle\mu\rangle=\meanMuPU$), \nTruthPU{} truth vertices with $n_\mathrm{trk}\ge2$.
\end{itemize}
\src{docs/\allowbreak{}additional\_\allowbreak{}studies/\allowbreak{}paper\_\allowbreak{}questions\_\allowbreak{}2026-11-11.md Part 3; \allowbreak{}outputs/\allowbreak{}11\_\allowbreak{}11\_\allowbreak{}2026\_\allowbreak{}output/\allowbreak{}datasets/\allowbreak{}hllhc\_\allowbreak{}dataset\_\allowbreak{}inventory.json}
\end{frame}

\begin{frame}{HL-LHC Density Model: Training, Validation and Test}
\footnotesize
\begin{columns}[T]
\begin{column}{.46\linewidth}
{\scriptsize\renewcommand{\arraystretch}{1.15}
\begin{tabular}{@{}lrr@{}}
\toprule
Role & 40 mm intervals & Event equiv. \\
\midrule
Training (96\%) & \nTrainSub & \nTrainEq \\
Validation (3\%) & \nValSub & \nValEq \\
Test (1\%, unused) & \nTestSub & \nTestEq \\
\bottomrule
\end{tabular}\par}
\par\vspace{.5em}
\begin{itemize}\setlength\itemsep{.25em}
\item Seed-42 random permutation of the intervals of the PU200 pool. The roles therefore share its pile-up shape: uniform over 190--210, with $\mu=196$ doubled and $\mu=197$ nearly empty (production feature).
\item Results never use the test intervals. They use the independent files r16443 and r16638.
\end{itemize}
\end{column}
\begin{column}{.52\linewidth}\centering
\plot{mu_training}{.3}\par\vspace{.2em}
\begin{minipage}{.48\linewidth}\centering\plot{mu_validation}{.26}\end{minipage}\hfill
\begin{minipage}{.48\linewidth}\centering\plot{mu_test}{.26}\end{minipage}%
\end{column}
\end{columns}
\end{frame}

\begin{frame}{HL-LHC Held-out Files and the GNN v2 Samples}
\footnotesize
\begin{columns}[T]
\begin{column}{.55\linewidth}
\scriptsize\renewcommand{\arraystretch}{1.08}
\begin{tabular}{@{}lp{.4\linewidth}r@{}}
\toprule
GNN v2 set & Source & Events \\
\midrule
Training & r16443 entries $\ge25{,}000$ & \gnnTrainA \\
 & r16638 (stratified in 10-unit $\mu$ bins) & \gnnTrainB \\
 & total & \gnnTrain \\
Validation ($t$ choice) & r16443 entries $\ge25{,}000$ & \gnnVal \\
\quad monitor (checkpoint) & subset of validation & \gnnMon \\
Reserved test & r16443 entries $\ge25{,}000$ & \gnnResTest \\
Reserved test & r16638 & \gnnResTestB \\
\midrule
Paper test & r16443 entries 0--24,999 & 25{,}000 \\
\bottomrule
\end{tabular}
\par\vspace{.4em}
\begin{itemize}\setlength\itemsep{.2em}
\item Graphs use truth-vertex candidates, so no density model enters the GNN training data.
\item No r16638 truth-$z$ or HS $(x,y,z)$ fingerprint occurs in r16443.
\item Resolution vs $\mu$ (before the GNN): \nEvtResMu{} events of both files, fixed per-bin counts chosen before looking at results.
\end{itemize}
\end{column}
\begin{column}{.42\linewidth}\centering
\plot{mu_r16443}{.33}\par\vspace{.2em}
\plot{mu_r16638}{.33}
\end{column}
\end{columns}
\src{outputs/\allowbreak{}11\_\allowbreak{}11\_\allowbreak{}2026\_\allowbreak{}output/\allowbreak{}gnn\_\allowbreak{}fullmu\_\allowbreak{}v2/\allowbreak{}split\_\allowbreak{}manifest.json; \allowbreak{}docs/\allowbreak{}evaluation/\allowbreak{}gnn\_\allowbreak{}full\_\allowbreak{}pileup\_\allowbreak{}training\_\allowbreak{}v2\_\allowbreak{}2026-11-11.md}
\end{frame}

\begin{frame}{Low Pile-up Samples}
\footnotesize
\begin{columns}[T]
\begin{column}{.54\linewidth}
\scriptsize\renewcommand{\arraystretch}{1.08}
\begin{tabular}{@{}lrl@{}}
\toprule
Sample & Events & Use \\
\midrule
Run 2 MC ($t\bar t$, $\mu$ 0--80) & \lpMCTotal & \\
\quad training (70\%) & \lpMCTrain & density model \\
\quad validation (25\%) & \lpMCVal & density model \\
\quad test (5\%) & \lpMCTest & \textbf{every MC result} \\
\quad non-test (train + val) & \lpMCNonTest & GNN training \\
Run 2 data (ZeroBias, run 364076) & 2{,}500 & data \\
Run 3 data (\texttt{file\_3.root}) & 2{,}500 & data \\
\bottomrule
\end{tabular}
\par\vspace{.4em}
\begin{itemize}\setlength\itemsep{.2em}
\item MC test: h5 events 48,450--50,999 (the last 5\%), \lpTruthMC{} truth vertices with $n_\mathrm{trk}\ge2$, $\langle\mu\rangle=\lpMuMC$.
\item Run 2 data: entries 0--2,499. Run 3 data: the first 2,500 accepted entries of 0--3,158. Vertex finding and track association use the same events.
\item Data have no truth: AMVF vertices ($n_\mathrm{trk}\ge2$, beam-spot corrected) are the reference.
\end{itemize}
\end{column}
\begin{column}{.43\linewidth}\centering
\plot{lowpu_mu}{.5}
\fcap{Pile-up of the three evaluation samples.}
\end{column}
\end{columns}
\end{frame}

%======================================================================
\partdivider{3}{Metric Definitions}{How each number is computed, and the code that computes it}

\newcommand{\mtab}[1]{\scriptsize\renewcommand{\arraystretch}{1.22}%
\begin{tabular}{@{}p{.14\linewidth}p{.54\linewidth}>{\raggedright\arraybackslash\tiny}p{.25\linewidth}@{}}
\toprule
Quantity & Definition & Code \\
\midrule
#1
\bottomrule
\end{tabular}}

\begin{frame}{Before the GNN: Efficiency, Fakes, Splits}
\mtab{%
Efficiency (Eq.~1) & truth vertices ($n_\mathrm{trk}\ge2$) with $\ge1$ reco vertex within \winMM{} mm, over all truth vertices; ratio of sums over events. Equals clean + merged truth of the closest-first matcher (merged credit), event by event. & \texttt{compare\_res\_reco} in \texttt{efficiency\_res\_\allowbreak optimized\_atlas.py}; identity test \texttt{tests/\allowbreak test\_amvf\_fair\_\allowbreak matching.py} \\
One-to-one efficiency & truth vertices that win a dedicated reco vertex in the closest-first assignment & same \\
Fake fraction & reco vertices with no truth vertex within \winMM{} mm, over all reco vertices & same (\texttt{fake}) \\
Fakes per event & the same count over the number of events & same \\
Split & reco vertex with a truth vertex in its window that a closer reco vertex already took & same (\texttt{split}) \\
HS reconstruction & events with a reco vertex within \hsWinMM{} mm of the true HS ($z$ only, ITk Sec.~5.5), over all events & \texttt{gnn/\allowbreak evaluation/\allowbreak hs\_metrics.py} \\
}
\par\vspace{.4em}{\footnotesize Truth: \texttt{TruthVertex} with $\ge2$ selected tracks. AMVF: stored vertices with $n_\mathrm{trk}\ge2$. Errors: event bootstrap (2,000 replicas); PV-Finder minus AMVF is paired on the same resampled events.}
\end{frame}

\begin{frame}{Before the GNN: Resolution and the Window}
\mtab{%
$\svtx$ & all within-event pairs of reco vertices; $\Delta z$ in \pairBins{} bins over $\pm6$ mm; fit $a/(1+e^{b(r-|x|)})+c$; $\svtx=|r|$. Same peaks as the counts. Quoted at PU200, and per $\mu$ bin only for $\mu\ge100$. & \texttt{fit\_dz} in \texttt{run\_eval\_\allowbreak pvf\_run3.py}; \texttt{amvf\_fair\_\allowbreak comparison/\allowbreak scan.py} \\
68\% residual & 68th percentile of $|z_\mathrm{reco}-z_\mathrm{true}|$ after the optimal one-to-one assignment within \assignMM{} mm (most matches, then least total distance). The resolution versus $\mu$. & \texttt{matched\_residuals\_\allowbreak optimal}, \texttt{amvf\_fair\_\allowbreak comparison/\allowbreak matching.py} \\
Window & $3\times$ the larger 68\% residual of the two algorithms at PU200, rounded: $3\times\pseAMVF~\mu$m $=0.1694$ mm $\to$ \winMM{} mm & \texttt{working\_points.\allowbreak MATCH\_WINDOW\_MM} \\
Hard-scatter window & \hsWinMM{} mm, ITk performance paper & \texttt{working\_points.\allowbreak HS\_MATCH\_\allowbreak WINDOW\_MM} \\
}
\end{frame}

\begin{frame}{After the GNN: Vertex Categories}
\mtab{%
Assignment & each track keeps its top-1 GNN edge if the score is $\ge t=\hlT$; candidates with no track are dropped from the output & \texttt{gnn/evaluation/\allowbreak chain\_inference.py} \\
Track lists & PV-Finder: GNN-assigned tracks. AMVF: its TTVA links (\texttt{Nonprompt\_Hard\_MaxWeight}, \texttt{RecoVertex\_assocTracks}). One rule for both. & \\
Clean, Merged, Fake & per vertex, the plurality truth vertex of its tracks. Fake if the plurality is unlabelled tracks; Clean if its share is $\ge\purityCut$; Merged otherwise & \texttt{classify\_\allowbreak assignments}, \texttt{gnn/evaluation/\allowbreak classification.py} \\
Split & among vertices with the same plurality truth vertex, all but the largest $\sum p_\mathrm{T}^2$ & same \\
Per event & Clean, Merged, Split, Fake vertices per event (bar plot at PU200, versus $\mu$) & \texttt{hs\_metrics.\allowbreak classify\_vertices} \\
Clean/truth & Clean vertices over truth vertices with $n_\mathrm{trk}\ge2$ & \texttt{chain\_t\_summary} \\
Cross-checks & AMVF with its fit tracks and fit weights; PV-Finder with GNN scores as weights & \texttt{weighted\_\allowbreak classification.py} \\
Bounds & finder cap: peaks matched one to one within \winMM{} mm. Truth-guided: each labelled track to the peak nearest its own truth vertex & \texttt{hs\_metrics}, \texttt{chain\_gap\_\allowbreak decomposition.py} \\
}
\par\vspace{.3em}{\footnotesize Errors: event bootstrap (200 replicas), paired. Every category count is cross-checked by independent code.}
\end{frame}

\begin{frame}{After the GNN: Hard Scatter}
\mtab{%
HS selection & score $=\sum p_\mathrm{T}^2$ over a vertex's tracks with compatibility $\chi^2\le\chiCut$ (Athena 24.0 \texttt{SumPt2Sorting}); selected vertex = argmax. Reconstructed and selected: the selected vertex is the HS match (most HS tracks) and lies within \hsWinMM{} mm of the true HS. AMVF reference: its stored primary. & \texttt{gnn/evaluation/\allowbreak pv\_selection.py}; \texttt{hs\_selection\_\allowbreak identical\_rule.py} \\
HS resolution & events where both algorithms have a vertex within \hsWinMM{} mm of the true HS; $z_\mathrm{reco}-z_\mathrm{HS}$ of the closest; width = iterative RMS within $\pm3$ RMS until stable (Athena \texttt{iterRMS\_convergence}). PV-Finder position: 3D adaptive fit with beam-spot constraint of the GNN-assigned tracks. & \texttt{itk\_resolution.\allowbreak iterative\_rms}; \texttt{hs\_position\_\allowbreak estimators.\allowbreak adaptive\_\allowbreak beamline\_fit} \\
Clean-HS fraction & events with a Clean vertex whose plurality truth is the HS & \texttt{hs\_metrics} \\
HS-track fraction & HS tracks on the HS-matched vertex over all HS tracks & same \\
Local density & truth vertices within $\pm2$ mm of the true HS, HS excluded, over 4 mm (all truth vertices, ITk definition) & \texttt{itk\_resolution.\allowbreak local\_pu\_density} \\
}
\end{frame}

\begin{frame}{After the GNN: Track Association}
\mtab{%
Correct track & assigned to a vertex whose plurality truth vertex is the track's own & \texttt{chain\_t\_metrics} \\
Track efficiency & correct tracks / truth-labelled tracks; unassigned tracks stay in the denominator & same \\
Track purity & correct tracks / assigned tracks with a truth label & same \\
Track $F_1$ & $2\epsilon p/(\epsilon+p)$ & same \\
Wrong rate & $1-$ purity & same \\
Unassigned fraction & $1-$ (assigned labelled tracks / labelled tracks); efficiency $=$ purity $\times(1-$ unassigned fraction$)$ & \texttt{hllhc\_final\_\allowbreak chain\_extras.py} \\
}
\par\vspace{.4em}{\footnotesize Tracks without a truth label enter neither efficiency nor purity, only the assigned fraction.}
\end{frame}

%======================================================================
\partdivider{4}{The 29 Questions}{Answers with page and line of the 11 November draft}

% Answers to the 29 questions of the 11 November draft.
% Source: docs/additional_studies/paper_questions_2026-11-11.md, updated to the
% frozen configuration of 12 November. Page and line refer to
% final_paper/current_paper/paper/PVFinderPaper_11_11_2026.pdf.

\begin{frame}{Q1. Do the 40 mm Intervals Overlap?}
\draftref{p.~7, l.~190}
\qtext{Do the intervals overlap slightly at their boundaries? If so, by how much, and how is a track in an overlap region assigned?}
\small
\begin{itemize}\setlength\itemsep{.45em}
\item \textbf{No overlap.} Interval $k$ covers $[-240+40k,\,-240+40(k+1))$ mm, $k=0,\dots,11$.
\item Each track goes to the one interval that contains its $z_0$; the lower edge is included.
\item Tracks outside $[-240,240)$ mm are dropped. The ntuple selection already requires $|z_0|\le200$ mm.
\item One sentence worth adding: the target histogram is built on the full 12,000-bin grid and then cut into twelve pieces, so a truth vertex near a boundary can put target content into the neighbouring interval, which holds none of its tracks.
\end{itemize}
\vfill\src{src/\allowbreak{}pv\_\allowbreak{}finder/\allowbreak{}data/\allowbreak{}root\_\allowbreak{}to\_\allowbreak{}h5.py:34-36, 67-104, 108-140}
\end{frame}

\begin{frame}{Q2. Is Every Element of Fig.~3 Right, and Does It Cover Run 2/3?}
\draftref{p.~8, Fig.~3 caption}
\small
\begin{columns}[T]
\begin{column}{.52\linewidth}
\textbf{Corrections for the final HL-LHC configuration}
\begin{itemize}\setlength\itemsep{.35em}
\item \textbf{Latent map: $1\times1000$}, not $4\times1000$. The frozen model \hlModel{} has one latent channel (v6 had four).
\item \textbf{Inputs}: seven features per track, $d_0$, $z_0$, $\sigma(d_0)$, $\sigma(z_0)$, $\mathrm{cov}(d_0,z_0)$, interval start, interval end. The helix parameters $(\phi,\theta,q/p)$ of l.~183 do not enter the density network; $\theta$, $\phi$, $p_\mathrm{T}$ enter only the GNN.
\item \textbf{GNN update order}: track-to-vertex and vertex-to-track updates run in parallel from the previous layer (one \texttt{HeteroConv}), not vertices first (l.~307--308).
\end{itemize}
\end{column}
\begin{column}{.45\linewidth}
\textbf{Run 2/3 differences}
\begin{itemize}\setlength\itemsep{.35em}
\item U-Net v1 instead of v2 (Q10).
\item \lpUnetCh{} channels instead of \hlUnetCh.
\item Fully connected track-vertex graphs instead of $k=\gnnK$.
\end{itemize}
\par\vspace{.6em}
\final{one figure covers both regimes if the caption lists these three differences. The latent map is $1\times1000$ in both.}
\end{column}
\end{columns}
\vfill\src{root\_\allowbreak{}to\_\allowbreak{}h5.py:37; \allowbreak{}src/\allowbreak{}gnn/\allowbreak{}models/\allowbreak{}ttva\_\allowbreak{}gat.py:67-121; \allowbreak{}working\_\allowbreak{}points.py}
\end{frame}

\begin{frame}{Q3. Target Widths Come from AMVF}
\draftref{p.~9, l.~205}
\qtext{The current target-width model is derived from AMVF vertex resolutions. Has anything been done to remove this dependence?}
\small
\begin{itemize}\setlength\itemsep{.4em}
\item \textbf{Not for this paper.} The final models are trained on the AMVF-derived widths.
\item Width model: $\sigma(n_\mathrm{trk})=A\,n_\mathrm{trk}^{-B}+C$, floor 2~$\mu$m; truth vertices with fewer than two tracks get one bin (0.04 mm).
\item HL-LHC preset \texttt{hllhc\_alleta}: $A=0.14318$ mm, $B=0.362$, $C=-0.01797$ mm, a background-corrected fit of AMVF-truth residuals (Q29).
\item Only the widths come from AMVF. The target positions are the truth vertex positions.
\item What removing it takes: a preset fitted on the track-only precision (inverse-variance $z_0$ of the truth tracks, already computed as a diagnostic), a rebuild of the \nPoolEvt-event pool and a retrain (\hlModel: 13.7 h on one A100).
\end{itemize}
\vfill\src{root\_\allowbreak{}to\_\allowbreak{}h5.py:82-88; \allowbreak{}src/\allowbreak{}pv\_\allowbreak{}finder/\allowbreak{}data/\allowbreak{}resolution\_\allowbreak{}presets.py}
\end{frame}

\begin{frame}{Q4. U-Net v2 Details}
\draftref{p.~9, l.~223}
\footnotesize
\begin{columns}[T]
\begin{column}{.55\linewidth}
\textbf{Yes, the Appendix E design}, with \hlUnetCh{} channels in the frozen model (\hlModel):
\begin{itemize}\setlength\itemsep{.25em}
\item \textbf{Encoder}: ConvBNReLU($1\to\hlUnetCh$, kernel 15) at 1000 bins; residual block (kernel 5) at 1000; max-pool 2, residual block at 500; max-pool 2, residual block at 250 (bottleneck, factor 4).
\item \textbf{Decoder}, twice: nearest-neighbour upsample $\times2$ + Conv(kernel 3) + BN + ReLU, additive skip, residual block (kernel 5). 250 $\to$ 500 $\to$ 1000.
\item \textbf{Head}: ConvBNReLU($80\to80$, k5), ConvBNReLU($80\to40$, k3), Conv($40\to1$, k1), softplus.
\item Residual block: Conv + BN + ReLU + Dropout, plus identity.
\end{itemize}
\end{column}
\begin{column}{.42\linewidth}
\begin{itemize}\setlength\itemsep{.25em}
\item \textbf{Normalisation}: BatchNorm after every convolution except the last $1\times1$. ReLU throughout.
\item \textbf{Dropout} 0.25 after the ReLU of every convolution block; none in the MLP.
\item \textbf{MLP}: $7\to128\times5\to1000$, LeakyReLU (0.01), softplus; padded tracks masked, output summed over tracks and scaled by $10^{-3}$.
\item \textbf{Parameters}: \hlParams{} = MLP \hlParamsMLP{} + U-Net \hlParamsUnet.
\end{itemize}
\par\vspace{.4em}
\final{Table B.1 (280 channels) and the 3.55 M parameters describe v6. Replace by \hlUnetCh{} channels and \hlParams.}
\end{column}
\end{columns}
\vfill\src{src/\allowbreak{}pv\_\allowbreak{}finder/\allowbreak{}models/\allowbreak{}unet\_\allowbreak{}v2.py; \allowbreak{}autoencoder\_\allowbreak{}models.py:214-352; \allowbreak{}configs/\allowbreak{}vertex\_\allowbreak{}finding/\allowbreak{}config\_\allowbreak{}hllhc\_\allowbreak{}alleta\_\allowbreak{}v6\_\allowbreak{}small\_\allowbreak{}s8.yml}
\end{frame}

\begin{frame}{Q5. Concatenation or Smoothing? One Batch?}
\draftref{p.~9, l.~232}
\small
\begin{itemize}\setlength\itemsep{.5em}
\item \textbf{Plain concatenation in $z$ order.} No smoothing, normalisation or boundary treatment. The intervals do not overlap (Q1).
\item The twelve intervals run as \textbf{one batch of 12}, padded to the widest interval.
\item In Athena the input tensor is $(12, 7, N_\mathrm{trk})$, also one batch.
\item Packed-track inference (CPU timing, M6) runs the MLP only on real tracks and scatters each into its interval. The output equals the padded one to $4\times10^{-6}$; it is the same model.
\end{itemize}
\vfill\src{run\_\allowbreak{}eval\_\allowbreak{}pvf\_\allowbreak{}run3.py:195-215; \allowbreak{}src/\allowbreak{}pv\_\allowbreak{}finder/\allowbreak{}models/\allowbreak{}packed\_\allowbreak{}inference.py}
\end{frame}

\begin{frame}{Q6. Phase 1: U-Net Absent or Frozen?}
\draftref{p.~10, l.~244}
\small
\begin{itemize}\setlength\itemsep{.5em}
\item \textbf{Absent.} Phase 1 runs only the MLP; its latent output is compared directly with the target histogram.
\item The U-Net parameters are also set to \texttt{requires\_grad = False} in Phase 1 and switched on for Phase 2. They have no effect, because the U-Net is not in the forward pass.
\item \final{the paper says ``absent''.}
\end{itemize}
\vfill\src{hllhc\_\allowbreak{}helpers.py:185-195 (phase1\_\allowbreak{}forward); \allowbreak{}train\_\allowbreak{}hllhc\_\allowbreak{}e2e.py:81-90, 128}
\end{frame}

\begin{frame}{Q7. Datasets, Checkpoint, Warm-up, Time, Hardware}
\draftref{p.~10, l.~250}
\footnotesize
\begin{columns}[T]
\begin{column}{.5\linewidth}
\begin{itemize}\setlength\itemsep{.3em}
\item \textbf{Sample}: the six PU200 files of Part 2, \nPoolEvt{} events.
\item \textbf{Split}: seed-42 permutation of 62,083,200 intervals: \nTrainSub{} / \nValSub{} / \nTestSub{} (96/3/1\%). Results use r16443 and r16638 only.
\item \textbf{Schedule of \hlModel}: Phase 1 one epoch, Phase 2 one epoch (v6: 2 + 2).
\item \textbf{Checkpoint}: Phase 2 step 450,000 of 465,624 (97\%). The run stopped on a host-memory error after the last step, before the end-of-epoch checkpoint; step 450,000 is the last saved state. The learning rate there is $1.3\times10^{-6}$, at the cosine floor.
\end{itemize}
\end{column}
\begin{column}{.47\linewidth}
\begin{itemize}\setlength\itemsep{.3em}
\item \textbf{Learning rate}: Phase 1 constant $10^{-3}$. Phase 2 per-step linear warm-up from $10^{-6}$ to $10^{-4}$ over 3,000 steps, then per-step cosine decay to $10^{-6}$.
\item \textbf{Other}: Adam $(0.9, 0.999)$, batch 128, gradient-norm clip 1.0, MSE loss.
\item \textbf{Time and hardware}: one NVIDIA A100-PCIE-40GB. Phase 1 9.1 h, Phase 2 4.6 h (resumed once after a cuDNN initialisation crash on the shared GPU). 13.7 h in total.
\end{itemize}
\end{column}
\end{columns}
\vfill\src{outputs/\allowbreak{}logs/\allowbreak{}hllhc\_\allowbreak{}alleta\_\allowbreak{}v6\_\allowbreak{}s8\_\allowbreak{}1p1ep\_\allowbreak{}20261109\_\allowbreak{}211919.log, \_\allowbreak{}20261110\_\allowbreak{}211930.log; \allowbreak{}build\_\allowbreak{}phase2\_\allowbreak{}scheduler, hllhc\_\allowbreak{}helpers.py:100-160}
\end{frame}

\begin{frame}{Q8. Peak Width, $\sigma_\mathrm{PV}$ and Height}
\draftref{p.~10, l.~262}
\small
\begin{itemize}\setlength\itemsep{.4em}
\item \textbf{Width}: square root of the amplitude-weighted variance of the bin index over the whole above-threshold region (up to a split point), times 0.04 mm.
\item \textbf{$\sigma_\mathrm{PV}$}: this width, used directly in the first edge feature $(z_0-z_\mathrm{PV})/\sqrt{\sigma(z_0)^2+\sigma_\mathrm{PV}^2}$. Median at PU200: \sigmaPVmed{} mm.
\item \textbf{Height}: the maximum bin value of the region, the raw predicted density without normalisation. Second vertex-node feature.
\item \textbf{Training graphs differ}: truth $z$, $\sigma_\mathrm{PV}$ from the \texttt{hllhc\_alleta} $n_\mathrm{trk}$ parameterisation, height from the target recipe. The text near l.~441 (Eq.~4) describes $\sigma_\mathrm{PV}$ as the $n_\mathrm{trk}$ parameterisation; at inference it is the peak width.
\end{itemize}
\vfill\src{efficiency\_\allowbreak{}res\_\allowbreak{}optimized\_\allowbreak{}atlas.py:16-190; \allowbreak{}graph\_\allowbreak{}construction.py:105-168; \allowbreak{}chain\_\allowbreak{}inference.py}
\end{frame}

\begin{frame}{Q9. Final Peak-Finding Configurations; Local Centroid at Low Pile-up?}
\draftref{p.~10, l.~269}
\small
\begin{columns}[T]
\begin{column}{.55\linewidth}
\renewcommand{\arraystretch}{1.2}
\begin{tabular}{@{}lll@{}}
\toprule
 & HL-LHC & Low pile-up \\
\midrule
Amplitude & \hlAmp & \hlAmp \\
Minimum width & \hlWidth{} bins & \hlWidth{} bins \\
Integral & \hlIntegral & \lpIntegral \\
Minimum height & \hlHeight & none \\
Position & centroid $\pm$\hlCentroid & centroid $\pm$\hlCentroid \\
$\svtx$ peaks & counted peaks & counted peaks \\
\bottomrule
\end{tabular}
\end{column}
\begin{column}{.42\linewidth}
\begin{itemize}\setlength\itemsep{.35em}
\item \textbf{Yes}, the local centroid is now used at low pile-up too. On Run 2 MC it changes the efficiency by \lpCentroidShift{} points.
\item The HL-LHC cuts are those of \hlModel. The draft's 0.30 / 0.03 are the v6 values.
\item The note's ``0.50 for the HL-LHC resolution fit'' described the old code. It is gone: the fit uses the counted peaks.
\end{itemize}
\end{column}
\end{columns}
\vfill\src{working\_\allowbreak{}points.py; \allowbreak{}lowpu\_\allowbreak{}final\_\allowbreak{}numbers\_\allowbreak{}2026-11-12.md Sec.~4}
\end{frame}

\begin{frame}{Q10. Run 2/3 Density Model}
\draftref{p.~11, l.~281}
\footnotesize
\begin{columns}[T]
\begin{column}{.52\linewidth}
\textbf{\lpModel{}, every low pile-up result.}
\begin{itemize}\setlength\itemsep{.25em}
\item \textbf{MLP}: $7\to100\times5\to1000$, LeakyReLU (0.01), softplus, 1 latent channel.
\item \textbf{U-Net v1}, \lpUnetCh{} channels, the original design: encoder ConvBNReLU kernels 25, 7, 5, 5 with max-pool, $1000\to500\to250\to125$ (bottleneck 125); decoder three ConvTranspose(2, 2) + ConvBNReLU(k5) with concatenated skips; output Conv($128\to64$, k5), Conv($64\to1$, k5), softplus; dropout 0.25 in every convolution block.
\item \textbf{Parameters}: \lpParams{} trainable.
\item \textbf{Inputs}: the same seven track features as HL-LHC.
\end{itemize}
\end{column}
\begin{column}{.45\linewidth}
\begin{itemize}\setlength\itemsep{.25em}
\item \textbf{Training}: Phase 1, 50 epochs MLP to histogram; Phase 2, 400 epochs end to end; Adam, constant $10^{-3}$ in both phases, batch 128, MSE. Checkpoint: last epoch (400).
\item \textbf{Data}: \lpMCTotal{} Run 2 MC events, 70 / 25 / 5\% (Part 2).
\item \textbf{No KDE in any step}: two-phase, the same procedure as HL-LHC.
\item \textbf{Target preset}: Qi Bin's file; upstream constants (preset \texttt{run3}) $A=0.23817$ mm, $B=0.49491$, $C=-0.00079$ mm.
\end{itemize}
\end{column}
\end{columns}
\vfill\src{configs/\allowbreak{}vertex\_\allowbreak{}finding/\allowbreak{}config\_\allowbreak{}mlp\_\allowbreak{}hist\_\allowbreak{}e2e.yml; \allowbreak{}autoencoder\_\allowbreak{}models.py:801-1000; \allowbreak{}train\_\allowbreak{}mlp\_\allowbreak{}hist\_\allowbreak{}then\_\allowbreak{}e2e.py}
\end{frame}

\begin{frame}{Q11. GNN Details}
\draftref{p.~12, l.~314}
\footnotesize
\begin{columns}[T]
\begin{column}{.5\linewidth}
\begin{itemize}\setlength\itemsep{.3em}
\item \textbf{Layers}: PyG \texttt{GATConv} (original GAT, not GATv2), 4 heads averaged (\texttt{concat=False}), edge features in the attention (\texttt{edge\_dim=3}), no self-loops. Two \texttt{HeteroConv} layers.
\item \textbf{Activations}: LeakyReLU (0.01) after the two linear encoders and the two hidden edge-MLP layers; ReLU after each residual sum; sigmoid score.
\item \textbf{Normalisation}: none. \textbf{Dropout}: none. The stored \texttt{dropout=0.25} is not used in \texttt{forward}, so the trained model has no dropout.
\end{itemize}
\end{column}
\begin{column}{.47\linewidth}
\begin{itemize}\setlength\itemsep{.3em}
\item \textbf{Parameters}: \gnnParams{} (v2 epoch \hlGnnEpoch).
\item \textbf{$d_0$}: as stored in the ntuple. In the HL-LHC simulation the beam spot is at $(0,0,0)$ with $\sigma_{xy}=12~\mu$m, so ``relative to the beam spot'' and ``relative to the origin'' are the same.
\item \textbf{Tracks}: exactly the tracks of the density network; the chain builds both inputs from the same arrays.
\item \textbf{Track features}: $d_0$, $z_0$, $\sigma(d_0)$, $\sigma(z_0)$, cov, $\theta$, $\phi$, $p_\mathrm{T}/50$ GeV.
\end{itemize}
\end{column}
\end{columns}
\vfill\src{src/\allowbreak{}gnn/\allowbreak{}models/\allowbreak{}ttva\_\allowbreak{}gat.py:54-137; \allowbreak{}pu200\_\allowbreak{}chain\_\allowbreak{}graphs.py:189-217; \allowbreak{}constants.py:10}
\end{frame}

\begin{frame}{Q12. Eq.~(1) and the Fake Fraction}
\draftref{p.~14, l.~361}
\small
\begin{itemize}\setlength\itemsep{.45em}
\item \textbf{Eq.~(1) equals the note's efficiency exactly}: truth with $\ge1$ reco vertex within the window = clean + merged truth after the closest-first assignment. Checked event by event on every model and on AMVF. No recalculation of the definition is needed; the numbers are recomputed for the final configuration.
\item \textbf{Fake fraction} (reco vertices without truth within \winMM{} mm, over reco vertices), PU200, \nEvtPU{} events:
\end{itemize}
\begin{center}
\renewcommand{\arraystretch}{1.2}
\begin{tabular}{@{}lccc@{}}
\toprule
 & Fake / reco [\%] & Fakes / event & Splits / event \\
\midrule
PV-Finder (\hlModel) & \fakeFracPVF & \fakeEvtPVF & \splitEvtPVF \\
AMVF & \fakeFracAMVF & \fakeEvtAMVF & \splitEvtAMVF \\
\bottomrule
\end{tabular}
\end{center}
\vfill\src{metric\_\allowbreak{}definitions.md Sec.~3.3a; \allowbreak{}model\_\allowbreak{}freeze\_\allowbreak{}candidates\_\allowbreak{}2026-11-12.md; \allowbreak{}HLDOC}
\end{frame}

\begin{frame}{Q13. One Window for Everything?}
\draftref{p.~14, l.~367}
\qtext{We think we agreed on a single window of 0.1 mm for all results. Is that correct?}
\small
\begin{itemize}\setlength\itemsep{.45em}
\item \textbf{One window, \winMM{} mm}, for every vertex result, both algorithms, both regimes. \textbf{\hsWinMM{} mm for hard-scatter results only}, as in the ITk performance paper.
\item Not 0.1 mm for all vertices: the 68\% single-vertex residual at PU200 is 55--56~$\mu$m for both algorithms, and 0.1 mm is less than twice that.
\item Every other window in the draft is replaced: 0.1694 mm, the 0.5 mm one-to-one match of the chain, the fitted 0.416 mm (Run 2 MC) and the fitted data $\svtx$ are all \winMM{} mm in the final numbers.
\item The 1 mm of the single-vertex residuals is the assignment radius used to measure the residual that sets the window. It is not a matching window and must stay wider than it.
\end{itemize}
\vfill\src{working\_\allowbreak{}points.py; \allowbreak{}matching\_\allowbreak{}window\_\allowbreak{}audit\_\allowbreak{}2026-11-09.md}
\end{frame}

\begin{frame}{Q14. Hard-Scatter Vertex Resolution}
\draftref{p.~15, l.~392}
\small
\begin{columns}[T]
\begin{column}{.55\linewidth}
\textbf{Include it.} Definition (ITk Fig.~28, Athena \texttt{iterRMS\_convergence}):
\begin{itemize}\setlength\itemsep{.3em}
\item events where both algorithms have a vertex within \hsWinMM{} mm of the true HS;
\item residual $z_\mathrm{reco}-z_\mathrm{HS}$ of the closest such vertex;
\item width = iterative RMS: restrict to the mean $\pm3$ RMS and recompute until the selection is stable.
\item PV-Finder position: 3D adaptive vertex fit with a beam-spot constraint of the tracks the GNN assigns to the HS candidate at $t=\hlT$.
\end{itemize}
\end{column}
\begin{column}{.42\linewidth}
\renewcommand{\arraystretch}{1.25}
\begin{tabular}{@{}lc@{}}
\toprule
PU200, \hsResN{} events & Resolution [$\mu$m] \\
\midrule
PV-Finder + GNN & $\hsResGNN\pm\hsResGNNErr$ \\
AMVF & $\hsResAMVF\pm\hsResAMVFErr$ \\
\bottomrule
\end{tabular}
\par\vspace{.6em}
In the note, say only how it is computed (M9).
\end{column}
\end{columns}
\vfill\src{itk\_\allowbreak{}hs\_\allowbreak{}resolution\_\allowbreak{}2026-11-09.md; \allowbreak{}CH hs\_\allowbreak{}resolution\_\allowbreak{}um}
\end{frame}

\begin{frame}{Q15. Local Pile-up Density}
\draftref{p.~15, l.~396}
\small
\begin{itemize}\setlength\itemsep{.45em}
\item \textbf{Include it.} Density as in the ITk paper: truth vertices within $\pm2$ mm of the true HS, HS excluded, over 4 mm. All truth vertices are counted, from Trey's side ntuple (same events, all truth vertices).
\item HS reconstruction at \hsWinMM{} mm is at least \hsRecoDensMin\% in every density bin for both algorithms.
\item HS selection (same Athena rule): PV-Finder + GNN leads AMVF by \hsSelDensLo{} to \hsSelDensHi{} points in every bin.
\item HS resolution: PV-Finder + GNN narrower by \hsResDensLo{} to \hsResDensHi{} $\mu$m in every bin, flat in density.
\item Also available (backup): vertex efficiency before and after the GNN versus local density.
\end{itemize}
\vfill\src{hllhc\_\allowbreak{}supplementary\_\allowbreak{}2026-11-12.md; \allowbreak{}SUP hs\_\allowbreak{}vs\_\allowbreak{}density.json}
\end{frame}

\begin{frame}{Q16. Track Efficiency and Purity Entering $F_1$}
\draftref{p.~16, l.~428}
\small
\begin{itemize}\setlength\itemsep{.45em}
\item \textbf{Correct track}: assigned to a vertex whose plurality truth vertex is the track's own truth vertex.
\item \textbf{Efficiency} $\epsilon$ = correct tracks / all tracks with a truth vertex. Unassigned tracks count in the denominator.
\item \textbf{Purity} $p$ = correct tracks / assigned tracks that have a truth vertex.
\item \textbf{Tracks without truth} are excluded from both; they enter only the assigned fraction.
\item $F_1=2\epsilon p/(\epsilon+p)$.
\end{itemize}
\begin{center}
\renewcommand{\arraystretch}{1.2}
\begin{tabular}{@{}lccc@{}}
\toprule
PU200 & Efficiency [\%] & Purity [\%] & $F_1$ [\%] \\
\midrule
PV-Finder + GNN & \trkEffGNN & \trkPurGNN & \trkFOneGNN \\
AMVF & \trkEffAMVF & \trkPurAMVF & \trkFOneAMVF \\
\bottomrule
\end{tabular}
\end{center}
\vfill\src{threshold\_\allowbreak{}scan.py:149-164, 240-252; \allowbreak{}chain\_\allowbreak{}scan.py:96-130}
\end{frame}

\begin{frame}{Q17. Wrong-Assignment and Unassigned-Track Metrics}
\draftref{p.~16, l.~434}
\small
\begin{itemize}\setlength\itemsep{.4em}
\item \textbf{Include them} (backup figures in Part 6). With $N_\mathrm{lab}$ labelled tracks, $N_\mathrm{at}$ assigned labelled tracks, $N_c$ correct:
\begin{itemize}
\item wrong-assignment rate $=(N_\mathrm{at}-N_c)/N_\mathrm{at}=1-p$: PV-Finder + GNN \wrongGNN\%, AMVF \wrongAMVF\% (PU200);
\item unassigned fraction $u=1-N_\mathrm{at}/N_\mathrm{lab}$;
\item lost unassigned: labelled tracks left unassigned whose truth vertex has no candidate within \winMM{} mm: \lostUnassigned;
\item unlabelled tracks correctly left unassigned / all unlabelled tracks: \unlabelledLeft.
\end{itemize}
\item \textbf{Relation to $F_1$}: $\epsilon=p\,(1-u)$. The wrong rate is the complement of the purity; the unassigned fraction is what separates efficiency from purity.
\end{itemize}
\vfill\src{hllhc\_\allowbreak{}final/\allowbreak{}chain\_\allowbreak{}extras/\allowbreak{}README.md; \allowbreak{}CX}
\end{frame}

\begin{frame}{Q18. Significance-Based Association, Eq.~(4)}
\draftref{p.~16, l.~445}
\small
\begin{columns}[T]
\begin{column}{.52\linewidth}
\begin{itemize}\setlength\itemsep{.35em}
\item \textbf{Eq.~(4) confirmed}: $s=|z_0-z_\mathrm{PV}|/\sqrt{\sigma(z_0)^2+\sigma_\mathrm{PV}^2}$, the same combination as the GNN edge feature.
\item Run on the same graph: each track's $k=\gnnK$ nearest peaks; the track goes to the smallest-$s$ peak if $s<s_\mathrm{max}$. $\sigma_\mathrm{PV}$ = peak width.
\item The best $s_\mathrm{max}$ is chosen on the reported events, which favours the baseline.
\end{itemize}
\end{column}
\begin{column}{.45\linewidth}
\renewcommand{\arraystretch}{1.2}
\begin{tabular}{@{}lc@{}}
\toprule
PU200 & Clean/truth [\%] \\
\midrule
PV-Finder + GNN & \cleanTruthGNN \\
Eq.~(4), best $s_\mathrm{max}$ & \eqFourBest \\
Eq.~(4), $s_\mathrm{max}=3$ & \eqFourThree \\
nearest $z_0$ & \zNear \\
$|z_0-z_\mathrm{PV}|<3\sigma(z_0)$ & \zThree \\
AMVF & \cleanTruthAMVF \\
\bottomrule
\end{tabular}
\end{column}
\end{columns}
\vfill\src{graph\_\allowbreak{}construction.py:128; \allowbreak{}chain\_\allowbreak{}extras README; \allowbreak{}CX}
\end{frame}

\begin{frame}{Q19. Why 25,000 Events and the 1,920 PU200 Events?}
\draftref{p.~17, l.~453}
\small
\begin{itemize}\setlength\itemsep{.45em}
\item \textbf{Compute budget}, not physics. The 25,000 are the first entries of r16443 (100,000 events, flat $\mu$ 0--200, about 500 per integer $\mu$).
\item The \nEvtPU{} are those with $185\le\mu\le215$ (actual 185--200, $\langle\mu\rangle=\meanMuPU$). The whole file has 7,680 events in that range.
\item \textbf{The full chain cannot use more of these files}: r16443 entries $\ge25{,}000$ and r16638 are GNN v2 training and validation data.
\item Before the GNN, more events are used where it helps: r16638 is the second held-out file (\nEvtPUb{} PU200 events) and the resolution versus $\mu$ uses \nEvtResMu{} events of both files.
\end{itemize}
\vfill\src{paper\_\allowbreak{}questions\_\allowbreak{}2026-11-11.md Q19; \allowbreak{}gnn\_\allowbreak{}full\_\allowbreak{}pileup\_\allowbreak{}training\_\allowbreak{}v2\_\allowbreak{}2026-11-11.md}
\end{frame}

\begin{frame}{Q20. Chance Matches, Window Scan, Efficiency vs $\mu$, $n_\mathrm{trk}$, Separation}
\draftref{p.~17, l.~470}
\small
\begin{itemize}\setlength\itemsep{.45em}
\item \textbf{Window scan}: show it as an auxiliary figure. PV-Finder minus AMVF is between \scanLo{} and \scanHi{} points for windows of 0.10--0.20 mm.
\item \textbf{3 mm-shift chance match} at \winMM{} mm: PV-Finder \chancePVF\%, AMVF \chanceAMVF\%. Keep it in the note.
\item \textbf{Efficiency vs $\mu$}: available, proposed for the paper.
\item \textbf{Efficiency vs truth $n_\mathrm{trk}$} (PUB-note style): available before and after the GNN.
\item \textbf{Efficiency vs separation}: available as efficiency versus $\Delta z$ to the nearest truth vertex, and versus local density (backup). The merged credit of Eq.~(1) credits both vertices of a shared peak, so it stays high at small separation; the one-to-one efficiency falls there.
\end{itemize}
\vfill\src{h1\_\allowbreak{}vertex\_\allowbreak{}finding/\allowbreak{}plots; \allowbreak{}chain\_\allowbreak{}extras/\allowbreak{}plots}
\end{frame}

\begin{frame}{Q21. $\svtx$ Threshold, AMVF Value and Fit}
\draftref{p.~17, l.~481}
\small
\begin{itemize}\setlength\itemsep{.45em}
\item \textbf{One threshold.} The fit uses the counted peaks (\hlModel{}: integral \hlIntegral, height \hlHeight). The draft's 0.229 mm came from a second peak finding at integral 0.50; v6 at its counted peaks gives 0.218 mm.
\item \textbf{Final}: PV-Finder $\svtx=\sigPVF\pm\sigPVFErr$ mm, AMVF $\sigAMVF\pm\sigAMVFErr$ mm, same \nEvtPU{} events.
\item \textbf{Fit}: all within-event pairs of vertices ($n_\mathrm{trk}\ge2$ for AMVF), $\Delta z$ histogram with \pairBins{} bins over $\pm6$ mm (0.04 mm), sigmoid $a/(1+e^{b(r-|x|)})+c$, $\svtx=|r|$. Same code for both. The error is the fit error.
\item The sigmoid assumes a flat plateau of pile-up pairs. It describes the data at high $\mu$ only, which is why $\svtx$ is quoted at PU200 (M10).
\end{itemize}
\vfill\src{run\_\allowbreak{}eval\_\allowbreak{}pvf\_\allowbreak{}run3.py fit\_\allowbreak{}dz; \allowbreak{}model\_\allowbreak{}freeze\_\allowbreak{}candidates\_\allowbreak{}2026-11-12.md Sec.~2}
\end{frame}

\begin{frame}{Q22. Merged and Split after the GNN at PU200; Empty Candidates}
\draftref{p.~18, l.~500}
\small
\begin{columns}[T]
\begin{column}{.5\linewidth}
\renewcommand{\arraystretch}{1.2}
\begin{tabular}{@{}lcc@{}}
\toprule
PU200, per reco [\%] & PVF + GNN & AMVF \\
\midrule
Clean & \cleanRecoGNN & \cleanRecoAMVF \\
Merged & \mergedRecoGNN & \mergedRecoAMVF \\
Split & \splitRecoGNN & \splitRecoAMVF \\
Fake & \fakeRecoGNN & \fakeRecoAMVF \\
\midrule
Clean / truth & \cleanTruthGNN & \cleanTruthAMVF \\
\bottomrule
\end{tabular}
\end{column}
\begin{column}{.47\linewidth}
\begin{itemize}\setlength\itemsep{.35em}
\item Table 4 now has the four categories of Table 5.
\item \textbf{Remove candidates without tracks from the output}, not only from the rates. A vertex without tracks cannot be written as an xAOD vertex with tracks. \emptyDropped{} candidates are dropped at PU200.
\item The paper figure is the number of vertices per event in each category (Part 6).
\end{itemize}
\end{column}
\end{columns}
\vfill\src{CH table4; \allowbreak{}SUP classes}
\end{frame}

\begin{frame}{Q23. GNN Thresholds at PU200; One Working Point?}
\draftref{p.~18, l.~514}
\small
\begin{itemize}\setlength\itemsep{.45em}
\item \textbf{One working point, $t=\hlT$}, for vertex classification, HS selection, HS resolution and track $F_1$, in both regimes. The 0.98 and 0.95 values of the draft are dropped.
\item Chosen on the \gnnVal{} GNN validation events (disjoint from the test set) as the most $\mu$-robust value; details in M4.
\item At $t=\hlT$, HS reconstructed and selected at PU200 is \hsSelGNN\% (AMVF stored primary \hsSelAMVF\%).
\end{itemize}
\vfill\src{gnn\_\allowbreak{}v2\_\allowbreak{}freeze\_\allowbreak{}2026-11-12.md Sec.~5; \allowbreak{}working\_\allowbreak{}points.py}
\end{frame}

\begin{frame}{Q24. AMVF Run 2 Resolution and Efficiency on the Same Events}
\draftref{p.~20, l.~528}
\small
\textbf{Both are now on the same \lpMCTest{} Run 2 MC events with the same code.}
\begin{center}\footnotesize
\renewcommand{\arraystretch}{1.25}
\begin{tabular}{@{}lcccc@{}}
\toprule
 & $\svtx$ [mm] & Eff., \winMM{} mm [\%] & Fake/reco [\%] & 68\% res. [$\mu$m] \\
\midrule
PV-Finder (\lpModel) & $\lpSigPVF\pm\lpSigPVFErr$ & $\lpEffPVF\pm\lpEffPVFErr$ & \lpFakeFracPVF & \lpPsePVF \\
AMVF & $\lpSigAMVF\pm\lpSigAMVFErr$ & \lpEffAMVF & \lpFakeFracAMVF & \lpPseAMVF \\
\bottomrule
\end{tabular}
\end{center}
\begin{itemize}\setlength\itemsep{.3em}
\item The draft's AMVF 0.793 mm was on 2,500 other events; the PV-Finder 0.416 mm used 60 bins and the ep300 model.
\end{itemize}
\vfill\src{LPDOC R1; \allowbreak{}VF mc}
\end{frame}

\begin{frame}{Q25. Low Pile-up GNN at Integral 0.2, Vertex Finding at 0.50}
\draftref{p.~20, l.~535}
\small
\begin{itemize}\setlength\itemsep{.45em}
\item \textbf{Final: integral \lpIntegral{} for both}, on \lpModel{} peaks with the local centroid.
\item The difference was more than the threshold: the draft's GNN numbers (77.43\% Clean/truth, $t=0.99$) came from PVF-e400 at integral 0.2 and the full-region mean; the vertex finding from the ep300 model at 0.50.
\item \textbf{Recalculated}: graphs rebuilt from the frozen model at \lpIntegral, GNN at $t=\lpT$:
\begin{itemize}
\item Clean/truth \lpCleanTruth\% (AMVF \lpCleanTruthAMVF\%);
\item Clean / Merged / Split / Fake per reco: \lpClean{} / \lpMerged{} / \lpSplit{} / \lpFake\%.
\end{itemize}
\item Integral \lpIntegral{} finds fewer candidates than 0.2, mostly small vertices; that and the lower $t$ explain the change from 77.43\%.
\end{itemize}
\vfill\src{LPDOC R2, Sec.~4; \allowbreak{}GM}
\end{frame}

\begin{frame}{Q26. GNN Thresholds in Run 2 Simulation}
\draftref{p.~21, l.~548}
\small
\begin{itemize}\setlength\itemsep{.45em}
\item \textbf{One $t=\lpT$} for vertex classification, track $F_1$ and HS identification. The 0.99 / 0.95 / 0.5 of the draft are dropped.
\item Rule: the HL-LHC value unless clearly worse at low pile-up. No $t$ beats \lpT{} in Clean/truth, HS identification and $F_1$ together by more than two paired standard errors.
\item At $t=\lpT$: Clean/truth \lpCleanTruth\%, $F_1$ \lpFOne\% (AMVF \lpFOneAMVF\%), HS identification \lpHSID\% (AMVF \lpHSIDAMVF\%).
\item $t=0.99$ gives Clean/truth \lpCleanTruthNN\%, with a lower $F_1$; its gain comes from $\mu>40$ only, and at $\mu<20$ $t=\lpT$ is better.
\item The draft's 98.98\% HS identification was obtained at $t=0.99$, not at the 0.5 the note lists.
\end{itemize}
\vfill\src{LPDOC Sec.~3; \allowbreak{}L/\allowbreak{}gnn\_\allowbreak{}mc/\allowbreak{}t\_\allowbreak{}options.md}
\end{frame}

\begin{frame}{Q27. Run 3 Data Association: $t$, Peak Threshold, Events}
\draftref{p.~22, l.~574}
\small
\begin{columns}[T]
\begin{column}{.5\linewidth}
\begin{itemize}\setlength\itemsep{.4em}
\item \textbf{Final}: $t=\lpT$, \lpModel{} at integral \lpIntegral{} with the local centroid.
\item \textbf{Events}: the same 2,500 Run 3 events and the same peaks as the vertex finding. The draft's 2,000 events were a random sample of \texttt{file\_3.root}; only 8 were in the vertex-finding set.
\item The draft's model was the two-stage KDE model at $t=0.98$; both are replaced.
\end{itemize}
\end{column}
\begin{column}{.47\linewidth}
\renewcommand{\arraystretch}{1.2}
\begin{tabular}{@{}lc@{}}
\toprule
Run 3 data, $t=\lpT$ & [\%] \\
\midrule
Tracks assigned (GNN) & \lpRthreeAssigned \\
Agreement with AMVF, 1 mm & \lpRthreeAgree \\
Agreement with AMVF, \winMM{} mm & \lpRthreeAgreeW \\
Leading-vertex agreement & \lpRthreeLead \\
\bottomrule
\end{tabular}
\par\vspace{.4em}
{\footnotesize Agreement is over tracks assigned by both. These are agreement numbers, not efficiencies.}
\end{column}
\end{columns}
\vfill\src{LPDOC R4; \allowbreak{}R3A}
\end{frame}

\begin{frame}{Q28. ``integral=0.2'' in the Fig.~10 Panel Titles}
\draftref{p.~23, Fig.~10 caption}
\small
\begin{itemize}\setlength\itemsep{.45em}
\item \textbf{0.5 was used}; the titles were stale. Re-running the old checkpoint on 25 Run 2 MC events reproduces the stored peak lists in 23 of 25 events at 0.5 and in 7 of 25 at 0.2; the other 2 differ by CPU/GPU rounding.
\item The label came from the April 2026 plotting code, which carried the then-default value.
\item \textbf{Final}: all panels are regenerated with \lpModel{} at the final configuration. Every panel states ``integral threshold 0.50''; simulation panels say ``Simulation Internal'', data panels ``Internal''; the Run 3 panels name the simulation Run 2 MC.
\end{itemize}
\vfill\src{LPDOC Sec.~6; \allowbreak{}L/\allowbreak{}summary/\allowbreak{}figures/\allowbreak{}fig10*}
\end{frame}

\begin{frame}{Q29. AMVF Dependence; Sample of \texttt{hllhc\_alleta}; Run 2/3 Parameters}
\draftref{p.~28, l.~652}
\small
\begin{itemize}\setlength\itemsep{.45em}
\item \textbf{AMVF dependence}: Q3.
\item \textbf{\texttt{hllhc\_alleta}} was fitted on the \textbf{full-$|\eta|$ sample} ($|\eta|<4$): 601229 r16633, 99,800 events, 2.31 M truth vertices with $n_\mathrm{trk}\ge2$, background-corrected fit of AMVF-truth residuals. Not the $|\eta|<2.5$ sample.
\item Parameters: $A=0.14318$ mm, $B=0.362$, $C=-0.01797$ mm.
\item \textbf{Run 2/3} (preset \texttt{run3}, Qi Bin's file): $A=0.23817443$ mm, $B=0.49491396$, $C=-0.000787436$ mm.
\end{itemize}
\vfill\src{src/\allowbreak{}pv\_\allowbreak{}finder/\allowbreak{}data/\allowbreak{}resolution\_\allowbreak{}presets.py; \allowbreak{}outputs/\allowbreak{}07\_\allowbreak{}31\_\allowbreak{}2026\_\allowbreak{}output/\allowbreak{}amvf\_\allowbreak{}resolution\_\allowbreak{}alleta/\allowbreak{}}
\end{frame}


%======================================================================
\partdivider{5}{Mattia's Questions}{Decisions behind the configuration}

% Answers to Mattia's own questions (meeting of 11 November 2026).

\begin{frame}{M1. Why a \winMM{} mm Window}
\small
\begin{columns}[T]
\begin{column}{.5\linewidth}\footnotesize
\begin{itemize}\setlength\itemsep{.3em}
\item The window asks whether a vertex was put where the truth vertex is. It is set by how precisely one vertex is placed, the same for both algorithms.
\item PU200, \nEvtPU{} events: reco and truth vertices assigned one to one within \assignMM{} mm (most matches, then least total distance). 68\% residual: PV-Finder \pseVsix{} $\mu$m, AMVF \pseAMVF{} $\mu$m.
\item Window $=3\times$ the larger $=3\times\pseAMVF~\mu\mathrm{m}=0.1694$ mm, rounded to \winMM{} mm. Not a Gaussian $3\sigma$: the residuals have tails.
\item A greedy assignment would bias the width about \greedyBias\% low.
\end{itemize}
\end{column}
\begin{column}{.47\linewidth}\centering
\plot{window_decision}{.42}
\fcap{Single-vertex residuals at PU200 (v6, 9 November). The \hlModel{} version is in Part 6.}
\par\vspace{.3em}
{\footnotesize Not $\svtx$: it measures how close two vertices can be and still be separated, and differs by 40\% between the algorithms. At low pile-up the window is \lpWinRatio{} times the PV-Finder residual.}
\end{column}
\end{columns}
\end{frame}

\begin{frame}{M2. Hard Scatter at \hsWinMM{} mm, for ITk Consistency}
\small
\begin{itemize}\setlength\itemsep{.5em}
\item The ITk performance paper (arXiv:2412.15090, Sec.~5.5, Fig.~26b) reconstructs the hard-scatter vertex with a \hsWinMM{} mm criterion and quotes only the hard scatter. The HS has tens of tracks and a resolution near 10~$\mu$m, so \hsWinMM{} mm is wide for it.
\item The all-vertex population is dominated by low-multiplicity pile-up vertices with a 68\% residual of 55--56~$\mu$m; for them \hsWinMM{} mm would mix position precision into the efficiency. They keep \winMM{} mm.
\item Every HS number uses \hsWinMM{} mm: reconstruction, selection, resolution.
\item HS reconstruction at PU200: PV-Finder \hsRecoPVF\%, AMVF \hsRecoAMVF\%.
\end{itemize}
\takeaway{Sentence for the paper: ``For consistency with ITk performance studies, which quote only the hard-scatter vertex, we also give the hard-scatter results with a 0.1~mm window.''}
\end{frame}

\begin{frame}{M3. Why \hlModel{} Was Frozen}
\footnotesize
\begin{columns}[T]
\begin{column}{.55\linewidth}
\scriptsize\renewcommand{\arraystretch}{1.05}
\begin{tabular}{@{}p{.5\linewidth}p{.44\linewidth}@{}}
\toprule
\hlModel{} minus v6, same events & \\
\midrule
Efficiency at the v6 fake rate (r16443, r16638) & \sEightDEffA{}, \sEightDEffB{} pt \\
$\svtx$ & \sEightDSig{} $\mu$m \\
68\% residual & \sEightDpse{} $\mu$m \\
Clean/truth after the GNN ($t=\hlT$) & \sEightDClean{} pt \\
HS selection after the GNN & \sEightDHSsel{} pt \\
HS resolution after the GNN & \sEightDHSres{} $\mu$m \\
\midrule
CPU speed-up, packed, vs v6 as deployed & \spdSEightNov$\times$ (9 Nov), \spdPack$\times$ (12 Nov) \\
Parameters & \hlParams{} vs 3,547,993 \\
\bottomrule
\end{tabular}
\end{column}
\begin{column}{.42\linewidth}\scriptsize
\begin{itemize}\setlength\itemsep{.25em}
\item Bar: efficiency loss below 0.2 points at the v6 fake rate, and no worse resolution. Only \hlModel{} meets it.
\item s5 (112 channels) separates close vertices worse: $\svtx$ \sFiveDSig{} $\mu$m at 6--7 standard deviations, with either schedule.
\item Cuts: v6's 0.30 / 0.03 scaled until the fakes per event match v6; the two-file mean is \hlIntegral{} / \hlHeight.
\item After the GNN the trends with $\mu$ are the same for both models.
\end{itemize}
\end{column}
\end{columns}

\vfill\src{model\_\allowbreak{}freeze\_\allowbreak{}candidates\_\allowbreak{}2026-11-12.md; \allowbreak{}gnn\_\allowbreak{}v2\_\allowbreak{}freeze\_\allowbreak{}2026-11-12.md Sec.~7; \allowbreak{}cpu\_\allowbreak{}small\_\allowbreak{}model\_\allowbreak{}2026-11-09.md}
\end{frame}

\begin{frame}{M3. \hlModel{} and v6 after the GNN}
\centering
\plot{v6_vs_s8_gnn}{.72}
\fcap{Chain metrics versus $\mu$ for both density models with GNN v2, test events (12 November freeze). The trends are the same; the differences are at most a few tenths of a point.}
\end{frame}

\begin{frame}{M4. The GNN Score Cut $t=\hlT$}
\footnotesize
\begin{columns}[T]
\begin{column}{.6\linewidth}
\scriptsize\setlength{\tabcolsep}{3pt}\renewcommand{\arraystretch}{1.1}
\begin{tabular}{@{}llccccc@{}}
\toprule
Rule (on validation) & $t$ & \multicolumn{3}{c}{Clean/truth [\%]} & $F_1$ & HS res \\
 & & incl. & PU200 & $\mu<20$ & [\%] & [$\mu$m] \\
\midrule
A: max Clean/truth & 0.98 & \optAclean & \optAcleanPU & \optAcleanLow & \optAFone & \optAHSres \\
\textbf{D: most $\mu$-robust} & \textbf{0.97} & \optDclean & \optDcleanPU & \optDcleanLow & \optDFone & \optDHSres \\
B: max HS selection & 0.94 & \optBclean & & & & \\
C: max track $F_1$ & 0.92 & \optCclean & & & & \\
E: no cut & 0 & \optEclean & & & & \\
AMVF & & \cleanTruthAMVFincl & \cleanTruthAMVF & \optAMVFcleanLow & \optAMVFFone & \optAMVFHSres \\
\bottomrule
\end{tabular}
\par\vspace{.3em}
{\tiny Test events, v6 peaks (12 November freeze). D: largest worst-case ratio, over five $\mu$ bins, of Clean/truth to its best value in the bin.}
\end{column}
\begin{column}{.37\linewidth}\scriptsize
\begin{itemize}\setlength\itemsep{.2em}
\item Checkpoint: epoch \hlGnnEpoch, the lowest validation-monitor loss (\gnnValLoss) of 540 shard epochs; plateaued at the learning-rate floor.
\item $t=\hlT$ is within \optDvsAclean{} points of the maximum Clean/truth, gains \optDvsAlow{} points at $\mu<20$, \optDvsAFone{} points of $F_1$ and \optDvsAres~$\mu$m of HS resolution.
\item $t=0.99$ is worse than 0.98 in Clean/truth, $F_1$, HS selection and HS resolution.
\item Chosen by Mattia, 12 November; also used at low pile-up.
\end{itemize}
\end{column}
\end{columns}

\vfill\src{gnn\_\allowbreak{}v2\_\allowbreak{}freeze\_\allowbreak{}2026-11-12.md Sec.~1, 5; \allowbreak{}G/\allowbreak{}options.md}
\end{frame}

\begin{frame}{M4. The $t$ Scan in $\mu$ Bins}
\centering
\plot{t_scan_val_s8}{.72}
\fcap{All chain metrics versus $t$ per $\mu$ bin, GNN validation events, \hlModel{} peaks.}
\end{frame}

\begin{frame}{M5. Behaviour of the GNN at Low $\mu$}
\footnotesize
\begin{columns}[T]
\begin{column}{.55\linewidth}\scriptsize
\textbf{Seen} at $\mu<20$: Clean/truth peaks at $t=\lmCleanPeakT$ and falls at higher $t$; above $\mu=50$ it rises up to 0.98.
\par\vspace{.3em}
\textbf{Why}, measured:
\begin{itemize}\setlength\itemsep{.15em}
\item At $\mu<10$, \lmTopOneLow\% of top-1 edges are correct (\lmTopOneHigh\% at $\mu$ 150--200). A higher $t$ then removes mostly good tracks.
\item The network is less confident on small vertices at low $\mu$: \lmSmallBelowLow\% of correct tracks of 5--10-track vertices score below 0.99 at $\mu<10$ (\lmSmallBelowMid\% at $\mu$ 20--50, \lmSmallBelowHigh\% at 150--200), on truth-candidate graphs too.
\item $\mu<20$ events hold about \lmTrackShare\% of the tracks of a flat-$\mu$ sample, so they weigh little in the loss.
\item v1 was worse (\lmVoneBelow\% vs \lmVtwoBelow\% of correct tracks below 0.99 at $\mu<10$): the behaviour Mattia saw.
\item Splits at low $\mu$ are HS satellites: \lmSplitHS\% have the HS as dominant truth, \lmSplitOne\% hold one track.
\end{itemize}
\end{column}
\begin{column}{.42\linewidth}
\begin{itemize}\setlength\itemsep{.25em}
\item Fake/reco stays at or below \lmFakeMax\%; there is no failure mode.
\item The satellite splits are weak density peaks next to the large HS peak, a larger share of the list when there are few vertices. They are not GNN errors.
\item $t=\hlT$ is the compromise: within \optDvsAclean{} points of the best Clean/truth at high $\mu$, \optDvsAlow{} points better than 0.98 at $\mu<20$.
\end{itemize}
\end{column}
\end{columns}
\par\vspace{.2em}
{\scriptsize At the working point, $\mu$ 0--10: track $F_1$ \lmFOneGNN\% (PV-Finder + GNN) vs \lmFOneAMVF\% (AMVF). No track-$F_1$ claim at low $\mu$ in the paper.}
\end{frame}

\begin{frame}{M5. Low-$\mu$ Event Displays}
\centering
\plot{lowmu_displays}{.72}
\fcap{Validation events at $\mu\approx17$: isolated 7- and 11-track vertices whose tracks all score 0.9--0.99, so the vertex vanishes at $t=0.99$.}
\end{frame}

\begin{frame}{M6. CPU Timing against AMVF}
\footnotesize
\begin{columns}[T]
\begin{column}{.5\linewidth}
\textbf{Projected Athena CPU time} (Xeon 4216, one thread), from speed-ups measured on sneezy and the note's Athena table:
\[T_\mathrm{chain}=\frac{\athCPUinf~\mathrm{ms}}{S}+\athOther~\mathrm{ms}+\athGNN~\mathrm{ms}\]
\renewcommand{\arraystretch}{1.2}
\begin{tabular}{@{}lcc@{}}
\toprule
Chain & $S$ & $T_\mathrm{chain}$ [ms] \\
\midrule
v6 padded (Athena today) & 1 & \athCPUms \\
\hlModel{} packed + GNN v2 & \spdPack & \projChain \\
\hlModel{} padded + GNN v2 & \spdPad & \projChainPad \\
AMVF (note figure, $\mu=200$) & & \athAMVF{} (\athAMVFlo--\athAMVFhi) \\
\bottomrule
\end{tabular}
\end{column}
\begin{column}{.47\linewidth}
\begin{itemize}\setlength\itemsep{.35em}
\item With packed inference the chain is \projFactor$\times$ faster than AMVF on CPU. The worst case tried, \projWorst{} ms, is still below \athAMVFlo{} ms.
\item Packed inference needs a new Athena tensor builder: a track list plus an interval index.
\item With today's padded tensor, \hlModel{} is comparable to AMVF, not faster.
\item The GNN is now the largest step (\athGNN{} ms in Athena).
\item On GPU the full chain takes \athGPUms{} ms (note).
\item AMVF cannot run on sneezy: its time is read off the note's figure. The chain time is a projection, not an Athena measurement.
\end{itemize}
\end{column}
\end{columns}
\vfill\src{H/\allowbreak{}cpu\_\allowbreak{}timing/\allowbreak{}README.md; \allowbreak{}note sections/\allowbreak{}04b\_\allowbreak{}timing.tex, tab:athena-timing-breakdown}
\end{frame}

\begin{frame}{M7. Binning: Is the Network or the 0.04 mm Grid the Limit?}
\footnotesize
\qtext{Professor's question: does the smaller Run 2 $\to$ HL-LHC gain of the AMVF/PV-Finder ratio mean that PV-Finder does not use the better HL-LHC tracks, so that finer bins are needed?}
\begin{columns}[T]
\begin{column}{.52\linewidth}
Test: the training target is what a perfect network would output. It goes through the same peak finder and the same metrics.
\par\vspace{.3em}
\renewcommand{\arraystretch}{1.15}
\begin{tabular}{@{}lcc@{}}
\toprule
$\svtx$ [mm] & Run 2 MC & HL-LHC PU200 \\
\midrule
Target, 0.04 mm grid & \tgtSigRunTwo & \tgtSigHL{} (\tgtSigHLwone$^\ast$) \\
Target, 0.02 mm grid$^\ast$ & & \tgtSigHLfine \\
PV-Finder & \pvfSigRunTwo & \pvfSigHLvsix{} (v6) \\
AMVF & \amvfSigRunTwo & \sigAMVF \\
\bottomrule
\end{tabular}
\par{\scriptsize $^\ast$1-bin minimum width, so that every HS target spike is found.}
\end{column}
\begin{column}{.45\linewidth}
\begin{itemize}\setlength\itemsep{.3em}
\item PV-Finder is 2.0 times its target in Run 2 and 1.5--1.7 times at HL-LHC: closer to its ideal at HL-LHC.
\item A 0.02 mm grid improves the ideal by 10\%, far from the gap to PV-Finder.
\item The 0.04 mm grid carries the HS position to \tgtHSum~$\mu$m; PV-Finder reaches \pvfHSum~$\mu$m.
\item Track-based ideal precision of pile-up vertices: \idealPURunTwo{} $\to$ \idealPUHL~$\mu$m (1.7$\times$). At the same peak setting PV-Finder $\svtx$ also improves 1.7$\times$ (\pvfSigRunTwo{} $\to$ \pvfSigHLrunTwoCfg{} mm).
\end{itemize}
\end{column}
\end{columns}
\takeaway{The headroom is in the network, not in the grid. Finer bins are not needed.}
\end{frame}

\begin{frame}{M8. Hard-Scatter Selection with One Identical Rule}
\footnotesize
\begin{columns}[T]
\begin{column}{.5\linewidth}
\textbf{The Athena 24.0 rule} (read from source):
\begin{itemize}\setlength\itemsep{.25em}
\item vertices sorted by $\sum p_\mathrm{T}^2$ of their linked tracks (\texttt{SumPt2Sorting}); the first is \texttt{PriVtx};
\item a track is linked if its compatibility with the vertex is $\chi^2\le\chiCut$; no weights, no $p_\mathrm{T}$ cap;
\item HS match: the vertex with the most HS tracks (\texttt{bestHardScatterMatch}).
\end{itemize}
\textbf{Validation on AMVF} (first 3,000 entries): on its stored lists our implementation agrees with the stored primary in \ruleAgreeAll\% of events (\ruleAgreePU\% at PU200); plain $\sum p_\mathrm{T}^2$ gives \plainAgreeAll\%.
\end{column}
\begin{column}{.47\linewidth}
\renewcommand{\arraystretch}{1.2}
\begin{tabular}{@{}lc@{}}
\toprule
PU200, reco and selected (\hsWinMM{} mm) & [\%] \\
\midrule
PV-Finder + GNN, Athena rule & $\hsSelGNN\pm\hsSelGNNErr$ \\
AMVF, stored primary & \hsSelAMVF \\
AMVF, same rule, our code & \hsSelAMVFrule \\
AMVF, plain $\sum p_\mathrm{T}^2$ (old note) & \hsSelAMVFplain \\
\midrule
PV-Finder + GNN $-$ AMVF stored & $\hsSelDiff\pm\hsSelDiffErr$ pt \\
\bottomrule
\end{tabular}
\par\vspace{.4em}
\begin{itemize}\setlength\itemsep{.25em}
\item Headline: GNN with the Athena rule against AMVF's own stored primary. Our code can only lose AMVF information, so this favours AMVF.
\item It replaces the note's +3.4 points, which compared different rules.
\end{itemize}
\end{column}
\end{columns}
\vfill\src{hs\_\allowbreak{}selection\_\allowbreak{}identical\_\allowbreak{}rule\_\allowbreak{}2026-11-11.md; \allowbreak{}src/\allowbreak{}gnn/\allowbreak{}evaluation/\allowbreak{}pv\_\allowbreak{}selection.py}
\end{frame}

\begin{frame}{M9. Hard-Scatter Resolution: Why PV-Finder + GNN Reaches AMVF}
\footnotesize
\begin{columns}[T]
\begin{column}{.5\linewidth}
\renewcommand{\arraystretch}{1.15}
\begin{tabular}{@{}lc@{}}
\toprule
HS position (GNN v4 study) & [$\mu$m] \\
\midrule
Histogram peak, local centroid & \hsResPeak \\
1D adaptive fit of $z_0$, GNN tracks & \hsResOneD \\
3D adaptive fit, all tracks within 1 mm & \hsResNoGNN \\
3D adaptive fit, AMVF's own HS tracks & \hsResAMVFlist \\
3D adaptive fit, true HS tracks & \hsResTruthTrk \\
AMVF stored vertex & \hsResAMVFstudy \\
\midrule
\textbf{Final}, 3D fit, GNN tracks, $t=\hlT$ & $\hsResGNN\pm\hsResGNNErr$ \\
\textbf{Final}, AMVF & $\hsResAMVF\pm\hsResAMVFErr$ \\
\bottomrule
\end{tabular}
\end{column}
\begin{column}{.47\linewidth}
\begin{itemize}\setlength\itemsep{.3em}
\item The missing piece was the transverse vertex position. ITk tracks reach $|\eta|=4$ (median $|\cot\theta|=3.2$), so a 10--15~$\mu$m offset from the beam axis moves $z_0$ by tens to hundreds of $\mu$m.
\item The fit: $(x,y,z)$ with a beam-spot constraint, $(d_0,z_0)$ covariance, adaptive weights with annealing, seeded at the histogram peak.
\item GNN tracks do as well as the true HS tracks; on AMVF's own tracks the same fit beats AMVF's stored position. The fit model, not the association, made the difference.
\item In the note: say only how it is computed.
\end{itemize}
\end{column}
\end{columns}
\vfill\src{itk\_\allowbreak{}hs\_\allowbreak{}resolution\_\allowbreak{}2026-11-09.md (2026-11-11 sections); \allowbreak{}hs\_\allowbreak{}position\_\allowbreak{}estimators.adaptive\_\allowbreak{}beamline\_\allowbreak{}fit}
\end{frame}

\begin{frame}{M10. Resolution versus $\mu$: the 68\% Residual}
\footnotesize
\begin{columns}[T]
\begin{column}{.5\linewidth}
\textbf{Choice}: the resolution versus $\mu$ in the paper is the 68\% residual $|z_\mathrm{reco}-z_\mathrm{true}|$. $\svtx$ is quoted at PU200 only, or per bin for $\mu\ge100$.
\par\vspace{.4em}
\textbf{Why}:
\begin{itemize}\setlength\itemsep{.15em}
\item Both algorithms make more close vertex pairs than truth has: close-pair rate over flat expectation \rvmExcessPVF{} (PV-Finder), \rvmExcessAMVF{} (AMVF), \rvmExcessTruth{} (truth) at $\mu<20$. The hard-scatter neighbourhood is the largest source.
\item This excess is about constant per event while random pairs grow as $\mu^2$, so it dominates at low $\mu$. Below $\mu=100$ the sigmoid fit measures a notch inside it, not the separation ($\chi^2/\mathrm{ndf}$ \rvmChiLo--\rvmChiHi).
\item $\svtx$ is not flat in $\mu$ for either algorithm, so a single all-$\mu$ value is wrong.
\item The 68\% residual is defined at every $\mu$; the two algorithms agree within 2~$\mu$m in every bin.
\end{itemize}
\end{column}
\begin{column}{.47\linewidth}\centering
\plot{resolution_vs_mu}{.53}
\fcap{68\% residual (top) and $\svtx$ (bottom; open markers below $\mu=100$), \hlModel.}
\end{column}
\end{columns}
\vfill\src{resolution\_\allowbreak{}vs\_\allowbreak{}mu\_\allowbreak{}2026-11-11.md; \allowbreak{}H/\allowbreak{}resolution\_\allowbreak{}vs\_\allowbreak{}mu}
\end{frame}


%======================================================================
\partdivider{6}{Results}{Every paper figure and number; PV-Finder and AMVF on the same events}

% Results: every paper figure and number at the frozen configuration.
% Figures through images.txt; placeholders until the recompute agents finish.

%---------------------------------------------------------------------- HL-LHC before the GNN
\begin{frame}{HL-LHC before the GNN: Table 3 at PU200}
\mainbadge\hfill{\footnotesize r16443, \nEvtPU{} events, $\langle\mu\rangle=\meanMuPU$, \nTruthPU{} truth vertices ($n_\mathrm{trk}\ge2$), window \winMM{} mm}
\par\vspace{.3em}\small
\begin{center}
\renewcommand{\arraystretch}{1.12}
\begin{tabular}{@{}lccc@{}}
\toprule
 & PV-Finder (\hlModel) & AMVF & PVF $-$ AMVF \\
\midrule
Efficiency, Eq.~(1) [\%] & $\effPVF\pm\effPVFErr$ & \effAMVF & $\effDiff\pm\effDiffErr$ pt \\
One-to-one efficiency [\%] & \effOneOnePVF & \effOneOneAMVF & \\
Fake / reco [\%] & \fakeFracPVF & \fakeFracAMVF & \\
Fakes / event & \fakeEvtPVF & \fakeEvtAMVF & \\
Splits / event & \splitEvtPVF & \splitEvtAMVF & \\
Reco vertices / event & \recoEvtPVF & \recoEvtAMVF & (truth \truthEvt) \\
$\svtx$ [mm] & $\sigPVF\pm\sigPVFErr$ & $\sigAMVF\pm\sigAMVFErr$ & \\
68\% residual [$\mu$m] & \psePVF & \pseAMVF & \\
HS reconstruction, \hsWinMM{} mm [\%] & \hsRecoPVF & \hsRecoAMVF & \\
\midrule
Efficiency on r16638 (\nEvtPUb{} events) [\%] & \effPVFb & \effAMVFb & \\
\bottomrule
\end{tabular}
\end{center}
\src{HLDOC; \allowbreak{}H1 = H/\allowbreak{}h1\_\allowbreak{}vertex\_\allowbreak{}finding/\allowbreak{}h1\_\allowbreak{}summary.json (table3.tex)}
\end{frame}

\begin{frame}{HL-LHC before the GNN: Efficiency and Fakes versus $\mu$}
\mainbadge
\twofig{eff_vs_mu}{Vertex-finding efficiency, Eq.~(1), \winMM{} mm, r16443 entries 0--24,999.}%
{fake_fraction_vs_mu}{Fake fraction: reco vertices without truth within \winMM{} mm.}
\end{frame}

\begin{frame}{HL-LHC before the GNN: Fakes and Yield versus $\mu$}
\mainbadge
\twofig{fakes_per_event_vs_mu}{Fake vertices per event.}%
{yield_vs_mu}{Reconstructed vertices per event, with truth (draft Fig.~6).}
\end{frame}

\begin{frame}{HL-LHC before the GNN: Vertex Separation and Position at PU200}
\mainbadge
\twofig{deltaz_pu200}{$\Delta z$ between reconstructed vertices with the $\svtx$ fit (draft Fig.~7a).}%
{residuals_pu200}{Single-vertex residual $z_\mathrm{reco}-z_\mathrm{true}$, one-to-one within 1 mm (draft Fig.~7b).}
\end{frame}

\begin{frame}{HL-LHC before the GNN: Resolution versus $\mu$}
\mainbadge
\begin{columns}[T]
\begin{column}{.5\linewidth}\centering
\plot{p68_vs_mu}{.6}
\fcap{68\% residual versus $\mu$: the paper's resolution versus pile-up. \nEvtResMu{} events, r16443 and r16638.}
\end{column}
\begin{column}{.47\linewidth}
\footnotesize
\renewcommand{\arraystretch}{1.2}
\begin{tabular}{@{}lcc@{}}
\toprule
 & PV-Finder & AMVF \\
\midrule
68\% residual, $\mu$ 0--10 [$\mu$m] & \pseMuLoPVF & \pseMuLoAMVF \\
68\% residual, $\mu$ 180--200 [$\mu$m] & \pseMuHiPVF & \pseMuHiAMVF \\
$\svtx$, $\mu$ 100--120 [$\mu$m] & \sigMuHundredPVF & \sigMuHundredAMVF \\
$\svtx$, $\mu$ 180--200 [$\mu$m] & \sigMuTopPVF & \sigMuTopAMVF \\
AMVF / PVF $\svtx$, $\mu\ge100$ & \multicolumn{2}{c}{\sigRatioMu} \\
\bottomrule
\end{tabular}
\par\vspace{.5em}
$\svtx$ per bin only for $\mu\ge100$ (M10). The two-panel version is on the M10 slide.
\end{column}
\end{columns}
\end{frame}

\begin{frame}{HL-LHC before the GNN: Window Justification and Scan}
\twofig{window_justification}{\mainbadge{} $|z_\mathrm{reco}-z_\mathrm{true}|$ with the 68\% points and $3\times$ the larger one: the \winMM{} mm window.}%
{window_scan}{\backupbadge{} PV-Finder minus AMVF efficiency versus window.}
\end{frame}

\begin{frame}{HL-LHC before the GNN: Hard-Scatter Reconstruction}
\mainbadge
\begin{columns}[T]
\begin{column}{.58\linewidth}\centering
\plot{hs_reco_vs_mu}{.62}
\end{column}
\begin{column}{.38\linewidth}
\small
\fcap{Fraction of events with a reconstructed vertex within \hsWinMM{} mm of the true hard scatter, versus $\mu$.}
\par\vspace{.6em}
PU200: PV-Finder \hsRecoPVF\%, AMVF \hsRecoAMVF\%.
\end{column}
\end{columns}
\end{frame}

\begin{frame}{HL-LHC before the GNN: Backup}
\backupbadge
\twofig{eff_one_to_one_vs_mu}{One-to-one efficiency versus $\mu$.}%
{splits_per_event_vs_mu}{Split vertices per event versus $\mu$.}
\end{frame}

\begin{frame}{HL-LHC before the GNN: Positional Categories (Backup)}
\backupbadge
\twofig{pos_categories_pu200}{Positional categories at PU200, both algorithms.}%
{amvf_pos_categories_vs_mu}{AMVF positional categories versus $\mu$.}
\end{frame}

\begin{frame}{HL-LHC before the GNN: Efficiency vs $n_\mathrm{trk}$, Density, Separation (Backup)}
\backupbadge
\threefig{eff_ntrk_pre_pu200}{versus truth $n_\mathrm{trk}$, PU200.}%
{eff_density_pre_pu200}{versus local density, PU200.}%
{eff_dz_pre_pu200}{versus $\Delta z$ to the nearest truth vertex, PU200.}
\end{frame}

%---------------------------------------------------------------------- HL-LHC after the GNN
\begin{frame}{HL-LHC after the GNN: Table 4 at PU200}
\mainbadge\hfill{\footnotesize \hlModel{} + GNN v2 ($t=\hlT$) vs AMVF TTVA links, \nEvtPU{} events, vertices with $\ge1$ track}
\par\vspace{.2em}\scriptsize
\begin{center}
\renewcommand{\arraystretch}{1.1}
\begin{tabular}{@{}lcc@{}}
\toprule
 & PV-Finder + GNN & AMVF \\
\midrule
Clean / Merged / Split / Fake per event & \clnEvtGNN{} / \mrgEvtGNN{} / \splEvtGNN{} / \fakEvtGNN & \clnEvtAMVF{} / \mrgEvtAMVF{} / \splEvtAMVF{} / \fakEvtAMVF \\
Clean / Merged / Split / Fake per reco [\%] & \cleanRecoGNN{} / \mergedRecoGNN{} / \splitRecoGNN{} / \fakeRecoGNN & \cleanRecoAMVF{} / \mergedRecoAMVF{} / \splitRecoAMVF{} / \fakeRecoAMVF \\
Clean / truth [\%] & $\cleanTruthGNN\pm\cleanTruthGNNErr$ & \cleanTruthAMVF \\
\midrule
HS reconstructed and selected [\%] & $\hsSelGNN\pm\hsSelGNNErr$ & \hsSelAMVF{} (stored primary) \\
HS longitudinal resolution [$\mu$m] & $\hsResGNN\pm\hsResGNNErr$ & $\hsResAMVF\pm\hsResAMVFErr$ \\
Clean-HS fraction [\%] & \hsCleanGNN & \hsCleanAMVF \\
HS-track fraction [\%] & \hsTrkFracGNN & \hsTrkFracAMVF \\
\midrule
Track efficiency / purity / $F_1$ [\%] & \trkEffGNN{} / \trkPurGNN{} / \trkFOneGNN & \trkEffAMVF{} / \trkPurAMVF{} / \trkFOneAMVF \\
Wrong-assignment rate [\%] & \wrongGNN & \wrongAMVF \\
\midrule
All $\mu$ (25,000 events): Clean/truth [\%] & \cleanTruthGNNincl & \cleanTruthAMVFincl \\
All $\mu$: HS reconstructed and selected [\%] & \hsSelGNNincl & \hsSelAMVFincl \\
\bottomrule
\end{tabular}
\end{center}
\src{HLDOC; \allowbreak{}CH, CHI, SUP}
\end{frame}

\begin{frame}{HL-LHC after the GNN: Pile-up Vertex Quality at PU200}
\mainbadge
\begin{columns}[T]
\begin{column}{.6\linewidth}\centering
\plot{classes_pu200}{.62}
\end{column}
\begin{column}{.37\linewidth}
\small
\fcap{Clean, Merged, Split and Fake vertices per event at PU200, track-content classification (70\% rule), PV-Finder + GNN tracks vs AMVF TTVA links.}
\par\vspace{.5em}
Clean per event: \clnEvtGNN{} vs \clnEvtAMVF.
\end{column}
\end{columns}
\end{frame}

\begin{frame}{HL-LHC after the GNN: Vertex Categories versus $\mu$}
\mainbadge
\fourfig{clean_evt_vs_mu}{Clean / event}{merged_evt_vs_mu}{Merged / event}%
{split_evt_vs_mu}{Split / event}{fake_evt_vs_mu}{Fake / event}
\end{frame}

\begin{frame}{HL-LHC after the GNN: Clean/truth and Efficiency versus $n_\mathrm{trk}$}
\mainbadge
\twofig{clean_truth_vs_mu}{Clean vertices over truth vertices versus $\mu$.}%
{eff_ntrk_post_pu200}{Efficiency after the GNN versus truth $n_\mathrm{trk}$, PU200 (PUB-note style).}
\end{frame}

\begin{frame}{HL-LHC after the GNN: Hard-Scatter Selection and Resolution}
\mainbadge
\twofig{hs_sel_vs_mu}{HS reconstructed and selected (\hsWinMM{} mm), Athena rule for PV-Finder + GNN; AMVF stored primary and same rule.}%
{hs_res_vs_mu}{HS longitudinal resolution (iterative 3-RMS), 3D adaptive fit on GNN tracks vs AMVF.}
\end{frame}

\begin{frame}{HL-LHC after the GNN: Hard Scatter versus Local Pile-up Density}
\mainbadge
\threefig{hs_reco_vs_density}{HS reconstruction}%
{hs_sel_vs_density}{HS selection}%
{hs_res_vs_density}{HS resolution}
\par\vspace{.3em}{\footnotesize Density: truth vertices within $\pm2$ mm of the true HS, HS excluded, per mm (ITk). All $\mu$, r16443 entries 0--24,999.}
\end{frame}

\begin{frame}{HL-LHC after the GNN: Track Association}
\threefig{trk_f1_vs_mu}{\mainbadge{} Track $F_1$}%
{trk_eff_vs_mu}{\backupbadge{} Track efficiency}%
{trk_pur_vs_mu}{\backupbadge{} Track purity}
\par\vspace{.3em}{\footnotesize At PU200 the GNN trades efficiency (\trkEffGNN{} vs \trkEffAMVF\%) for purity (\trkPurGNN{} vs \trkPurAMVF\%); $F_1$ \trkFOneGNN{} vs \trkFOneAMVF\%. At low $\mu$ the GNN $F_1$ is below AMVF (M5).}
\end{frame}

\begin{frame}{HL-LHC after the GNN: Hard-Scatter Tracks (Backup)}
\backupbadge
\twofig{hs_clean_vs_mu}{Clean-HS fraction versus $\mu$.}%
{hs_trkfrac_vs_mu}{HS tracks on the HS-matched vertex over all HS tracks.}
\end{frame}

\begin{frame}{HL-LHC after the GNN: HS Resolution vs $n_\mathrm{trk}$, Density Definition (Backup)}
\backupbadge
\twofig{hs_res_vs_ntrk}{HS resolution versus HS truth $n_\mathrm{trk}$, PU200.}%
{density_definitions}{Local density: ITk definition (all truth vertices) and the stored-truth one.}
\end{frame}

\begin{frame}{HL-LHC after the GNN: Classification Cross-checks (Backup)}
\backupbadge
\twofig{classes_variants}{Categories per event at PU200: unweighted (headline), GNN-score weights, AMVF fit tracks with fit weights.}%
{amvf_fit_weights}{AMVF fit-weight distribution.}
\par\vspace{.2em}{\footnotesize With its own fit weights AMVF has \clnEvtAMVFfit{} Clean per event; PV-Finder + GNN leads by \cleanLeadFit$\times$ instead of \cleanLeadTTVA$\times$.}
\end{frame}

\begin{frame}{HL-LHC after the GNN: Association Baselines (Backup)}
\backupbadge
\twofig{eq4_vs_mu}{Eq.~(4) significance association, Clean/truth versus $\mu$.}%
{z0_baselines_vs_mu}{Nearest $z_0$ and $3\sigma(z_0)$ baselines.}
\end{frame}

\begin{frame}{HL-LHC after the GNN: Bounds and Stages (Backup)}
\backupbadge
\threefig{truth_guided_vs_mu}{Truth-guided association on the same peaks.}%
{finder_cap_vs_mu}{Finder cap (one to one, \winMM{} mm).}%
{stages_pu200}{Clean yield by stage, PU200.}
\end{frame}

\begin{frame}{HL-LHC after the GNN: Unassigned and Wrong Tracks (Backup)}
\backupbadge
\threefig{wrong_rate_vs_mu}{Wrong-assignment rate ($1-$ purity).}%
{lost_unassigned_vs_mu}{Labelled tracks left unassigned with no candidate within \winMM{} mm.}%
{unlabelled_unassigned_vs_mu}{Unlabelled tracks correctly left unassigned.}
\end{frame}

\begin{frame}{HL-LHC after the GNN: Density, Separation, Peak Width (Backup)}
\backupbadge
\threefig{eff_density_post_pu200}{Efficiency after the GNN versus local density.}%
{eff_dz_post_pu200}{versus $\Delta z$ to the nearest truth vertex.}%
{sigma_pv_pu200}{Peak width ($\sigma_\mathrm{PV}$) vs the $n_\mathrm{trk}$ parameterisation.}
\end{frame}

\begin{frame}{GNN v2: Learning Curve and Metrics versus $\mu$ (Backup)}
\backupbadge
\twofig{gnn_learning_curve}{Validation-monitor loss, v1 and v2.}%
{vs_mu_test_s8}{All chain metrics versus $\mu$ for several $t$, \hlModel{} peaks, test events.}
\end{frame}

%---------------------------------------------------------------------- Low pile-up
\begin{frame}{Low Pile-up: Vertex Finding (Table 6)}
\mainbadge\hfill{\footnotesize \lpModel, integral \lpIntegral, window \winMM{} mm}
\par\vspace{.4em}\footnotesize
\begin{center}
\renewcommand{\arraystretch}{1.2}
\begin{tabular}{@{}lccc@{}}
\toprule
Run 2 MC, \lpMCTest{} events, \lpTruthMC{} truth vertices & PV-Finder & AMVF & PVF $-$ AMVF \\
\midrule
Efficiency [\%] & $\lpEffPVF\pm\lpEffPVFErr$ & \lpEffAMVF & \lpEffDiff{} pt \\
Fake / reco [\%] & \lpFakeFracPVF & \lpFakeFracAMVF & \\
Fakes / event & \lpFakeEvtPVF & \lpFakeEvtAMVF & \\
$\svtx$ [mm] & $\lpSigPVF\pm\lpSigPVFErr$ & $\lpSigAMVF\pm\lpSigAMVFErr$ & \\
68\% residual [$\mu$m] & \lpPsePVF & \lpPseAMVF & \\
HS reconstruction, \hsWinMM{} mm [\%] & \lpHSPVF & \lpHSAMVF & \\
\bottomrule
\end{tabular}
\par\vspace{.6em}
\begin{tabular}{@{}lccc@{}}
\toprule
Collision data, AMVF as reference & Run 2 data & Run 3 data & Run 2 MC, same convention \\
\midrule
AMVF-vertex recovery [\%] & \lpRecRunTwo & \lpRecRunThree & \lpRecMCref \\
PVF vertices without AMVF vertex / event & \lpNoRefRunTwo & \lpNoRefRunThree & \\
$\svtx$ PV-Finder / AMVF [mm] & \lpSigRunTwo{} / \lpSigAMVFRunTwo & \lpSigRunThree{} / \lpSigAMVFRunThree & \\
\bottomrule
\end{tabular}
\end{center}
\src{LPDOC R1, R3; \allowbreak{}VF}
\end{frame}

\begin{frame}{Low Pile-up: Run 2 MC Vertex Finding (Fig.~8)}
\mainbadge
\twofig{lp_fig8a}{(a) $\Delta z$ between vertices, PV-Finder and AMVF.}%
{lp_fig8b}{(b) Vertices per event versus $\mu$.}
\end{frame}

\begin{frame}{Low Pile-up: Run 2 MC after the GNN (Table 5, Fig.~9)}
\mainbadge
\begin{columns}[T]
\begin{column}{.36\linewidth}
\footnotesize
\renewcommand{\arraystretch}{1.18}
\begin{tabular}{@{}lcc@{}}
\toprule
{[\%]} & PVF + GNN & AMVF \\
\midrule
Clean & \lpClean & \lpCleanAMVF \\
Merged & \lpMerged & \lpMergedAMVF \\
Split & \lpSplit & \lpSplitAMVF \\
Fake & \lpFake & \lpFakeAMVF \\
Clean/truth & \lpCleanTruth & \lpCleanTruthAMVF \\
Track $F_1$ & \lpFOne & \lpFOneAMVF \\
HS ident. & \lpHSID & \lpHSIDAMVF \\
\bottomrule
\end{tabular}
\par\vspace{.3em}{\scriptsize $t=\lpT$; per reco vertex with $\ge1$ track; \lpEmpty{} empty candidates dropped.}
\end{column}
\begin{column}{.31\linewidth}\centering
\plot{lp_fig9a}{.5}\fcap{(a) Matched-vertex efficiency vs truth $n_\mathrm{trk}$.}
\end{column}
\begin{column}{.31\linewidth}\centering
\plot{lp_fig9b}{.5}\fcap{(b) Categories.}
\end{column}
\end{columns}
\end{frame}

\begin{frame}{Low Pile-up: Collision Data (Fig.~10)}
\mainbadge
\fourfig{lp_fig10a}{(a) Run 2 data $\Delta z$, with Run 2 MC}{lp_fig10b}{(b) Run 2 data vertices vs $\mu$}%
{lp_fig10c}{(c) Run 3 data $\Delta z$, with Run 2 MC}{lp_fig10d}{(d) Run 3 data vertices vs $\mu$}
\end{frame}

\begin{frame}{Low Pile-up: Run 3 Data Association and the $t$ Scan}
\begin{columns}[T]
\begin{column}{.45\linewidth}
\mainbadge\par\vspace{.4em}
\footnotesize
\renewcommand{\arraystretch}{1.2}
\begin{tabular}{@{}lc@{}}
\toprule
Run 3 data, 2,500 events, $t=\lpT$ & [\%] \\
\midrule
Tracks assigned by the GNN & \lpRthreeAssigned \\
Agreement with AMVF (1 mm) & \lpRthreeAgree \\
Agreement with AMVF (\winMM{} mm) & \lpRthreeAgreeW \\
Leading-vertex agreement & \lpRthreeLead \\
\bottomrule
\end{tabular}
\par\vspace{.3em}{\scriptsize Same events and peaks as the vertex finding. Agreement, not efficiency.}
\end{column}
\begin{column}{.52\linewidth}\centering
\backupbadge\par
\plot{lp_t_scan}{.55}\fcap{Run 2 MC metrics versus $t$.}
\end{column}
\end{columns}
\end{frame}

%---------------------------------------------------------------------- Timing
\begin{frame}{Timing: CPU and GPU}
\begin{columns}[T]
\begin{column}{.55\linewidth}\centering
\mainbadge\par
\plot{athena_timing}{.6}
\fcap{Athena timing versus pile-up: AMVF on CPU, PV-Finder chain on GPU (note figure).}
\end{column}
\begin{column}{.42\linewidth}
\footnotesize
\renewcommand{\arraystretch}{1.2}
\begin{tabular}{@{}lc@{}}
\toprule
$\mu=200$, per event & [ms] \\
\midrule
PV-Finder chain, GPU (Athena) & \athGPUms \\
PV-Finder chain, CPU (Athena, v6) & \athCPUms \\
\hlModel{} chain, CPU, projected & \projChain \\
AMVF, CPU (note figure) & \athAMVF \\
\bottomrule
\end{tabular}
\par\vspace{.5em}
The CPU claim in the paper needs packed inference in Athena (M6). The GPU claim stands as measured.
\end{column}
\end{columns}
\end{frame}


%======================================================================
\partdivider{7}{Proposed Internal-Note Edits}{Old text and new [MM: ...] text, with page and line}

% Proposed internal-note edits. To be filled by the note-edit agent:
% replace each \noteedit placeholder by the page, line, old text and the new
% red [MM: ...] text. The topics are those of
% docs/additional_studies/meeting_2026-11-11_plan.md ("Note edits").

\newcommand{\editbox}[2]{% label, colour
  {\setlength{\fboxsep}{5pt}\fcolorbox{#2!60!black}{#2!4}{\parbox[c][.22\textheight][t]{.93\linewidth}{%
  {\scriptsize\bfseries\textcolor{#2!60!black}{#1}}\par\vspace{.3em}{\footnotesize\textcolor{red}{TBD}}}}}}
\newcommand{\noteedit}[2]{% number, topic
\begin{frame}{Note Edit #1: #2}
\vspace*{.8em}{\small\textbf{Note page, line}: \textcolor{red}{TBD} \hfill \textbf{Section}: \textcolor{red}{TBD}}
\par\vspace{.5em}\centering
\editbox{Old text}{black}\par\vspace{.6em}
\editbox{New text, in red [MM: ...]}{red}
\end{frame}}

\begin{frame}{Note Edits: List}
\small
Internal note, GitLab main. Changes in red as \textcolor{red}{[MM: ...]} on a branch with a merge request.
\begin{enumerate}\setlength\itemsep{.3em}
\item The \winMM{} mm window with its justification plot; the hard scatter also at \hsWinMM{} mm.
\item Split-inclusive efficiency: one sentence.
\item HS longitudinal resolution from the 3D adaptive fit: one paragraph and one figure.
\item Timing claim restricted to GPU.
\item Simulation on the Fig.~6.2 panels, and the AMVF-reference bias.
\item $z_0$ baselines in Fig.~7.2.
\item AMVF HS selection \hsSelAMVF\% and the HS numbers.
\item The AMVF-weights paragraph.
\item The frozen configuration: \hlModel, GNN v2 epoch \hlGnnEpoch, $t=\hlT$; \lpModel{} at low pile-up.
\end{enumerate}
\end{frame}

\noteedit{1}{Matching Window}
\noteedit{2}{Split-Inclusive Efficiency}
\noteedit{3}{Hard-Scatter Longitudinal Resolution}
\noteedit{4}{Timing Claim}
\noteedit{5}{Fig.~6.2 and the AMVF Reference}
\noteedit{6}{$z_0$ Baselines in Fig.~7.2}
\noteedit{7}{AMVF Hard-Scatter Selection}
\noteedit{8}{AMVF Weights}
\noteedit{9}{Frozen Configuration}


\end{document}
